% Les effets des mouvements migratoires — Géographie 1re Tech — Cours signé OnBuch+
% Programme officiel MINESEC. Compiler avec : tectonic N02-effets-mouvements-migratoires.tex
\newcommand{\DOCMATIERE}{Géographie}
\newcommand{\DOCNIVEAU}{1re Tech}
\newcommand{\DOCMODULE}{Géographie – classe de Première technique}
\newcommand{\DOCLECON}{N02}
\newcommand{\DOCTITRE}{Les effets des mouvements migratoires}
% OnBuch+ — préambule « cours signés » ENRICHI (compilé avec tectonic / XeLaTeX)
% Système visuel « Pop » d'OnBuch+ : fond crème (jamais blanc), bordures encre
% épaisses, ombres « dures » décalées, titres Archivo Black en capitales.
% Version MATHÉMATIQUES : graphes (pgfplots/tikz), boîtes pédagogiques
% (définition, propriété, méthode, attention, le savais-tu, exercice résolu,
% exercice, TP, bilan), page de garde et signature OnBuch+ ancrée en bas.
%
% Macros attendues dans main.tex AVANT \input{preamble} :
%   \DOCMATIERE  \DOCNIVEAU  \DOCMODULE  \DOCLECON (numéro)  \DOCTITRE
\documentclass[11pt]{article}

\usepackage{fontspec}
\usepackage[french]{babel}
\usepackage[a4paper,margin=2cm,top=3.4cm,bottom=2.8cm,headheight=1.4cm,footskip=1.3cm]{geometry}
\usepackage{amsmath,amssymb,amsfonts}
\usepackage{mathtools} % \xmapsto, \coloneqq... souvent utilisés par les modèles
\usepackage{cancel} % simplifications barrées (bilans de forces)
% \abs{z} (module, valeur absolue) et \norm{u} (norme), utilisés par les modèles
\providecommand{\abs}{}\let\abs\relax\DeclarePairedDelimiter{\abs}{\lvert}{\rvert}
\providecommand{\norm}{}\let\norm\relax\DeclarePairedDelimiter{\norm}{\lVert}{\rVert}
\usepackage[table]{xcolor} % [table] : fournit aussi \rowcolors (alternance de lignes)
\usepackage{pagecolor}
\usepackage{tikz}
\usetikzlibrary{arrows.meta,positioning,calc,shapes.geometric,shapes.misc,shapes.arrows,decorations.pathmorphing,decorations.pathreplacing,decorations.markings,patterns,shadows,mindmap,trees,fit,backgrounds,matrix,intersections,angles,quotes}
\usepackage{pgfplots}
\usepackage[version=4]{mhchem} % équations nucléaires \ce{^{235}_{92}U}
\usepackage[european,siunitx]{circuitikz} % schémas électriques
\pgfplotsset{compat=1.18}
\usepgfplotslibrary{fillbetween,groupplots} % aires entre courbes (calcul intégral)
\usepackage{enumitem}
\usepackage{fancyhdr}
% [most] charge toutes les bibliothèques tcolorbox (listings, minted, raster,
% documentation...) et gonfle inutilement la mémoire TeX sur un document long
% (~300 boîtes) : on ne charge que ce qui est réellement utilisé.
\usepackage[skins,breakable]{tcolorbox}
\usepackage{graphicx}
\usepackage{colortbl}
\usepackage{booktabs}
\usepackage{tabularx}
\usepackage{multirow}
\usepackage{diagbox} % case d'en-tête diagonale des tableaux à double entrée
% Colonnes extensibles centrée / gauche / droite (C, L, R), souvent utilisées par les modèles
\newcolumntype{C}{>{\centering\arraybackslash}X}
\newcolumntype{L}{>{\raggedright\arraybackslash}X}
\newcolumntype{R}{>{\raggedleft\arraybackslash}X}
\usepackage{array}
\usepackage{multicol}
\usepackage{siunitx}
% parse-numbers=false : accepte aussi \SI{2,5 \times 10^{-3}}{...}
\sisetup{locale=FR,output-decimal-marker={,},group-minimum-digits=4,parse-numbers=false}
\DeclareSIUnit\ohm{\ensuremath{\symup{\Omega}}} % U+2126 absent de la police
\DeclareSIUnit\division{div} % réglages d'oscilloscope (ms/div, V/div)
\DeclareSIUnit\tr{tr} % tours (vitesse de rotation, ex. \SI{1500}{\tr\per\minute})
\DeclareSIUnit\molar{M} % concentration molaire (ex. \SI{3}{\molar}, \SI{1}{\micro\molar})
\DeclareSIUnit\an{an} % année (le modèle confond parfois avec \year, un compteur TeX) — ex. \SI{5}{\kilo\meter\per\an}
\DeclareSIUnit\FCFA{FCFA} % franc CFA (ex. \SI{500}{\FCFA\per\kilo\gram})
\DeclareSIUnit\degreeBrix{\textdegree Brix} % teneur en sucre d'un sirop/jus (réfractométrie)
\DeclareSIUnit\month{mois} % le modèle utilise parfois \month, un compteur TeX, comme unité de durée
\DeclareSIUnit\microliter{\micro\liter} % le modèle écrit parfois \microliter en un seul token
\DeclareSIUnit\million{millions} % dénombrement (ex. \SI{15}{\million\per\mL} spermatozoïdes)
\usepackage{qrcode}
% \degree : commande fréquemment utilisée par les modèles pour un angle ou une
% température en dehors de \SI{}{} (non fournie par les paquets chargés).
\providecommand{\degree}{^{\circ}}
% \qed (fin de démonstration), utilisé par les modèles sans amsthm
\providecommand{\qed}{\hfill$\square$}
% \barre : double barre des valeurs interdites dans un tableau de variations (tkz-tab)
\providecommand{\barre}{\ensuremath{\|}}
% environnement proof (amsthm non chargé) utilisé par les modèles
\ifdefined\proof\else\newenvironment{proof}[1][Démonstration]{\par\noindent\textit{#1.}\ }{\qed\par}\fi
\usepackage{lastpage}
\usepackage{hyperref}
\hypersetup{hidelinks}

% ============================================================
% Palette « Pop » OnBuch+ — mêmes hex que la classe P (plus_ui.dart)
% ============================================================
\definecolor{poporange}{HTML}{F59321}
\definecolor{popdark}{HTML}{E07A0C}
\definecolor{popcream}{HTML}{FFF6E8}
\definecolor{popink}{HTML}{1C1714}
\definecolor{poppurple}{HTML}{7A5AE0}
\definecolor{popgreen}{HTML}{2E9E6A}
\definecolor{poppink}{HTML}{E0457B}
\definecolor{popblue}{HTML}{2F7DE1}
\definecolor{popgold}{HTML}{C98A12}
\definecolor{popmuted}{HTML}{8A7F76}
\definecolor{popline}{HTML}{EADFD2}
\definecolor{popgray}{HTML}{8A7F76}
\colorlet{popgrey}{popgray}
% Alias tolérants (le modèle nomme parfois les couleurs différemment)
\colorlet{popwhite}{white}
\colorlet{popblack}{popink}
\colorlet{popred}{poppink}
\colorlet{popyellow}{popgold}
% Teintes claires (fonds de tableaux/encadrés) — le modèle nomme parfois
% les couleurs avec un suffixe L (ex. popblueL) sans les avoir définies.
\colorlet{poporangeL}{poporange!20}
\colorlet{popdarkL}{popdark!20}
\colorlet{poppurpleL}{poppurple!20}
\colorlet{popgreenL}{popgreen!20}
\colorlet{poppinkL}{poppink!20}
\colorlet{popblueL}{popblue!20}
\colorlet{popgoldL}{popgold!20}
\colorlet{popredL}{popred!20}
\colorlet{popgrayL}{popgray!20}
% Teintes claires pour fonds de tableaux / zones
\colorlet{popblueL}{popblue!10!white}
\colorlet{poporangeL}{poporange!14!white}
\colorlet{popgreenL}{popgreen!12!white}
\colorlet{poppurpleL}{poppurple!12!white}
\colorlet{poppinkL}{poppink!12!white}
\colorlet{popgoldL}{popgold!14!white}

\pagecolor{popcream}

% ============================================================
% Typographie
% ============================================================
\defaultfontfeatures{Ligatures=TeX}
\setmainfont{PlusJakartaSans-Medium.ttf}[Path=../../../fonts/,BoldFont=PlusJakartaSans-Bold.ttf]
% Maths en sans-serif (Fira Math) avec chiffres et lettres droites de Plus
% Jakarta, pour que les formules se fondent dans le texte.
\usepackage[math-style=ISO,bold-style=ISO]{unicode-math}
\setmathfont{FiraMath-Regular.otf}[Path=../../../fonts/]
\setmathfont{PlusJakartaSans-Medium.ttf}[Path=../../../fonts/,range={up/num,up/{Latin,latin},\mathup}]
\usepackage{icomma} % virgule décimale sans espace en mode math
% Caractères mathématiques Unicode absents des polices de texte (Plus Jakarta,
% Archivo Black) : rendus par leur équivalent mathématique, en texte comme en
% formule — sinon ils s'affichent en carré vide (ex. « Arithmétique dans ℤ »).
\usepackage{newunicodechar}
\newunicodechar{ℝ}{\ensuremath{\mathbb{R}}}
\newunicodechar{ℤ}{\ensuremath{\mathbb{Z}}}
\newunicodechar{ℂ}{\ensuremath{\mathbb{C}}}
\newunicodechar{ℕ}{\ensuremath{\mathbb{N}}}
\newunicodechar{ℚ}{\ensuremath{\mathbb{Q}}}
\newunicodechar{𝔻}{\ensuremath{\mathbb{D}}}
\newunicodechar{α}{\ensuremath{\alpha}}
\newunicodechar{β}{\ensuremath{\beta}}
\newunicodechar{γ}{\ensuremath{\gamma}}
\newunicodechar{δ}{\ensuremath{\delta}}
\newunicodechar{ε}{\ensuremath{\varepsilon}}
\newunicodechar{θ}{\ensuremath{\theta}}
\newunicodechar{λ}{\ensuremath{\lambda}}
\newunicodechar{μ}{\ensuremath{\mu}}
\newunicodechar{σ}{\ensuremath{\sigma}}
\newunicodechar{τ}{\ensuremath{\tau}}
\newunicodechar{φ}{\ensuremath{\varphi}}
\newunicodechar{ψ}{\ensuremath{\psi}}
\newunicodechar{ω}{\ensuremath{\omega}}
\newunicodechar{Σ}{\ensuremath{\Sigma}}
% majuscules produites par \MakeUppercase dans les titres (α-aminés -> Α)
\newunicodechar{Α}{\ensuremath{\alpha}}
\newunicodechar{Β}{\ensuremath{\beta}}
\newunicodechar{Γ}{\ensuremath{\Gamma}}
\newunicodechar{Δ}{\ensuremath{\Delta}}
\newunicodechar{Ε}{\ensuremath{\varepsilon}}
\newunicodechar{Θ}{\ensuremath{\Theta}}
\newunicodechar{Λ}{\ensuremath{\Lambda}}
\newunicodechar{Μ}{\ensuremath{\mu}}
\newunicodechar{Τ}{\ensuremath{\tau}}
\newunicodechar{Φ}{\ensuremath{\Phi}}
\newunicodechar{Ψ}{\ensuremath{\Psi}}
\newunicodechar{Ω}{\ensuremath{\Omega}}
\newunicodechar{↦}{\ensuremath{\mapsto}}
\newunicodechar{∈}{\ensuremath{\in}}
\newunicodechar{∉}{\ensuremath{\notin}}
\newunicodechar{∧}{\ensuremath{\wedge}}
\newunicodechar{⇔}{\ensuremath{\Leftrightarrow}}
\newunicodechar{⟺}{\ensuremath{\Longleftrightarrow}}
\newunicodechar{⇒}{\ensuremath{\Rightarrow}}
\newunicodechar{⟹}{\ensuremath{\Longrightarrow}}
\newunicodechar{∘}{\ensuremath{\circ}}
\newunicodechar{∪}{\ensuremath{\cup}}
\newunicodechar{∩}{\ensuremath{\cap}}
\newunicodechar{≡}{\ensuremath{\equiv}}
\newunicodechar{⊂}{\ensuremath{\subset}}
\newunicodechar{⊥}{\ensuremath{\perp}}
\newunicodechar{∃}{\ensuremath{\exists}}
\newunicodechar{∀}{\ensuremath{\forall}}
\newunicodechar{∛}{\ensuremath{\sqrt[3]{\,}}}
\newunicodechar{∜}{\ensuremath{\sqrt[4]{\,}}}
\newunicodechar{⃗}{\ensuremath{{}^{\to}}}
\newunicodechar{ⁿ}{\ifmmode^{n}\else\textsuperscript{\ensuremath{n}}\fi}
\newunicodechar{ˣ}{\ifmmode^{x}\else\textsuperscript{\ensuremath{x}}\fi}
\newunicodechar{ᵇ}{\ifmmode^{b}\else\textsuperscript{\ensuremath{b}}\fi}
\newunicodechar{ʳ}{\ifmmode^{r}\else\textsuperscript{\ensuremath{r}}\fi}
\newunicodechar{ᵘ}{\ifmmode^{u}\else\textsuperscript{\ensuremath{u}}\fi}
\newunicodechar{ʸ}{\ifmmode^{y}\else\textsuperscript{\ensuremath{y}}\fi}
\newunicodechar{ᵃ}{\ifmmode^{a}\else\textsuperscript{\ensuremath{a}}\fi}
\newunicodechar{ᵉ}{\ifmmode^{e}\else\textsuperscript{\ensuremath{e}}\fi}
\newunicodechar{ᵏ}{\ifmmode^{k}\else\textsuperscript{\ensuremath{k}}\fi}
\newunicodechar{ᵖ}{\ifmmode^{p}\else\textsuperscript{\ensuremath{p}}\fi}
\newunicodechar{ᴺ}{\ifmmode^{N}\else\textsuperscript{\ensuremath{N}}\fi}
\newunicodechar{ᶜ}{\ifmmode^{c}\else\textsuperscript{\ensuremath{c}}\fi}
\newunicodechar{ʰ}{\ifmmode^{h}\else\textsuperscript{\ensuremath{h}}\fi}
\newunicodechar{ᵠ}{\ifmmode^{q}\else\textsuperscript{\ensuremath{q}}\fi}
\newunicodechar{ₙ}{\ifmmode_{n}\else\textsubscript{\ensuremath{n}}\fi}
\newunicodechar{ₐ}{\ifmmode_{a}\else\textsubscript{\ensuremath{a}}\fi}
\newunicodechar{ᵢ}{\ifmmode_{i}\else\textsubscript{\ensuremath{i}}\fi}
\newunicodechar{ᵣ}{\ifmmode_{r}\else\textsubscript{\ensuremath{r}}\fi}
\newunicodechar{ₖ}{\ifmmode_{k}\else\textsubscript{\ensuremath{k}}\fi}
\newunicodechar{ᵦ}{\ifmmode_{\beta}\else\textsubscript{\ensuremath{\beta}}\fi}
\newunicodechar{ₓ}{\ifmmode_{x}\else\textsubscript{\ensuremath{x}}\fi}
\newfontfamily\popdisplay{ArchivoBlack-Regular.ttf}[Path=../../../fonts/]
\newfontfamily\poplogo{SpaceGrotesk-Bold.ttf}[Path=../../../fonts/]
\newfontfamily\popsemi{PlusJakartaSans-SemiBold.ttf}[Path=../../../fonts/]
\newfontfamily\popxbold{PlusJakartaSans-ExtraBold.ttf}[Path=../../../fonts/]
\newfontfamily\popxboldls{PlusJakartaSans-ExtraBold.ttf}[Path=../../../fonts/,LetterSpace=3.0]

\setlist[enumerate]{leftmargin=1.7em,itemsep=5pt,topsep=4pt}
\setlist[itemize]{leftmargin=1.7em,itemsep=4pt,topsep=4pt,label={\color{poporange}\textbullet}}
\setlength{\parindent}{0pt}
\setlength{\parskip}{5pt}
\renewcommand{\arraystretch}{1.25}

% pgfplots : style Pop par défaut
\pgfplotsset{
  every axis/.append style={axis line style={popink,line width=1pt},tick style={popink},
    grid style={popline,line width=0.5pt},label style={font=\small},tick label style={font=\footnotesize}},
  popcurve/.style={poporange,line width=1.8pt,no markers,smooth},
}
% Même style disponible dans un \draw TikZ simple (hors pgfplots)
\tikzset{popcurve/.style={poporange,line width=1.8pt}}
\tikzset{popgrid/.style={popline,very thin}}
\colorlet{popcurve}{poporange} % le nom du style est parfois utilisé comme couleur

% ============================================================
% Pill / squiggle
% ============================================================
\newcommand{\poppill}[3]{%
  \tikz[baseline=-0.55ex]\node[draw=popink,line width=1.2pt,rounded corners=2.6pt,%
    fill=#2,inner xsep=6.5pt,inner ysep=3pt]{%
    \popxboldls\fontsize{7}{7}\selectfont\color{#3}\MakeUppercase{#1}};%
}
\newcommand{\popsquiggle}[1]{%
  \tikz[baseline=-0.4ex]{
    \draw[#1,line width=2.6pt,line cap=round]
      plot [smooth] coordinates {(0,0.09) (0.34,0.01) (0.68,0.21) (1.02,0.01) (1.36,0.10)};
  }%
}
% Mot-clé surligné (marqueur orange sous le texte)
\newcommand{\cle}[1]{{\popxbold\color{popdark}#1}}

% ============================================================
% En-tête / pied
% ============================================================
\pagestyle{fancy}
\fancyhf{}
\renewcommand{\headrulewidth}{0pt}
\renewcommand{\footrulewidth}{0pt}
\lhead{\raisebox{-0.15ex}{\poplogo\fontsize{13}{13}\selectfont\color{popink}OnBuch\color{poporange}+}}
\rhead{\poppill{\DOCMATIERE\ \textbullet\ \DOCNIVEAU\ \textbullet\ Leçon \DOCLECON}{white}{popink}}
\lfoot{\popxbold\fontsize{7.6}{7.6}\selectfont\color{popmuted}\MakeUppercase{Cours signé OnBuch+}\ \textbullet\ {\popsemi\DOCTITRE}}
\rfoot{\tikz[baseline=-0.6ex]\node[rounded corners=8.5pt,fill=popink,minimum height=17pt,minimum width=17pt,inner xsep=5pt,inner sep=0]{\popxbold\fontsize{8}{8}\selectfont\color{popcream}\thepage\,/\,\pageref*{LastPage}};}

% ============================================================
% Titres
% ============================================================
\newcommand{\chaptitre}[2]{%
  \begin{tcolorbox}[enhanced,colback=poporange,colframe=popink,boxrule=2.6pt,arc=11pt,
      drop shadow=popink,left=15pt,right=15pt,top=12pt,bottom=11pt,before skip=3pt,after skip=18pt]
    \poppill{#2}{popink}{popcream}\par\vspace{9pt}
    {\raggedright\hyphenpenalty=10000\exhyphenpenalty=10000\popdisplay\fontsize{22}{24}\selectfont\color{popcream}\MakeUppercase{#1}\par}
  \end{tcolorbox}
}
% Section numérotée : pastille orange avec le numéro + titre Archivo
\newcounter{popsec}
\newcommand{\coursec}[1]{%
  \stepcounter{popsec}\setcounter{popsub}{0}%
  \par\vspace{16pt}\noindent
  \begin{minipage}{\linewidth}
  \tikz[baseline=-0.7ex]\node[draw=popink,line width=1.6pt,fill=poporange,rounded corners=4pt,minimum size=22pt,inner sep=0]{\popdisplay\fontsize{12}{12}\selectfont\color{popcream}\thepopsec};%
  \hspace{8pt}{\popdisplay\fontsize{15}{17}\selectfont\color{popink}\MakeUppercase{#1}}\par
  \vspace{4pt}\noindent{\color{poporange}\rule{36pt}{3.6pt}}
  \end{minipage}\par\nopagebreak\vspace{8pt}
}
\newcounter{popsub}
\newcommand{\courssub}[1]{%
  \stepcounter{popsub}%
  \par\vspace{9pt}\noindent{\popxbold\fontsize{11.5}{13}\selectfont\color{popink}{\color{poporange}\thepopsec.\thepopsub}\hspace{6pt}#1}\par\nopagebreak\vspace{4pt}
}

% ============================================================
% Boîtes pédagogiques — même recette : fond blanc, bordure épaisse colorée,
% ombre dure encre, badge plein. Titre optionnel [#1].
% ============================================================
\tcbset{popbase/.style={breakable,enhanced,colback=white,boxrule=2.3pt,arc=9pt,
  shadow={3pt}{-3pt}{0pt}{popink},left=14pt,right=14pt,top=11pt,bottom=11pt,before skip=11pt,after skip=15pt,
  fontupper=\popsemi\fontsize{10.6}{14.5}\selectfont\color{popink},
  fontlower=\popsemi\fontsize{10.6}{14.5}\selectfont\color{popink}}}
% Coche dessinée (le glyphe ✓ n'existe pas dans les polices du document)
\newcommand{\popcheck}[1]{\tikz[baseline=-0.5ex]\draw[#1,line width=1.6pt,line cap=round,line join=round](-0.13,0.0)--(-0.04,-0.09)--(0.13,0.1);}
\providecommand{\coche}{\popcheck{popgreen}}
\providecommand{\resultat}[1]{\par\smallskip\noindent\fcolorbox{popgreen}{popgreenL}{\parbox{\dimexpr\linewidth-2\fboxsep-2\fboxrule\relax}{\textbf{Résultat.} #1}}\par\smallskip}
\newcommand{\popboxhead}[3]{\poppill{#1}{#2}{white}\ifx\relax#3\relax\else\hspace{6pt}{\popxbold\fontsize{10.6}{12}\selectfont\color{popink}#3}\fi\par\vspace{7pt}}

\newtcolorbox{aretenir}[1][]{popbase,colframe=popgreen,before upper={\popboxhead{À retenir}{popgreen}{#1}}}
\newtcolorbox{exemplebox}[1][]{popbase,colframe=poporange,before upper={\popboxhead{Exemple}{poporange}{#1}}}
\newtcolorbox{definition}[1][]{popbase,colframe=popblue,before upper={\popboxhead{Définition}{popblue}{#1}}}
\newtcolorbox{propriete}[1][]{popbase,colframe=poppurple,before upper={\popboxhead{Propriété}{poppurple}{#1}}}
\newtcolorbox{methode}[1][]{popbase,colframe=popgold,before upper={\popboxhead{Méthode}{popgold}{#1}}}
\newtcolorbox{attention}[1][]{popbase,colframe=poppink,before upper={\popboxhead{Attention}{poppink}{#1}}}
\newtcolorbox{experience}[1][]{popbase,colframe=popblue,colback=popblueL,before upper={\popboxhead{Expérience}{popblue}{#1}}}
% « Le savais-tu ? » : carte violette pleine (contexte, culture, Cameroun)
\newtcolorbox{savaistu}[1][]{popbase,colback=poppurple,colframe=popink,
  fontupper=\popsemi\fontsize{10.4}{14.2}\selectfont\color{white},
  before upper={\poppill{Le savais-tu\,?}{white}{poppurple}\ifx\relax#1\relax\else\hspace{6pt}{\popxbold\color{white}#1}\fi\par\vspace{7pt}}}
% Exercice résolu : énoncé en haut, \tcblower puis la solution sur fond crème
\newcounter{popexo}\newcounter{popexr}
\newtcolorbox{exoresolu}[1][]{popbase,colframe=poporange,colbacklower=popcream,
  segmentation style={popink,line width=1.2pt,dashed},
  before upper={\stepcounter{popexr}\popboxhead{Exercice résolu \thepopexr}{poporange}{#1}},
  before lower={\poppill{Solution}{popink}{popcream}\par\vspace{6pt}}}
% Exercice (non résolu en place) — la correction est dans la partie Corrigés
\newtcolorbox{exercice}[1][]{popbase,colframe=popink,
  before upper={\stepcounter{popexo}\popboxhead{Exercice \thepopexo}{popink}{#1}}}
% Corrigé
\newtcolorbox{corrige}[1][]{popbase,colframe=popgreen,colback=popgreenL,
  before upper={\popboxhead{Corrigé}{popgreen}{#1}}}
% Objectifs / prérequis / bilan
\newtcolorbox{objectifs}{popbase,colframe=popink,colback=poporangeL,
  before upper={\popboxhead{Objectifs de la leçon}{popink}{}}}
\newtcolorbox{prerequis}{popbase,colframe=popink,colback=popblueL,
  before upper={\popboxhead{Prérequis}{popblue}{}}}
\newtcolorbox{bilan}{popbase,colframe=popink,colback=white,boxrule=2.8pt,
  before upper={\popboxhead{Fiche bilan}{poporange}{L'essentiel à connaître pour l'examen}}}
% Figure encadrée avec légende
\newtcolorbox{popfigure}[1][]{enhanced,breakable=false,colback=white,colframe=popline,boxrule=1.4pt,arc=8pt,
  left=10pt,right=10pt,top=10pt,bottom=8pt,before skip=10pt,after skip=12pt,halign=center,
  lower separated=false,halign lower=center,
  fontlower=\popsemi\fontsize{9}{11.5}\selectfont\color{popmuted},
  #1}
\newcommand{\legende}[1]{\par\vspace{6pt}{\popsemi\fontsize{9}{11.5}\selectfont\color{popmuted}#1}}

% ============================================================
% Signature OnBuch+
% ============================================================
\newcommand{\poplogoinline}[1]{{\poplogo\fontsize{#1}{#1}\selectfont\color{popink}OnBuch\color{poporange}+}}
\newcommand{\signatureonbuch}{%
  \begin{tcolorbox}[enhanced,colback=popink,colframe=popink,boxrule=2.2pt,arc=10pt,
      drop shadow=poporange,left=14pt,right=14pt,top=10pt,bottom=10pt,before skip=0pt,after skip=0pt]
    \tikz[baseline=-0.7ex]\node[circle,fill=popgreen,draw=popcream,line width=1.3pt,minimum size=22pt,inner sep=0]{\popcheck{white}};%
    \hspace{9pt}\begin{minipage}[c]{0.62\linewidth}
      {\popxbold\fontsize{12}{13}\selectfont\color{popcream}Signé\ \textbullet\ {\poplogo OnBuch\color{poporange}+}}\par\vspace{2pt}
      {\raggedright\popsemi\fontsize{8}{10.5}\selectfont\color{popline}\MakeUppercase{Équipe pédagogique OnBuch+}\\\MakeUppercase{Conforme au programme officiel MINESEC}\par}
    \end{minipage}\hfill
    {\poplogo\fontsize{22}{22}\selectfont\color{popcream}OnBuch\color{poporange}+}
  \end{tcolorbox}%
}

% ============================================================
% Page de garde
% #1 titre · #2 matière · #3 niveau · #4 module · #5 n° leçon · #6 durée ·
% #7 description / objectifs (texte ou itemize)
% ============================================================
\newcommand{\pagegarde}[7]{%
  \thispagestyle{empty}
  \newgeometry{a4paper,margin=1.7cm,top=1.6cm,bottom=1.4cm}%
  \begin{tikzpicture}[remember picture,overlay]
    \fill[poporange!18] ([xshift=-3.2cm,yshift=-3.6cm]current page.north east) circle (4.6cm);
    \fill[poppurple!14] ([xshift=2.4cm,yshift=7.6cm]current page.south west) circle (3.8cm);
  \end{tikzpicture}%
  {\poplogo\fontsize{30}{30}\selectfont\color{popink}OnBuch\color{poporange}+}\hfill
  \raisebox{0.6ex}{\poppill{Cours signé \textbullet\ Premium}{popink}{popcream}}\par
  \vspace{1pt}\popsquiggle{poppurple}\par
  \vspace{8pt}
  \begin{tcolorbox}[enhanced,colback=poporange,colframe=popink,boxrule=2.8pt,arc=13pt,
      drop shadow=popink,left=18pt,right=18pt,top=12pt,bottom=13pt,after skip=10pt]
    \poppill{#2 \textbullet\ #3}{popink}{popcream}\hspace{5pt}\poppill{#4}{white}{popink}\par\vspace{8pt}
    {\popdisplay\fontsize{14}{14}\selectfont\color{popink}LEÇON #5}\par\vspace{4pt}
    {\raggedright\hyphenpenalty=10000\exhyphenpenalty=10000\popdisplay\fontsize{25}{27}\selectfont\color{popcream}\MakeUppercase{#1}\par}
    \vspace{8pt}
    {\popxbold\fontsize{9.5}{9.5}\selectfont\color{popink}\MakeUppercase{Durée conseillée : #6}}
  \end{tcolorbox}
  \begin{tcolorbox}[enhanced,colback=white,colframe=popink,boxrule=2.2pt,arc=10pt,
      drop shadow=popink,left=16pt,right=16pt,top=10pt,bottom=10pt,after skip=8pt]
    \poppill{Dans cette leçon}{poppurple}{white}\par\vspace{5pt}
    \popsemi\fontsize{9.3}{12.6}\selectfont\color{popink}#7
  \end{tcolorbox}
  \noindent\begin{minipage}[t]{0.487\linewidth}
    \begin{tcolorbox}[enhanced,colback=popgreen,colframe=popink,boxrule=2.2pt,arc=10pt,
        drop shadow=popink,left=11pt,right=11pt,top=10pt,bottom=10pt,height=3.1cm,valign=center]
      \tcbox[colback=white,colframe=popink,boxrule=1.3pt,arc=5pt,boxsep=0pt,nobeforeafter,
        left=3pt,right=3pt,top=3pt,bottom=3pt,valign=center]{\qrcode[height=1.9cm]{https://play.google.com/store/apps/details?id=com.luvvix.onbuch&hl=fr}}%
      \hspace{8pt}\begin{minipage}[c]{0.5\linewidth}
        {\popdisplay\fontsize{11}{12.5}\selectfont\color{white}\MakeUppercase{Télécharge OnBuch}}\par\vspace{4pt}
        {\raggedright\popsemi\fontsize{8.4}{11}\selectfont\color{white}Gratuit \textbullet\ Google Play\\Tuteur IA, annales, concours, résultats\par}
      \end{minipage}
    \end{tcolorbox}
  \end{minipage}\hfill
  \begin{minipage}[t]{0.487\linewidth}
    \begin{tcolorbox}[enhanced,colback=poppurple,colframe=popink,boxrule=2.2pt,arc=10pt,
        drop shadow=popink,left=11pt,right=11pt,top=10pt,bottom=10pt,height=3.1cm,valign=center]
      \tikz[baseline=-0.7ex]\node[circle,fill=white,draw=popink,line width=1.3pt,minimum size=1.6cm,inner sep=0]{\popdisplay\fontsize{24}{24}\selectfont\color{poppurple}+};%
      \hspace{8pt}\begin{minipage}[c]{0.6\linewidth}
        {\popdisplay\fontsize{11}{12.5}\selectfont\color{white}\MakeUppercase{Passe à OnBuch+}}\par\vspace{4pt}
        {\raggedright\popsemi\fontsize{8.4}{11}\selectfont\color{white}Tous les cours signés, Tuteur IA illimité, fiches PDF — onglet OnBuch+ de l'app\par}
      \end{minipage}
    \end{tcolorbox}
  \end{minipage}
  \vspace{8pt}
  \popcontenu
  \vfill
  \signatureonbuch
  \restoregeometry
  \newpage
}

% Rangée « Ce cours contient » (page de garde)
\newcommand{\poptile}[3]{% #1 couleur #2 gros texte #3 légende
  \begin{tcolorbox}[enhanced,colback=#1,colframe=popink,boxrule=2pt,arc=9pt,shadow={2.5pt}{-2.5pt}{0pt}{popink},
      left=6pt,right=6pt,top=9pt,bottom=9pt,halign=center,valign=center,height=1.75cm,width=\linewidth,nobeforeafter]
    {\popdisplay\fontsize{12.5}{13}\selectfont\color{white}\MakeUppercase{#2}}\par\vspace{3pt}
    {\centering\popsemi\fontsize{7.8}{9.5}\selectfont\color{white}#3\par}
  \end{tcolorbox}}
\newcommand{\popcontenu}{%
  \par\noindent{\popxboldls\fontsize{7.5}{7.5}\selectfont\color{popmuted}CE COURS SIGNÉ CONTIENT}\par\vspace{7pt}
  \noindent\begin{minipage}[t]{0.235\linewidth}\poptile{popblue}{Cours}{complet, illustré, pas à pas}\end{minipage}\hfill
  \begin{minipage}[t]{0.235\linewidth}\poptile{poporange}{Méthodes}{et exercices résolus}\end{minipage}\hfill
  \begin{minipage}[t]{0.235\linewidth}\poptile{poppink}{Exercices}{type examen + corrigés}\end{minipage}\hfill
  \begin{minipage}[t]{0.235\linewidth}\poptile{popgreen}{Fiche bilan}{l'essentiel à retenir}\end{minipage}\par}

% Sommaire
\newtcolorbox{sommaire}{enhanced,colback=white,colframe=popink,boxrule=2.2pt,arc=10pt,
  drop shadow=popink,left=18pt,right=18pt,top=13pt,bottom=13pt,before skip=6pt,after skip=20pt,
  before upper={\poppill{Sommaire}{poppurple}{white}\par\vspace{9pt}\popsemi\fontsize{10.6}{14.5}\selectfont\color{popink}}}

% Signature de fin de cours — poussée tout en bas de la dernière page
\newcommand{\findecours}{%
  \par\vfill\nopagebreak
  \begin{center}\popsquiggle{poporange}\end{center}\vspace{4pt}
  \signatureonbuch
}
\AtBeginDocument{\def\llbracket{\mathopen{[\mkern-3.5mu[}}\def\rrbracket{\mathclose{]\mkern-3.5mu]}}} % intervalles d'entiers (glyphe ⟦ absent de la police)
% Glyphes absents de Fira Math : redéfinis à partir de glyphes disponibles
\AtBeginDocument{%
  \def\setminus{\mathbin{\backslash}}\let\smallsetminus\setminus
  \def\vdots{\mathord{\vbox{\baselineskip3.2pt\lineskiplimit0pt\kern4pt\hbox{.}\hbox{.}\hbox{.}}}}%
  \def\ddots{\mathinner{\mkern1mu\raise6.5pt\hbox{.}\mkern2mu\raise3.6pt\hbox{.}\mkern2mu\raise0.7pt\hbox{.}\mkern1mu}}%
  \def\triangle{\mathord{\scalebox{0.9}{$\symup{\Delta}$}}}%
  \def\nabla{\mathord{\raisebox{\depth}{\scalebox{1}[-1]{$\symup{\Delta}$}}}}%
  \def\checkmark{\popcheck{popdark}}%
  \def\star{\mathbin{\ast}}\def\bigstar{\mathord{\ast}}%
  \def\diamond{\mathbin{\scalebox{0.6}{$\Diamond$}}}%
  \def\bigcup{\mathop{\mathchoice{\raisebox{-0.3em}{\scalebox{1.7}{$\cup$}}}{\scalebox{1.25}{$\cup$}}{\cup}{\cup}}\displaylimits}%
  \def\bigcap{\mathop{\mathchoice{\raisebox{-0.3em}{\scalebox{1.7}{$\cap$}}}{\scalebox{1.25}{$\cap$}}{\cap}{\cap}}\displaylimits}%
  \def\lesseqqgtr{\mathrel{\lessgtr}}\def\bigoplus{\mathop{\oplus}\displaylimits}%
}
\providecommand{\ssi}{si et seulement si} % abréviation fréquente
\AtBeginDocument{\providecommand{\wideparen}{\overparen}} % arc de cercle

%% FIGFIX-BEGIN v3 — correctifs globaux des figures (docs/onbuch-generation/tools/figfix)
\usepackage{adjustbox}\usepackage{etoolbox}
% toute figure TikZ/pgfplots plus large que la ligne est réduite à la largeur disponible
\BeforeBeginEnvironment{tikzpicture}{\begin{adjustbox}{max width=\linewidth}}
\AfterEndEnvironment{tikzpicture}{\end{adjustbox}}
% légendes pgfplots lisibles par-dessus les courbes
\pgfplotsset{every axis/.append style={legend style={fill=white,fill opacity=0.92,text opacity=1,draw=black!15,inner sep=3pt}}}
% étiquettes d'arêtes (syntaxe edge["..."]) avec fond blanc
\tikzset{every edge quotes/.append style={fill=white,inner sep=1.2pt}}
% bulles de cartes mentales : texte à la ligne dans la bulle, bulle un peu plus grande
\tikzset{every concept/.append style={text width=2.0cm,align=center,minimum size=2.3cm,inner sep=2pt}}
%% FIGFIX-END
\begin{document}

\pagegarde{Les effets des mouvements migratoires}{Géographie}{1re Tech}{Géographie – classe de Première technique}{N02}{2 séances (fiche de progression harmonisée)}{Cette leçon analyse les effets socio-économiques des migrations (exode rural, urbanisation, transferts de fonds, fuite des cerveaux) et présente les politiques nationales et internationales de gestion des flux migratoires. Elle s'appuie sur des données camerounaises et africaines actualisées et sur la construction de cartes et de diagrammes.
\vspace{4pt}
\begin{multicols}{2}\small
\begin{itemize}
\item À la fin de cette leçon, je sais définir les principaux types de migrations et leurs effets
\item identifier les problèmes socio-économiques posés par l'exode rural au Cameroun
\item expliquer les incidences positives et négatives des migrations sur les régions de départ et d'accueil
\item interpréter des données chiffrées et des cartes de flux migratoires
\item construire un croquis ou un diagramme illustrant un flux migratoire
\item décrire les tentatives de règlement des problèmes migratoires aux niveaux national, régional et international
\end{itemize}\end{multicols}}

\begin{sommaire}
\begin{multicols}{2}
\begin{enumerate}
\item Rappels et typologie des effets migratoires
\item Les problèmes socio-économiques dans les régions de départ
\item Les incidences des migrations dans les zones d'accueil
\item Les migrations internationales : gains et pertes pour le Cameroun
\item Les tentatives de règlement des problèmes migratoires
\item TP guidé : cartographier et analyser un flux migratoire
\item Activité d'intégration
\item Méthodes et exercices résolus
\item Exercices
\item Corrigés des exercices
\item Fiche bilan
\end{enumerate}
\end{multicols}
\end{sommaire}

\begin{objectifs}
\begin{itemize}
\item À la fin de cette leçon, je sais définir les principaux types de migrations et leurs effets
\item identifier les problèmes socio-économiques posés par l'exode rural au Cameroun
\item expliquer les incidences positives et négatives des migrations sur les régions de départ et d'accueil
\item interpréter des données chiffrées et des cartes de flux migratoires
\item construire un croquis ou un diagramme illustrant un flux migratoire
\item décrire les tentatives de règlement des problèmes migratoires aux niveaux national, régional et international
\item rédiger et justifier une conclusion argumentée sur un cas migratoire concret
\item appliquer les notions de la leçon à une situation de la vie courante (départ d'un proche, installation en ville)
\end{itemize}
\end{objectifs}

\begin{prerequis}
\begin{itemize}
\item Notion de migration (définitions : émigration, immigration, migration interne/internationale) vue dans l'unité précédente
\item Les causes des mouvements migratoires (facteurs de répulsion et d'attraction)
\item Croissance démographique et répartition de la population du Cameroun (classe de 4e et début d'année)
\item Lecture d'une carte et d'un tableau statistique
\item Notions d'urbanisation et d'exode rural
\end{itemize}
\end{prerequis}

\newpage

\coursec{Rappels et typologie des effets migratoires}

Chaque année, des dizaines de milliers de jeunes quittent les villages du Nord, de l'Adamaoua et de l'Ouest camerounais pour Douala et Yaoundé, tandis que d'autres tentent l'aventure vers l'Europe, l'Amérique du Nord ou les pays du Golfe. Selon les projections de l'INS issues du RGPH (Recensement Général de la Population et de l'Habitat), la ville de Douala croît d'environ \num{100000} habitants par an, une part significative étant attribuable au solde migratoire positif. Mais que gagnent et que perdent réellement les régions de départ et d'accueil ? Comment ces mouvements transforment-ils l'économie, les familles, l'espace urbain et les équilibres politiques ? Et que font les États et les organisations internationales pour encadrer ces flux ? Cette leçon répond à ces questions à partir de documents concrets et de données chiffrées.

\textbf{Problématique :} quels sont les effets des mouvements migratoires sur les régions de départ et sur les régions d'accueil, et comment les États tentent-ils de régler les problèmes qu'ils posent ?

Dans cette première partie, nous posons les bases : les définitions essentielles, puis une \cle{typologie} (classification) des effets migratoires, enfin une grille d'analyse qui servira de méthode pour toute la leçon.

\courssub{Rappel des définitions fondamentales}

Avant d'analyser les effets, il faut s'accorder sur le vocabulaire précis. L'intuition est simple : une personne quitte un lieu pour s'installer durablement ailleurs. Mais ce déplacement ne devient une \cle{migration} au sens géographique que s'il respecte deux conditions cumulatives : il \textbf{franchit une limite spatiale} (limite de commune, de département, de région ou d'État) et il aboutit à un \textbf{séjour durable} (généralement plus de 6 mois, selon la définition de l'ONU), ce qui le distingue du simple voyage, du tourisme ou du déplacement quotidien (navettage).

\begin{definition}[Migration]
Une \cle{migration} est un déplacement de personnes qui franchit une limite administrative ou étatique en vue d'un séjour durable dans le lieu d'arrivée. On distingue :
\begin{itemize}
\item l'\cle{émigration} : le départ d'une personne hors de son espace d'origine (point de vue de la région de départ) ;
\item l'\cle{immigration} : l'arrivée d'une personne dans un espace d'accueil (point de vue de la région d'accueil) ;
\item la \cle{migration interne} : le déplacement à l'intérieur d'un même pays (exode rural, migration interurbaine, migration vers les fronts pionniers) ;
\item la \cle{migration internationale} : le déplacement d'un pays vers un autre pays.
\end{itemize}
\end{definition}

\begin{attention}[Un même migrant, deux points de vue]
Un jeune qui quitte Bafoussam (Ouest) pour Douala (Littoral) est \emph{émigrant} pour la région de l'Ouest et \emph{immigrant} pour la région du Littoral : c'est la même personne, vue des deux bouts du trajet. Ne jamais opposer émigration et immigration comme s'il s'agissait de deux catégories de personnes différentes ; ce sont deux regards sur un même fait.
\end{attention}

Au Cameroun, les migrations internes dominent largement en volume : l'\cle{exode rural} (départ des campagnes vers les villes) et les migrations \emph{rural-rural} (vers les zones de plantation ou les fronts pionniers de l'Est et de l'Adamaoua) alimentent la croissance urbaine. Les migrations internationales concernent un effectif plus restreint, mais souvent hautement qualifié (étudiants, cadres de santé, ingénieurs, informaticiens), ce qui pose avec acuité le problème de la \cle{fuite des cerveaux} (\emph{brain drain}).

\courssub{Les grandes catégories d'effets des migrations}

Observer un flux migratoire, c'est observer une transformation simultanée des deux territoires qu'il relie. Pour classer ces transformations, les géographes distinguent cinq natures d'effets, qui s'entrecroisent dans la réalité.

\begin{definition}[Typologie des effets migratoires]
Les effets des migrations se classent en cinq catégories :
\begin{itemize}
\item \textbf{Effets démographiques} : modification du nombre, de la structure par âge et par sexe de la population (vieillissement ou rajeunissement, déséquilibre du ratio de masculinité, croissance ou déclin démographique local) ;
\item \textbf{Effets économiques} : envois d'argent (\emph{transferts de fonds} ou \emph{remittances}), perte ou gain de main-d'œuvre (qualifiée ou non), fuite des cerveaux / gain de cerveaux (\emph{brain gain}), mise en valeur de nouveaux espaces agricoles, investissements productifs des migrants de retour ;
\item \textbf{Effets sociaux} : recomposition des familles (familles monoparentales, enfants « laissés-pour-compte »), évolution du niveau d'instruction (scolarisation financée par la diaspora), surpeuplement urbain, pression sur les services sociaux (santé, éducation, eau), tensions ou solidarités nouvelles (associations de développement villageois) ;
\item \textbf{Effets culturels} : diffusion des langues, des musiques, des cuisines, des modes vestimentaires, des pratiques religieuses ; brassage interculturel ou, à l'inverse, repli identitaire et xénophobie ;
\item \textbf{Effets politiques} : modification du poids électoral des circonscriptions, élaboration de politiques d'immigration/émigration, diplomatie migratoire entre États (accords de réadmission, gestion des réfugiés), rôle politique des diasporas (vote à l'étranger, financement de partis).
\end{itemize}
\end{definition}

\begin{exemplebox}[Des effets visibles dans la vie quotidienne au Cameroun]
\begin{itemize}
\item \emph{Démographique} : le département du Mayo-Danay (Extrême-Nord) perd chaque année une partie de sa jeunesse active au profit des villes du Sud, accélérant le vieillissement de la population villageoise.
\item \emph{Économique} : les transferts de fonds de la diaspora camerounaise (estimés à plus de \num{300} milliards de FCFA/an, soit près de 2\% du PIB) financent la construction de maisons, la scolarisation, la santé et le petit commerce au village.
\item \emph{Social} : à Douala, les quartiers périphériques (PK 12, PK 18, Makea, Ndogpassi) s'étendent sous l'effet de l'arrivée continue de migrants, souvent dans la précarité (habitat spontané).
\item \emph{Culturel} : le \emph{Ngondo} (fête du peuple Sawa) et les danses de l'Ouest (Bikutsi, Makossa modernisé) circulent et se recomposent dans les métropoles grâce aux migrants ; le \emph{Camfranglais} s'impose comme langue de rue transrégionale.
\item \emph{Politique} : l'accueil de plus de \num{400000} réfugiés centrafricains (Est, Adamaoua) et nigérians (Extrême-Nord, Nord) pèse sur les relations diplomatiques et le budget national, avec l'appui du HCR et de l'OIM.
\end{itemize}
\end{exemplebox}

\courssub{Régions de départ, régions d'accueil : une grille d'analyse}

L'idée directrice de la leçon est que \textbf{tout effet migratoire se lit à deux endroits} : là d'où l'on part (région de départ) et là où l'on arrive (région d'accueil). Le départ d'un jeune actif est une \emph{perte} de bras et de compétences pour le village mais un \emph{gain} de main-d'œuvre pour la ville. Réciproquement, l'argent qu'il envoie au village est un \emph{gain} financier pour la région de départ. La grille d'analyse croise donc trois questions fondamentales : \emph{Où ?} (départ ou accueil), \emph{Quoi ?} (nature de l'effet) et \emph{Quel signe ?} (positif / gain ou négatif / perte / problème).

\begin{methode}[Analyser un effet migratoire en trois étapes]
\begin{enumerate}
\item \textbf{Identifier l'espace concerné} : l'effet se manifeste-t-il dans la région de départ ou dans la région d'accueil ?
\item \textbf{Déterminer la nature de l'effet} : est-il démographique, économique, social, culturel ou politique ?
\item \textbf{Préciser son signe et le justifier} : s'agit-il d'un effet positif (gain, avantage) ou négatif (perte, coût, problème) ? Apporter un fait précis, un chiffre ou un exemple localisé.
\end{enumerate}
\end{methode}

\begin{popfigure}
\begin{center}
\begin{tikzpicture}[font=\small]
% Colonne départ
\draw[fill=poporangeL,draw=poporange,thick,rounded corners] (0,0) rectangle (5.6,6.4);
\node[font=\bfseries\color{popdark}] at (2.8,6.0) {Région de départ};
\node[align=left,anchor=north west,text width=5.2cm] at (0.2,5.5) {
\textcolor{popgreen}{\textbf{+}} Transferts de fonds (remittances)\\[2pt]
\textcolor{popgreen}{\textbf{+}} Baisse du chômage / pression foncière\\[2pt]
\textcolor{popgreen}{\textbf{+}} Retour de compétences, investissements\\[2pt]
\textcolor{popink}{\textbf{--}} Perte de jeunes actifs (force de travail)\\[2pt]
\textcolor{popink}{\textbf{--}} Vieillissement, féminisation de l'agriculture\\[2pt]
\textcolor{popink}{\textbf{--}} Champs délaissés, insécurité alimentaire\\[2pt]
\textcolor{popink}{\textbf{--}} Familles éclatées, enfants vulnérables};
% Colonne accueil
\draw[fill=popblueL,draw=popblue,thick,rounded corners] (9.4,0) rectangle (15,6.4);
\node[font=\bfseries\color{popdark}] at (12.2,6.0) {Région d'accueil};
\node[align=left,anchor=north west,text width=5.2cm] at (9.6,5.5) {
\textcolor{popgreen}{\textbf{+}} Main-d'œuvre abondante, peu coûteuse\\[2pt]
\textcolor{popgreen}{\textbf{+}} Dynamisme entrepreneurial, marchés\\[2pt]
\textcolor{popgreen}{\textbf{+}} Brassage culturel, innovation\\[2pt]
\textcolor{popink}{\textbf{--}} Surpeuplement, habitat précaire (bidonvilles)\\[2pt]
\textcolor{popink}{\textbf{--}} Pression sur emplois, salaires, services\\[2pt]
\textcolor{popink}{\textbf{--}} Tensions sociales, insécurité\\[2pt]
\textcolor{popink}{\textbf{--}} Pollution, embouteillages, gestion déchets};
% Flèches
\draw[-{Stealth[length=4mm]},line width=2.5pt,popink] (5.8,4.8) -- (9.2,4.8) node[midway,above,font=\footnotesize\bfseries]{Flux de personnes (départs)};
\draw[-{Stealth[length=4mm]},line width=2.5pt,popgreen] (9.2,1.6) -- (5.8,1.6) node[midway,below,font=\footnotesize\bfseries]{Flux de retour (argent, biens, compétences)};
\end{tikzpicture}
\end{center}
\legende{Schéma en double colonne des effets migratoires. La flèche orange représente le flux de personnes (départs) ; la flèche verte représente les flux de retour (transferts d'argent, investissements, retours de migrants qualifiés). Chaque espace cumule simultanément des effets positifs (vert) et négatifs (rose).}
\end{popfigure}

\begin{attention}[Erreur fréquente au Probatoire : confondre effets pour le migrant et effets pour les territoires]
Dans les copies, beaucoup d'élèves écrivent : « La migration permet au migrant de trouver un emploi et d'améliorer ses conditions de vie » alors que la question porte sur les effets \emph{pour les régions} (départ ou accueil). L'effet pour l'individu (trouver un emploi, gagner mieux sa vie, s'instruire) n'est \textbf{pas} l'effet pour le territoire (gagner/perdre de la main-d'œuvre, recevoir des transferts, subir le surpeuplement). Toujours se demander : \emph{Qui} subit ou bénéficie de l'effet — la personne physique ou l'espace géographique ?
\end{attention}

\courssub{Effets positifs, effets négatifs : la nécessité d'un bilan nuancé}

Un piège classique consiste à voir les migrations comme un mal absolu (misère, bidonvilles, déracinement) ou comme un bien absolu (richesse, développement, modernité). La réalité est \cle{nuancée} : un même flux produit simultanément des gains et des pertes, souvent répartis inégalement selon les groupes sociaux et les échelles de temps. L'exode rural soulage la pression démographique sur les terres du village (positif à court terme) mais le prive de ses forces vives, risquant d'hypothéquer son avenir agricole (négatif structurel) ; la ville gagne des travailleurs dynamiques (positif économique) mais voit exploser le surpeuplement et le déficit d'infrastructures (négatif social/environnemental). Un bon raisonnement de géographe \textbf{pèse les deux faces} avant de conclure, et précise toujours \emph{pour qui}, \emph{à quelle échelle} (locale, nationale) et \emph{à quel horizon temporel} (court / long terme) l'effet joue.

\begin{exemplebox}[Douala et Yaoundé, aimants démographiques du Cameroun]
Selon les projections de l'INS (années 2020-2023), Douala (environ \num{3,9} millions d'habitants) et Yaoundé (environ \num{4,3} millions) concentrent à elles deux près de \textbf{30\% de la population urbaine} et plus d'\textbf{un quart de la population totale} du Cameroun (estimée à \num{28,6} millions en 2023, source : World Bank / INS). Ce poids démographique exceptionnel s'explique en grande partie par les migrations internes : les deux métropoles attirent les jeunes de toutes les régions par l'emploi (secteur informel dominant), les études (universités, grandes écoles) et les services de santé. C'est la preuve concrète que les effets des migrations ne sont pas abstraits : ils redessinent la carte du peuplement national et posent le défi de l'urbanisation rapide.
\end{exemplebox}

\courssub{Mise en évidence à partir d'un document : la carte des flux internes}

Pour vérifier ces idées, lisons un document géographique : le croquis schématique des principaux flux migratoires internes au Cameroun. La méthode de lecture d'un croquis de flux est standardisée : repérer les \textbf{pôles d'attraction} (où convergent les flèches), les \textbf{bassins de départ} (d'où partent les flèches), l'\textbf{intensité} des flux (épaisseur des flèches) et les \textbf{obstacles ou relais} (montagnes, routes, villes-étapes).

\begin{popfigure}
\begin{center}
\begin{tikzpicture}[font=\footnotesize,scale=1.1]
% Contour très simplifié du Cameroun (approximatif)
\draw[fill=popcream,draw=popdark,thick] 
  (3.2,6.6) -- (4.6,6.2) -- (5.2,5.2) -- (6.6,4.8) -- (7.0,3.6) 
  -- (6.2,2.4) -- (6.6,1.2) -- (5.4,0.2) -- (3.4,0.0) -- (2.0,0.8) 
  -- (1.4,2.0) -- (0.6,2.6) -- (1.0,3.8) -- (2.0,4.4) -- (2.4,5.6) -- cycle;
% Villes principales (points)
\fill[popink] (1.3,1.9) circle (2.5pt) node[below left=2pt,font=\bfseries\small]{Douala};
\fill[popink] (3.4,1.5) circle (2.5pt) node[below=2pt,font=\bfseries\small]{Yaoundé};
\fill[popmuted] (2.6,3.4) circle (2pt) node[above left=1pt,font=\small]{Bafoussam};
\fill[popmuted] (2.1,4.3) circle (2pt) node[left=1pt,font=\small]{Bamenda};
\fill[popmuted] (3.6,5.6) circle (2pt) node[above=1pt,font=\small]{Ngaoundéré};
\fill[popmuted] (4.3,6.3) circle (2pt) node[above=1pt,font=\small]{Maroua};
\fill[popmuted] (5.6,2.6) circle (2pt) node[right=1pt,font=\small]{Bertoua};
% Flux principaux (flèches d'épaisseur proportionnelle à l'importance relative)
% Ouest -> Douala / Yaoundé (flux très forts)
\draw[-{Stealth[length=3.5mm]},line width=3.0pt,poporange] (2.6,3.3) -- (1.5,2.1);
\draw[-{Stealth[length=3.5mm]},line width=2.5pt,poporange] (2.6,3.3) -- (3.3,1.6);
% Nord-Ouest -> Douala / Yaoundé
\draw[-{Stealth[length=3.5mm]},line width=2.5pt,poporange] (2.15,4.2) -- (1.4,2.1);
\draw[-{Stealth[length=3.5mm]},line width=2.0pt,poporange] (2.15,4.2) -- (3.35,1.65);
% Adamaoua -> Yaoundé (fort)
\draw[-{Stealth[length=3.5mm]},line width=2.8pt,poporange] (3.55,5.5) -- (3.4,1.8);
% Nord -> Yaoundé (fort)
\draw[-{Stealth[length=3.5mm]},line width=2.5pt,poporange] (4.25,6.2) -- (3.5,1.8);
% Est -> Yaoundé / Douala (modéré)
\draw[-{Stealth[length=3.5mm]},line width=1.8pt,poporange] (5.5,2.55) -- (3.6,1.6);
\draw[-{Stealth[length=3.5mm]},line width=1.5pt,poporange] (5.5,2.55) -- (1.5,2.0);
% Nord -> Douala (transit via Yaoundé ou direct, modéré)
\draw[-{Stealth[length=3mm]},popdark,thick] (7.4,5.8) -- (7.4,6.5) node[above]{N};
% Échelle graphique approximative
\draw[popdark,thick] (0.5,0.3) -- (1.5,0.3) node[midway,below=1pt,font=\tiny]{~200 km};
\end{tikzpicture}
\end{center}
\legende{Croquis simplifié des principaux flux migratoires internes au Cameroun (contour approximatif). Les flèches orange (épaisseur proportionnelle à l'importance relative du flux) montrent la convergence des migrants de l'Ouest, du Nord-Ouest, de l'Adamaoua, du Nord et de l'Est vers les deux métropoles, Douala (capitale économique, port) et Yaoundé (capitale politique, administrations, universités). Les régions de l'Adamaoua et de l'Est jouent souvent un rôle d'étape (villes-relais).}
\end{popfigure}

\textbf{Lecture commentée du document.} Trois enseignements majeurs se dégagent :
\begin{enumerate}
\item Les flèches \textbf{convergent} massivement vers deux pôles : Douala (capitale économique, port autonome, zone industrielle de Bassa/Bonabéri) et Yaoundé (capitale politique, administrations centrales, pôles universitaires). Ce sont les deux \cle{pôles d'attraction} majeurs du système spatial camerounais.
\item Les flux les plus épais partent de l'\textbf{Ouest} (région la plus dense, pression foncière forte, tradition migratoire) et du \textbf{Nord} (pauvreté rurale, insécurité liée à Boko Haram, aléa climatique). Ce sont les principaux \cle{bassins de départ} (ou d'expulsion).
\item Les régions intermédiaires (Adamaoua, Est, Centre hors Yaoundé) jouent parfois un rôle d'\textbf{étape} (villes-relais comme Ngaoundéré, Bertoua) ou de \textbf{redistribution} : les migrants y séjournent quelques mois/années avant de repartir vers les métropoles.
\end{enumerate}

\begin{exoresolu}[Appliquer la grille d'analyse à un fait de la carte]
Énoncé : Sur le croquis, le flux Bafoussam $\rightarrow$ Douala est l'un des plus épais. En utilisant la méthode en trois étapes, analysez deux effets de ce flux : un pour la région de départ (Ouest), un pour la région d'accueil (Littoral / Douala).
\tcblower
Solution détaillée :
\begin{enumerate}
\item \textbf{Effet pour la région de départ (Ouest / Bafoussam).}
\begin{itemize}
\item \emph{Étape 1 (Où) :} Région de départ (département de la Mifi, région de l'Ouest).
\item \emph{Étape 2 (Nature) :} Effet \textbf{démographique} et \textbf{économique}.
\item \emph{Étape 3 (Signe + Justification) :} \textbf{Négatif} à court/moyen terme : départ massif des 15-35 ans (structure par âge creusée, vieillissement, féminisation de l'agriculture vivrière), terres en jachère faute de bras. \textbf{Positif} à moyen/long terme : transferts de fonds massifs (construction en dur, scolarisation, santé, petit commerce), financement des associations de développement (AD) qui réalisent des forages, routes rurales, salles de classe.
\end{itemize}
\item \textbf{Effet pour la région d'accueil (Littoral / Douala).}
\begin{itemize}
\item \emph{Étape 1 (Où) :} Région d'accueil (ville de Douala, département du Wouri).
\item \emph{Étape 2 (Nature) :} Effet \textbf{économique}, \textbf{social} et \textbf{spatial}.
\item \emph{Étape 3 (Signe + Justification) :} \textbf{Positif} : main-d'œuvre nombreuse, jeune, peu coûteuse et entreprenante (commerce informel, BTP, services, industrie portuaire) qui fait tourner l'économie urbaine. \textbf{Négatif} : arrivées massives non planifiées alimentant l'étalement urbain précaire (quartiers PK, Makea, Ndogpassi), saturation des voiries, déficit en eau/électricité/logements, tensions foncières avec populations autochtones (Bakoko, Duala).
\end{itemize}
\end{enumerate}
\textbf{Conclusion (bilan nuancé) :} Le même flux Bafoussam-Douala est simultanément une perte de capital humain pour le village et un gain de main-d'œuvre pour la métropole ; un coût social (familles éclatées) et une source de revenus (transferts). L'analyse géographique refuse le manichéisme : elle met en évidence la \textbf{complémentarité conflictuelle} des deux espaces.
\end{exoresolu}

\begin{aretenir}[L'essentiel de la section 1]
\begin{itemize}
\item Une migration est un déplacement franchissant une limite administrative pour un séjour durable ($\ge$ 6 mois) ; elle est interne ou internationale, et se lit comme émigration (point de vue départ) ou immigration (point de vue accueil).
\item Les effets migratoires sont de cinq natures : démographiques, économiques, sociaux, culturels, politiques.
\item Tout effet s'analyse en trois questions : \textbf{Où ?} (départ / accueil) — \textbf{Quoi ?} (nature) — \textbf{Quel signe ?} (positif / négatif, justifié).
\item Un bilan rigoureux est toujours \textbf{nuancé} : un flux produit des gains et des pertes des deux côtés ; il distingue les effets pour les \textbf{territoires} (objectif du géographe) des effets pour les \textbf{migrants} (trajectoire individuelle).
\item Au Cameroun, Douala et Yaoundé (ensemble $\approx$ 8,2 millions d'habitants, $>$ 25\% de la population nationale) polarisent les flux internes venus surtout de l'Ouest (densité, tradition) et du Nord (pauvreté, insécurité).
\end{itemize}
\end{aretenir}

\begin{exercice}[Application : À vous de jouer]
À partir du croquis des flux internes (figure ci-dessus), choisissez \textbf{un flux de votre choix} (ex. : Maroua $\rightarrow$ Yaoundé ; Bamenda $\rightarrow$ Douala ; Bertoua $\rightarrow$ Yaoundé).
Rédigez un paragraphe structuré (12-15 lignes) analysant, \textbf{avec la méthode en trois étapes}, \textbf{un effet positif et un effet négatif pour la région de départ} ET \textbf{un effet positif et un effet négatif pour la région d'accueil}.
Concluez par une phrase de bilan nuancé reprenant le vocabulaire : « complémentarité », « inégalités », « aménagement du territoire ».
\end{exercice}

\coursec{Les problèmes socio-économiques dans les régions de départ}

Dans la section précédente, nous avons rappelé les grandes catégories de migrations (migrations internes et internationales, exode rural, migrations de travail) et classé leurs effets en deux familles : effets dans les \cle{zones de départ} et effets dans les \cle{zones d'accueil}. Nous entrons maintenant dans le cœur de la leçon : que devient un village, une région, lorsqu'une partie importante de sa population s'en va ? C'est la question que posent chaque année des milliers de familles du Nord, de l'Est et de l'Ouest camerounais lorsque leurs jeunes partent vers Douala, Yaoundé ou l'étranger. Pour y répondre, nous procéderons comme un géographe : observation d'abord, explication des mécanismes ensuite, nuances enfin.

\courssub{Le dépeuplement et le vieillissement des campagnes}

\begin{definition}[Exode rural]
L'\cle{exode rural} est le départ \textbf{massif et durable} des populations rurales vers les villes. Il se distingue d'un simple déplacement temporaire (saisonnier) : celui qui part en exode rural s'installe en ville et ne revient au village que pour des visites.
\end{definition}

\textbf{Intuition.} Imaginez une classe de Première dans un lycée de village : à la fin de l'année, les trois quarts des élèves annoncent qu'ils partent « tenter leur chance » à Douala ou à Yaoundé. Dix ans plus tard, qui reste au village ? Les personnes âgées, les jeunes enfants, et quelques adultes. Le village n'a pas disparu, mais il a perdu sa force vive : c'est exactement ce que montre la géographie.

\textbf{Le mécanisme.} Les migrants sont rarement choisis au hasard : ce sont presque toujours les \cle{jeunes actifs} (15-35 ans), les plus instruits, les plus dynamiques, car ce sont eux qui ont la force de voyager et l'espoir de trouver un emploi. Leur départ produit deux effets combinés :

\begin{itemize}
\item le \cle{dépeuplement} : la population totale du village diminue ou cesse de croître ;
\item le \cle{vieillissement} : la proportion de personnes âgées augmente mécaniquement, puisque les jeunes partent et que les anciens restent.
\end{itemize}

\begin{exemplebox}[Un fait chiffré]
Dans certains villages de la région du \textbf{Nord-Ouest} du Cameroun, plus de \textbf{60 \%} des jeunes de 15 à 35 ans ont émigré vers les villes (Bamenda, Douala, Yaoundé) ou vers l'étranger. Dans ces villages, les champs sont de plus en plus cultivés par des personnes de plus de 55 ans.
\end{exemplebox}

Pour visualiser ce déséquilibre, comparons la structure par âge d'un village rural touché par l'exode et celle de la ville de Douala qui accueille les migrants.

\begin{popfigure}
\begin{center}
\begin{tikzpicture}
\begin{axis}[
  xbar, xmin=0, xmax=30,
  width=0.85\linewidth, height=7.5cm,
  xlabel={Part de la population (en \%)},
  symbolic y coords={75 ans et +,60-74 ans,45-59 ans,30-44 ans,15-29 ans,0-14 ans},
  ytick=data,
  nodes near coords, nodes near coords align={horizontal},
  legend style={at={(0.5,-0.22)},anchor=north,legend columns=2},
  bar width=6pt
]
\addplot[fill=popgreen!70,draw=popgreen] coordinates {(22,0-14 ans) (12,15-29 ans) (14,30-44 ans) (16,45-59 ans) (20,60-74 ans) (16,75 ans et +)};
\addplot[fill=popblue!70,draw=popblue] coordinates {(26,0-14 ans) (28,15-29 ans) (22,30-44 ans) (13,45-59 ans) (8,60-74 ans) (3,75 ans et +)};
\legend{Village rural (zone de départ), Douala (zone d'accueil)}
\end{axis}
\end{tikzpicture}
\end{center}
\legende{Pyramide des âges comparée (schématique) : dans le village de départ, les 15-29 ans ne représentent que 12 \% de la population contre 28 \% à Douala ; à l'inverse, les personnes âgées dominent au village. Les jeunes actifs ont « migré » d'une pyramide vers l'autre.}
\end{popfigure}

\textbf{Lecture du graphique.} La barre des 15-29 ans est le témoin de l'exode : elle est « creusée » au village et « gonflée » en ville. Le géographe dit que la migration est \cle{sélective} : elle sélectionne les âges, et donc les forces de travail.

\courssub{La baisse de la main-d'œuvre agricole et l'abandon des terres}

\textbf{Observation.} Dans de nombreux villages de l'Adamaoua ou de l'Est, on observe des champs envahis par les herbes, des plantations de café ou de cacao non entretenues, des récoltes qui pourrissent faute de bras pour les ramasser.

\textbf{Le mécanisme.} L'agriculture camerounaise est encore largement \cle{manuelle} : houe, machette, portage. Elle exige donc une main-d'œuvre nombreuse et jeune, surtout aux périodes de défrichage, de semis et de récolte. Quand les jeunes partent :

\begin{enumerate}
\item la \textbf{quantité} de travail disponible diminue : les surfaces cultivées se réduisent ;
\item la \textbf{qualité} du travail baisse : les personnes âgées ne peuvent plus défricher ni cultiver les terres lointaines ; on cultive seulement les champs proches des cases ;
\item certaines terres sont purement et simplement \textbf{abandonnées} : la brousse reprend ses droits.
\end{enumerate}

\begin{exemplebox}[Conséquence en chaîne]
Dans un village de la région de l'Est, le départ des jeunes a entraîné l'abandon de champs de cacao vieillissants : moins de main-d'œuvre $\rightarrow$ moins d'entretien $\rightarrow$ baisse des rendements $\rightarrow$ baisse des revenus $\rightarrow$ nouveaux départs. C'est un \cle{cercle vicieux} : la pauvreté provoque la migration, qui aggrave la pauvreté.
\end{exemplebox}

Pour mesurer qui part exactement, une enquête menée dans un village de 500 habitants donne la répartition suivante des 80 départs de l'année :

\begin{popfigure}
\begin{center}
\begin{tikzpicture}
\begin{axis}[
  width=0.7\linewidth, height=7cm,
  ybar, ymin=0, ymax=60,
  ylabel={Nombre de départs},
  symbolic x coords={Moins de 15 ans,15-25 ans,26-40 ans,41-60 ans,Plus de 60 ans},
  xtick=data, x tick label style={rotate=20,anchor=east,font=\small},
  nodes near coords,
  bar width=18pt
]
\addplot[fill=poporange!75,draw=poporange] coordinates {(Moins de 15 ans,5) (15-25 ans,38) (26-40 ans,24) (41-60 ans,10) (Plus de 60 ans,3)};
\end{axis}
\end{tikzpicture}
\end{center}
\legende{Répartition des 80 départs d'un village par tranche d'âge (enquête scolaire fictive mais réaliste) : les 15-40 ans représentent à eux seuls 62 départs sur 80, soit \textbf{77,5 \%} des migrants. La migration vide le village de ses forces vives.}
\end{popfigure}

\textbf{Calcul.} Part des 15-40 ans dans les départs : $\dfrac{38+24}{80}\times 100 = 77,5\,\%$. Cette sélectivité par l'âge est la clé de tous les problèmes qui suivent.

\courssub{La désorganisation des familles}

\textbf{Intuition.} La migration ne déplace pas seulement des travailleurs : elle déplace des pères, des mères, des fils. Toute la \cle{structure familiale} s'en trouve modifiée.

\textbf{Les situations concrètes observées au Cameroun :}

\begin{itemize}
\item \textbf{Les femmes chefs de ménage.} Quand l'homme part travailler en ville ou à l'étranger, sa femme reste au village et assume seule la direction du foyer : culture des champs, éducation des enfants, gestion de l'argent. On parle de \cle{féminisation} de l'agriculture et des responsabilités rurales. Cette charge supplémentaire est lourde, même si elle donne parfois aux femmes une nouvelle autonomie.
\item \textbf{Les enfants confiés aux grands-parents.} Quand les deux parents partent, les enfants sont souvent laissés aux grands-parents, parfois âgés et fatigués. Le suivi scolaire et l'éducation peuvent en souffrir ; certains enfants abandonnent l'école pour aider aux champs.
\item \textbf{L'éclatement des familles élargies.} La solidarité traditionnelle (travail collectif, entraide lors des deuils et des mariages) s'affaiblit quand les membres de la famille sont dispersés entre le village, Douala, Yaoundé et l'étranger.
\end{itemize}

\begin{attention}[Piège de rédaction]
Ne présentez jamais la famille rurale comme « détruite » par la migration : elle est \textbf{réorganisée}. Les liens se maintiennent par le téléphone mobile, les transferts d'argent et les retours au village lors des fêtes. Au Probatoire, employez le vocabulaire précis : \emph{désorganisation}, \emph{recomposition}, \emph{femmes chefs de ménage}, et non des jugements moraux.
\end{attention}

\courssub{La perte des élites rurales et l'affaiblissement du développement local}

\textbf{Le mécanisme.} Qui part en premier ? Souvent ceux qui ont réussi à l'école : les diplômés, les jeunes formés, les plus entreprenants. Le village perd ainsi ses \cle{élites} potentielles : futurs instituteurs, infirmiers, commerçants, élus locaux, leaders d'associations. Les géographes parlent de \cle{fuite des cerveaux} lorsqu'il s'agit de diplômés qui partent à l'étranger ; à l'échelle interne, on parle d'\cle{exode des élites rurales}.

\textbf{Les conséquences sur le développement local :}

\begin{itemize}
\item moins d'initiatives locales (coopératives, GIC, projets communautaires) faute d'animateurs ;
\item difficulté à trouver des jeunes pour les travaux communautaires (entretien des pistes, des points d'eau) ;
\item baisse des recettes des communes rurales (moins de contribuables, moins de marchés dynamiques) ;
\item découragement des investisseurs : qui ouvrirait un atelier ou un commerce dans un village qui se vide ?
\end{itemize}

\begin{exemplebox}[Cas des régions du Nord, de l'Est et de l'Ouest]
\begin{itemize}
\item \textbf{Régions du Nord (Nord, Extrême-Nord, Adamaoua)} : forts départs vers Garoua, Ngaoundéré, puis vers Douala et Yaoundé ; causes : sécheresses, pauvreté, faible industrialisation, insécurité dans l'Extrême-Nord. Conséquence : villages vieillis, champs de coton et de mil parfois sous-exploités.
\item \textbf{Région de l'Est} : exode vers Yaoundé et vers les zones minières ; villages enclavés qui perdent leurs jeunes, cacao et café peu entretenus.
\item \textbf{Région de l'Ouest} : terres denses et morcelées ; les jeunes, faute de terres suffisantes, partent massivement vers Douala et les villes du Littoral ; les hautes terres perdent leurs forces vives malgré une agriculture dynamique.
\end{itemize}
\end{exemplebox}

\courssub{Le contrepoint positif : un regard équilibré}

\begin{attention}[Ne pas noircir le tableau]
L'exode rural n'est pas \textbf{uniquement négatif}. Un bon devoir de géographie présente toujours les deux faces du phénomène : dire seulement « l'exode rural est une catastrophe » est une erreur d'analyse qui coûte des points.
\end{attention}

\textbf{Les effets positifs pour les régions de départ :}

\begin{enumerate}
\item \textbf{L'allègement de la pression sur les terres.} Dans les régions densément peuplées comme l'Ouest, le départ d'une partie des jeunes réduit la concurrence pour les terres : ceux qui restent disposent de parcelles plus grandes, et les conflits fonciers diminuent.
\item \textbf{La réduction du chômage rural.} En partant, les jeunes en surnombre ne restent pas sans emploi au village : le chômage et le sous-emploi ruraux diminuent.
\item \textbf{Les transferts de fonds.} Les migrants envoient de l'argent au village : ce \cle{transfert de fonds} finance la scolarité des enfants, les soins de santé, la construction de maisons modernes, parfois l'achat de matériel agricole. Certaines maisons en dur des villages du Nord-Ouest ou de l'Ouest sont payées par des fils partis à Douala ou en Europe.
\end{enumerate}

\begin{savaistu}[La diaspora camerounaise, une banque pour le pays]
Les transferts d'argent des migrants camerounais de la diaspora (Europe, Amérique du Nord, Afrique du Sud, Golfe) dépassent \textbf{200 milliards de FCFA par an} (estimations Banque mondiale, 2023). C'est plus que le budget de nombreux ministères ! Une partie de cet argent revient directement aux villages : frais scolaires, forages, églises, maisons. La migration, bien gérée, peut donc devenir un \emph{levier de développement} pour les régions de départ.
\end{savaistu}

\begin{exoresolu}[Analyser un tableau de données villageoises]
\textbf{Énoncé.} Un village de l'Adamaoua comptait \num{1200} habitants en 2010. En 2024, il n'en compte plus que 900. Sur les 300 départs, 210 concernent des jeunes de 15 à 35 ans. Par ailleurs, les migrants du village installés à Douala et à Yaoundé y ont envoyé 45 millions de FCFA en 2024, utilisés ainsi : 60 \% pour la consommation et la scolarité, 25 \% pour la construction, 15 \% pour l'agriculture.
1) Calculez le taux de dépeuplement du village entre 2010 et 2024.
2) Calculez la part des jeunes dans les départs et commentez.
3) Calculez le montant investi dans l'agriculture et concluez sur le rôle ambivalent (à la fois négatif et positif) de la migration pour ce village.
\tcblower
\textbf{Solution détaillée.}
1) Taux de variation de la population :
\[ \frac{900-1200}{1200}\times 100 = -25\,\%. \]
Le village a perdu \textbf{25 \%} de sa population en 14 ans : c'est un dépeuplement net, signe d'un exode rural important.

2) Part des jeunes dans les départs :
\[ \frac{210}{300}\times 100 = 70\,\%. \]
Les 15-35 ans représentent \textbf{70 \%} des départs : la migration est fortement \cle{sélective} par l'âge. Le village perd sa main-d'œuvre agricole et ses futurs cadres, ce qui explique le vieillissement de la population restante et le risque d'abandon des terres.

3) Montant investi dans l'agriculture :
\[ 45 \times \frac{15}{100} = 6,75 \text{ millions de FCFA}. \]
\textbf{Conclusion.} La migration joue un double rôle : \emph{négatif} (perte de 25 \% de la population et de 70 \% de jeunes parmi les partants, donc baisse de la main-d'œuvre) mais aussi \emph{positif} (45 millions de FCFA injectés dans l'économie du village, dont 6,75 millions pour moderniser l'agriculture et 27 millions pour la consommation et la scolarité, ce qui améliore le niveau de vie des familles restantes). Un bon raisonnement de géographe tient ces deux vérités ensemble.
\end{exoresolu}

\begin{aretenir}[L'essentiel de la section]
\begin{itemize}
\item L'\cle{exode rural} (départ massif et durable des ruraux vers les villes) est \textbf{sélectif} : il emporte surtout les jeunes actifs de 15 à 35 ans.
\item Conséquences dans les régions de départ : \textbf{dépeuplement} et \textbf{vieillissement}, baisse de la \textbf{main-d'œuvre agricole} et \textbf{abandon de terres}, \textbf{désorganisation des familles} (femmes chefs de ménage, enfants chez les grands-parents), \textbf{perte des élites} (exode des élites rurales) et affaiblissement du développement local.
\item Les régions du \textbf{Nord}, de l'\textbf{Est} et de l'\textbf{Ouest} camerounais sont particulièrement touchées.
\item Contrepoints positifs : allègement de la \textbf{pression sur les terres}, réduction du chômage rural, \textbf{transferts de fonds} (plus de 200 milliards de FCFA par an venant de la diaspora).
\end{itemize}
\end{aretenir}

\begin{exercice}[Pour s'entraîner]
Dans un village de la région de l'Ouest, une enquête montre que 55 \% des ménages ont au moins un membre parti en ville, et que les transferts reçus représentent en moyenne 40 \% du revenu annuel de ces ménages.
1) Expliquez en cinq lignes pourquoi les jeunes de ce village partent vers Douala.
2) Dégagez deux effets négatifs et deux effets positifs de ces départs pour le village.
3) Rédigez une conclusion de quatre lignes montrant que la migration est un phénomène « ambivalent ».
\end{exercice}

\begin{corrige}[Exercice]
1) Les jeunes partent parce que les terres familiales sont morcelées et insuffisantes, que l'emploi rural est rare et mal payé, et que la ville (Douala) fait miroiter des emplois, des études et un mode de vie moderne : ce sont les facteurs de répulsion (rurale) et d'attraction (urbaine).
2) Effets négatifs : baisse de la main-d'œuvre agricole et vieillissement du village ; désorganisation des familles (femmes seules cheffes de ménage). Effets positifs : transferts d'argent qui font vivre 55 \% des ménages ; allègement de la pression sur les terres et baisse du chômage rural.
3) Conclusion attendue : la migration vide le village de ses forces vives (effet négatif) mais le nourrit par les transferts de fonds (effet positif) ; elle est donc ambivalente, et c'est la gestion de ces transferts qui décidera si le village s'appauvrit ou se transforme.
\end{corrige}

Ainsi, les régions de départ paient un lourd tribut à la migration, tout en recevant en retour des ressources précieuses. Mais qu'advient-il de ces migrants une fois arrivés ? C'est l'objet de la section suivante : les incidences des migrations dans les zones d'accueil.

\coursec{Les incidences des migrations dans les zones d'accueil}

\courssub{Transition : de l'espace de départ à l'espace d'accueil}

Dans la section précédente, nous avons analysé comment le départ des migrants — souvent les plus jeunes, les plus dynamiques, les plus qualifiés — vide les campagnes de leurs forces vives, y ralentit l'agriculture, y fragilise la cohésion sociale et y creuse la dépendance aux transferts d'argent. Ce « coût du départ » a une contrepartie immédiate : l'arrivée massive de ces mêmes populations dans les villes et les zones d'accueil. Si le migrant quitte un problème, il en apporte souvent de nouveaux là où il s'installe. La ville africaine, et camerounaise en particulier, absorbe chaque année des flux qu'elle n'a pas la capacité d'intégrer harmonieusement. C'est cette face cachée, mais bien visible dans le paysage urbain, que nous allons décortiquer : les incidences, positives comme négatives, des migrations dans les zones d'accueil.

\courssub{Croissance urbaine rapide et urbanisation anarchique : l'explosion des bidonvilles}

\begin{definition}[Bidonville / Quartier précaire]
Un \cle{bidonville} (ou quartier spontané, quartier d'habitat précaire) est un espace urbanisé de fait, généralement en marge de la ville légale, caractérisé par : une occupation du sol sans titre foncier ni permis de construire ; une densité de population très élevée ; l'absence quasi totale d'équipements collectifs (voirie, eau potable, électricité, assainissement, évacuation des ordures) ; et un habitat de fortune en matériaux de récupération (tôles, planches, plastiques, banco non stabilisé).
\end{definition}

L'afflux migratoire, combiné à la croissance naturelle, fait doubler la population des grandes métropoles camerounaises tous les 15 à 20 ans. Douala et Yaoundé concentrent à elles deux plus de 40\% de la population urbaine du pays. Les municipalités, dépassées par l'ampleur de la demande, ne peuvent pas suivre : les lotissements officiels sont trop lents à produire, trop chers pour les nouveaux arrivants. Ces derniers s'installent donc « où il y a de la place » : zones marécageuses (Nylon, Makepe, Bépanda à Douala), flancs de collines instables (Messa, Etoudi, Nkolbisson à Yaoundé), emprises ferroviaires ou routières, berges de fleuves.

\begin{exemplebox}[Le quartier Nylon à Douala : l'archétype du bidonville de migration]
Né dans les années 1970 sur d'anciennes zones marécageuses du Wouri, Nylon (acronyme de \emph{New Layout}) a absorbé des vagues successives : migrants du Nord-Ouest et de l'Ouest (années 80), déplacés des crises économiques (années 90), réfugiés centrafricains et nigérians (années 2010), et aujourd'hui déplacés internes de la crise anglophone.
\begin{itemize}
    \item \textbf{Population estimée :} plus de \num{200 000} habitants (certaines ONG parlent de 300 000) sur moins de 5 km², soit une densité supérieure à \num{40 000} hab/km².
    \item \textbf{Habitat :} alignements de cases en tôle/banco sur des parcelles de 3 m × 5 m, mitoyennes, sans cour, sans fenêtre sur l'extérieur.
    \item \textbf{Services :} 1 forage pour 5 000 habitants, pas de réseau d'égouts (fosses septiques à ciel ouvert ou « latrines volantes »), électricité pirate sur lignes moyenne tension, voirie inexistante (rues de 1 m de large, impraticables en saison des pluies).
    \item \textbf{Risques :} inondations récurrentes (crue du Wouri + ruissellement), incendies dévastateurs (propagation instantanée), épidémies (choléra, typhoïde, paludisme).
\end{itemize}
Nylon n'est pas une exception : c'est la norme pour l'accueil des migrants pauvres dans les métropoles camerounaises.
\end{exemplebox}

\begin{popfigure}
\begin{tikzpicture}[scale=0.85]
% Fond : zone marécageuse
\fill[popblue!10] (0,0) rectangle (14,9);
% Rivière / Wouri
\fill[popblue!40] (0,2) .. controls (3,1) and (6,3) .. (14,4) -- (14,5) .. controls (6,4) and (3,2) .. (0,3) -- cycle;
\draw[popblue, thick] (0,2.5) .. controls (3,1.5) and (6,3.5) .. (14,4.5);
% Quartier spontané : parcelles denses
\foreach \x in {0.5,1.2,...,13.5} {
    \foreach \y in {0.8,1.5,...,8.5} {
        \pgfmathsetmacro{\rx}{rand*0.2}
        \pgfmathsetmacro{\ry}{rand*0.2}
        \fill[poporange!70!popdark] (\x+\rx,\y+\ry) rectangle ++(0.55,0.55);
        \draw[popdark, very thin] (\x+\rx,\y+\ry) rectangle ++(0.55,0.55);
    }
}
% Ruelles étroites (espace vide entre blocs)
% Quelques "rues" principales un peu plus larges
\draw[popmuted, line width=3pt] (0.2,5) -- (13.8,5.2);
\draw[popmuted, line width=2pt] (7,0.2) -- (7.2,9);
% Points d'eau (forages)
\foreach \pos in {(2,2), (5,6), (9,3), (11,7)} {
    \fill[popblue] \pos circle (3pt);
    \node[above, font=\tiny, text=popdark] at \pos {Forage};
}
% Latrines / fosses
\foreach \pos in {(3,8), (8,1), (12,5)} {
    \fill[poporange!50!black] \pos circle (4pt);
    \node[below, font=\tiny, text=popdark] at \pos {Fosse};
}
% Décharges
\fill[popgreen!40!black] (13,0.5) rectangle (13.8,1.5);
\node[rotate=90, font=\tiny, text=popdark] at (13.4,1) {Décharge};
% Légende intégrée
\node[align=left, font=\small, text=popdark, fill=white, fill opacity=0.8, draw=popline, rounded corners, inner sep=4pt] at (1.5,8.5) {
    \textbf{Légende simplifiée} \\
    \textcolor{poporange!70!popdark}{\rule{3mm}{3mm}} Cases en tôle/banco \\
    \textcolor{popblue}{\rule{3mm}{3mm}} Point d'eau (forage) \\
    \textcolor{poporange!50!black}{\rule{3mm}{3mm}} Fosse septique / latrine \\
    \textcolor{popgreen!40!black}{\rule{3mm}{3mm}} Dépôt d'ordures \\
    \textcolor{popmuted}{\rule{5mm}{2mm}} Ruelle principale
};
% Flèche Nord
\draw[popdark, ->, thick] (13,8) -- (13,8.8) node[above, font=\tiny] {N};
% Échelle
\draw[popdark, thick] (0.5,0.2) -- (2.5,0.2) node[midway, below, font=\tiny] {~200 m};
\end{tikzpicture}
\legende{Croquis schématique d'un quartier spontané type (inspiré de Nylon, Douala) : habitat dense en damier irrégulier, absence de voirie structurée, points d'eau rares, fosses septiques et décharges en bordure de rivière. L'orientation Nord et l'échelle indicative permettent de situer l'étendue du phénomène.}
\end{popfigure}

\courssub{Surpopulation et pression sur les services urbains de base}

L'arrivée massive de migrants transforme la demande en services en une urgence permanente. La ville « formelle » (lotissements, centre-ville, quartiers résidentiels) voit ses infrastructures saturées ; la ville « informelle » (bidonvilles) en est dépourvue.

\begin{propriete}[Effet de seuil sur les équipements urbains]
Lorsqu'un équipement (école, centre de santé, point d'eau, km de voirie) est dimensionné pour une population $P_0$ et que la population réelle devient $P = P_0 + \Delta P_{mig}$ (avec $\Delta P_{mig}$ la composante migratoire), le taux de couverture chute brutalement. Si le temps de réaction administratif (planification, financement, réalisation) est $T_{admin} \approx 5\text{--}10$ ans, et que le doublement migratoire se fait en $T_{mig} \approx 3\text{--}5$ ans, alors $T_{mig} < T_{admin}$ : la ville est structurellement en retard.
\end{propriete}

\begin{itemize}
    \item \textbf{Logement :} déficit estimé à plus de \num{1 000 000} de logements au Cameroun (MINHDU, 2023). Le migrant pauvre loue une « chambre » de 9 m² sans sanitaire privé, à \num{15 000}--\num{25 000} FCFA/mois, souvent partagée à 3 ou 4.
    \item \textbf{Eau et assainissement :} à Yaoundé, seulement 35\% des ménages sont raccordés au réseau CAMWATER ; les autres dépendent des forages privés, des vendeurs d'eau (prix multiplié par 10 à 20) ou de l'eau de surface polluée. L'assainissement collectif couvre moins de 10\% de la ville ; le reste est fosses septiques non vidangées ou déjection à l'air libre.
    \item \textbf{Éducation :} surcharge des classes (80 à 120 élèves/classe dans le public périphérique), écoles sous paillottes, turn-over enseignant élevé. Le migrant non francophone (anglophone, centrafricain, nigérian) ajoute une barrière linguistique.
    \item \textbf{Santé :} 1 lit pour 1 000 habitants (norme OMS : 5), pénurie de personnel, ruptures de stock. Les migrants, non assurés, retardent les soins → complications, surmortalité évitable.
\end{itemize}

\begin{exoresolu}[Analyse d'un tableau de pression urbaine : Yaoundé 2010 vs 2023]
\textbf{Énoncé :} À partir du tableau ci-dessous, calculez le taux de croissance annuel moyen (TCAM) de la population et de chaque équipement. Concluez sur l'écart.

\begin{center}
\begin{tabularx}{\linewidth}{|l|X|X|X|}\hline
\textbf{Indicateur} & \textbf{2010} & \textbf{2023} & \textbf{TCAM (\%/an)} \\ \hline
Population (hab) & 1 800 000 & 4 200 000 & ? \\ \hline
Lits d'hôpital & 1 800 & 2 500 & ? \\ \hline
Salles de classe publiques & 3 200 & 5 100 & ? \\ \hline
Km de voirie bitumée & 450 & 620 & ? \\ \hline
\end{tabularx}
\end{center}

\textbf{Solution détaillée :}
Formule : $TCAM = \left( \left(\frac{V_{fin}}{V_{deb}}\right)^{\frac{1}{n}} - 1 \right) \times 100$ avec $n = 13$ ans.
\begin{align*}
\text{Population : } & \left( \frac{4\,200\,000}{1\,800\,000} \right)^{1/13} = 2,333^{0,0769} \approx 1,067 \Rightarrow \textbf{6,7 \%/an} \\
\text{Lits : } & \left( \frac{2\,500}{1\,800} \right)^{1/13} = 1,389^{0,0769} \approx 1,026 \Rightarrow \textbf{2,6 \%/an} \\
\text{Salles : } & \left( \frac{5\,100}{3\,200} \right)^{1/13} = 1,594^{0,0769} \approx 1,037 \Rightarrow \textbf{3,7 \%/an} \\
\text{Voirie : } & \left( \frac{620}{450} \right)^{1/13} = 1,378^{0,0769} \approx 1,025 \Rightarrow \textbf{2,5 \%/an}
\end{align*}
\textbf{Conclusion :} La population croît près de \textbf{3 fois plus vite} que les équipements de santé, d'éducation et de voirie. L'écart se creuse chaque année : en 2023, on a 0,6 lit/1000 hab (vs 1,0 en 2010), 1,2 salle/1000 hab (vs 1,8), 0,15 km/1000 hab (vs 0,25). La migration (estimée à 60\% de la croissance) est le moteur principal de ce déséquilibre.
\end{exoresolu}

\courssub{Chômage, sous-emploi et explosion du secteur informel}

Le marché du travail formel (public + privé moderne) crée environ \num{30 000} à \num{40 000} emplois/an au Cameroun, pour une entrée annuelle sur le marché de \num{300 000} à \num{400 000} jeunes (dont une forte proportion de migrants ruraux et de diplômés chômeurs). L'écart est comblé par le secteur informel, qui emploie plus de 90\% de la population active urbaine.

\begin{aretenir}[Secteur informel urbain : caractéristiques et lien avec la migration]
\begin{itemize}
    \item \textbf{Accessibilité :} pas de diplôme requis, capital de départ faible (ex. moto-taxi : \num{300 000}--\num{500 000} FCFA en location-vente), entrée immédiate.
    \item \textbf{Flexibilité :} horaires libres, mobilité spatiale (suivi des chantiers, des marchés, des flux clients).
    \item \textbf{Précarité :} pas de contrat, pas de protection sociale (CNPS), pas de retraite, revenus journaliers variables (\num{2 000}--\num{5 000} FCFA/jour pour un moto-taximan après charges), exposition aux risques (accidents, racket, intempéries).
    \item \textbf{Profil migrant :} le migrant récent (moins de 5 ans en ville) est surreprésenté dans l'informel de survie (pousse-pousse, portefaix, petit commerce de rue, lavage de voitures). L'ancien migrant ou le natif accède davantage à l'informel « supérieur » (artisanat, commerce de gros, services).
\end{itemize}
\end{aretenir}

\begin{methode}[Analyser l'insertion professionnelle d'un migrant urbain : grille d'entretien]
Pour une enquête de terrain (TP géographie), utilisez cette grille :
\begin{enumerate}
    \item \textbf{Profil :} âge, sexe, origine (département/région), année d'arrivée, motif déclaré.
    \item \textbf{Parcours résidentiel :} quartiers successifs, type d'habitat, loyer.
    \item \textbf{Parcours professionnel :} 1er emploi en ville, emplois actuels (principal + secondaires), statut (indépendant, apprenti, salarié informel), revenus estimés, cotisation CNPS (oui/non).
    \item \textbf{Réseaux :} comment a-t-il trouvé le logement ? le travail ? (parenté, connaissance villageoise, association de développement, hasard).
    \item \textbf{Perception :} « La ville vous a-t-elle déçu ? » « Que manque-t-il pour réussir ? » « Voulez-vous retourner au village ? »
\end{enumerate}
Croisez les réponses pour tester l'hypothèse : « Le migrant récent est cantonné aux segments les plus précaires de l'informel. »
\end{methode}

\courssub{Problèmes sociaux : insécurité, délinquance, tensions identitaires}

La concentration de jeunes hommes sans emploi stable, sans encadrement familial, dans des espaces surpeuplés et mal éclairés, crée un terreau favorable à la déviance.

\begin{itemize}
    \item \textbf{Insécurité et délinquance :} vols à l'arraché (téléphones, sacs), braquages (moto-taxis), « coupeurs de route » urbains (bandes de \emph{microbes} à Abidjan, \emph{kulunas} à Kinshasa, \emph{bagarreurs} à Douala/Yaoundé). La police est sous-effectivée (1 policier pour 800--1 000 hab vs norme 1/400).
    \item \textbf{Prostitution et exploitation sexuelle :} migration féminine de plus en plus autonome (filles « placement », étudiantes en rupture, réfugiées). Prostitution de survie (\num{1 000}--\num{5 000} FCFA/passe), réseaux de traite (vers le Moyen-Orient, l'Europe, ou interne vers les zones minières/pétrolières).
    \item \textbf{Tensions interethniques / intercommunautaires :} la ville recompose les identités. Quartiers « ethniques » (Bamileke à Bépanda, Bassa à New-Bell, Haoussa/Foulbé à Haussaré, Centrafricains à PK12 Yaoundé). Conflits pour le contrôle des marchés (ex. marché Central Yaoundé : Bamiléké vs Haoussa), des parkings, des taxes informelles. Instrumentalisation politique (élections, crises anglophones).
\end{itemize}

\begin{attention}[Piège fréquent : le migrant bouc émissaire]
\textbf{Erreur :} « C'est la faute aux migrants si la ville est sale, insécure, embouteillée. »
\textbf{Réalité :} La croissance \textbf{naturelle} de la population urbaine (taux de fécondité encore élevé en ville : 4,2 enfants/femme à Yaoundé, 3,8 à Douala) contribue autant, voire plus, à la pression démographique que la migration nette. De plus, l'insécurité et l'insalubrité sont d'abord des \textbf{défaillances de gouvernance urbaine} (absence de planification, corruption foncière, sous-investissement chronique, justice lente). Le migrant est souvent la \emph{première victime} de ces dysfonctionnements (logement insalubre, racket policier, expulsion sans relogement). Ne confondez pas \emph{corrélation} (arrivée migrants + dégradation) et \emph{causalité unique}.
\end{attention}

\courssub{Incidences positives : main-d'œuvre, dynamisme, brassage}

Réduire le migrant à un « problème » est une erreur analytique majeure. L'histoire urbaine du Cameroun s'est écrite par et pour les migrants.

\begin{propriete}[Gains économiques et sociaux de l'immigration urbaine]
\begin{itemize}
    \item \textbf{Main-d'œuvre abondante et bon marché :} BTP (maçons, ferrailleurs, peintres — majoritairement Ouest/Nord-Ouest), industrie (usines de Bonabéri, Bassa — main-d'œuvre nationale et sous-régionale), services (domestiques, gardiens, commerce). Cette compétitivité-coût attire les investisseurs.
    \item \textbf{Dynamisme entrepreneurial :} Les migrants sont surreprésentés dans la création d'entreprises (formelles et informelles). Ex. : le réseau commercial bamiléké structure l'approvisionnement national en biens de consommation ; les Haoussa/Foulbé dominent le commerce du bétail et de la viande ; les Nigérians animent l'import-export informel (friperie, électronique, médicaments).
    \item \textbf{Brassage culturel et innovation :} métissage linguistique (Camfranglais, pidgin), culinaire (ndolé, eru, kilichi, suya, poisson braisé), musical (makossa, bikutsi, coupe-décalé, afrobeats nés du mélange urbain), religieux (églises de réveil, mosquées soufies). La ville est le creuset de l'identité nationale camerounaise.
    \item \textbf{Réseaux de solidarité :} tontines, associations de développement (AD), mutuelles villageoises financent l'école, la santé, l'eau, l'électricité au village \emph{et} en ville (rotations d'épargne pour loyer, maladie, mariage). Le migrant est un acteur de développement \emph{par le bas}.
\end{itemize}
\end{propriete}

\courssub{Cas des migrations internationales d'accueil : réfugiés centrafricains et nigérians}

Le Cameroun, par sa stabilité relative, est un pays d'asile majeur en Afrique centrale (loi n°2005/006 du 27 juillet 2005 sur le statut de réfugié).

\begin{center}
\begin{tabularx}{\linewidth}{|l|X|X|X|}\hline
\textbf{Flux} & \textbf{Origine / Cause} & \textbf{Zones d'accueil principales} & \textbf{Effectifs estimés (HCR/ONA, 2023--24)} \\ \hline
Réfugiés centrafricains & Guerres civiles successives (2013--présent), exactions Seleka/anti-Balaka & Région de l'Est (Gado-Badzéré, Lolo, Mbilé, Garoua-Boulaï, Batouri, Bertoua) & \num{~350 000} (dont \num{~120 000} en sites, le reste en milieu urbain/hôte) \\ \hline
Réfugiés nigérians & Insurrection Boko Haram / ISWAP (depuis 2014) & Région de l'Extrême-Nord (Minawao, Mora, Kousseri, Maroua, Kolofata) & \num{~120 000} (dont \num{~70 000} au camp de Minawao) \\ \hline
Déplacés internes (Cameroun) & Crise anglophone (N-O / S-O) + Boko Haram (Extrême-Nord) & Villes du Littoral, Centre, Ouest, Adamaoua, Nord & \num{~1 000 000} (total IDP) \\ \hline
\end{tabularx}
\end{center}

\begin{experience}[Étude de cas : Le camp de Minawao (Extrême-Nord) — de l'urgence à la ville-réfugiés]
\textbf{Contexte :} Ouvert en 2013, capacité initiale 15 000 places. Population 2024 : \num{~70 000} réfugiés nigérians (majorité Kanouri, Shuwa Arab, Foulbé).
\textbf{Observations (données HCR/ACNUR) :}
\begin{itemize}
    \item \textbf{Organisation spatiale :} quadrillage en « blocs » (A à Z), rues larges de 10 m, parcelles familiales 12 m × 15 m. Écoles (17 primaires, 2 secondaires), centre de santé (chirurgie, maternité), marché central, mosquées, église.
    \item \textbf{Économie du camp :} marché hebdomadaire attractif pour les populations hôtes (bétail, tissus, céréales, artisanat). Monnaie : FCFA + Naira. Activités génératrices de revenus (AGR) soutenues par ONG : couture, boulangerie, transformation agroalimentaire, mécanique moto.
    \item \textbf{Relations hôtes-réfugiés :} tensions foncières (camp sur terres coutumières), concurrence eau/bois/pâturage, mais aussi complémentarité économique (main-d'œuvre agricole réfugiée pour paysans hôtes, achats locaux).
    \item \textbf{Problèmes persistants :} dépendance aide alimentaire (ratio couverture 70\%), scolarisation filles < 50\% (mariages précoces), radicalisation résiduelle (recrutement Boko Haram), absence de perspectives de retour (insécurité persistante au Nord-Est Nigeria).
\end{itemize}
\textbf{Interprétation :} Un camp de réfugiés qui dure plus de 10 ans devient une \textbf{ville de fait}, avec ses services, son économie, sa gouvernance (comités de blocs, leaders traditionnels). La gestion passe de l'humanitaire d'urgence au développement local intégré (approche « Nexus humanitaire-développement »).
\end{experience}

\begin{savaistu}[Le Cameroun, « pays tampon » et « pays pivot » migratoire]
Le Cameroun est simultanément :
\begin{itemize}
    \item \textbf{Pays d'émigration :} vers la France, les USA, l'Allemagne, le Canada, le Qatar, le Koweït (élites, étudiants, main-d'œuvre qualifiée et non qualifiée) → transferts \num{> 1 000} milliards FCFA/an (1er source de devises hors pétrole).
    \item \textbf{Pays de transit :} route migratoire vers la Libye/Algérie/Méditerranée (Tchadiens, Soudanais, Centrafricains, Nigériens) → trafic de migrants, risques traite.
    \item \textbf{Pays d'immigration / asile :} comme vu ci-dessus (Centrafricains, Nigérians, mais aussi Tchadiens, Congolais (RDC), Rwandais, Maliens).
\end{itemize}
Cette triple position impose une politique migratoire cohérente (loi 2005, stratégie nationale 2020--2030) encore perfectible : régularisation, intégration socio-économique des réfugiés (droit au travail, à la terre, à la banque), lutte contre les réseaux criminels, protection des migrants camerounais à l'étranger.
\end{savaistu}

\courssub{Évolution démographique des deux métropoles : visualisation quantitative}

\begin{popfigure}
\begin{tikzpicture}
\begin{axis}[
    width=\linewidth,
    height=9cm,
    xlabel={Année},
    ylabel={Population (millions d'habitants)},
    title={Évolution de la population de Douala et Yaoundé (1976--2024)},
    legend pos=north west,
    grid=major,
    grid style={popline},
    xmin=1975, xmax=2025,
    ymin=0, ymax=5.5,
    xtick={1976,1987,2005,2010,2015,2020,2024},
    xticklabel style={rotate=45, anchor=east, font=\small},
    ytick={0,1,2,3,4,5},
    yticklabel style={font=\small},
    popcurve/.style={thick, mark=*, mark size=3},
]
% Données sources : Recensements 1976, 1987, 2005 (BUCREP) + Projections INS/World Urbanization Prospects 2024
% Douala
\addplot[popblue, popcurve] coordinates {
    (1976,0.45) (1987,0.81) (2005,1.91) (2010,2.44) (2015,3.05) (2020,3.65) (2024,4.15)
};
% Yaoundé
\addplot[poporange, popcurve] coordinates {
    (1976,0.38) (1987,0.68) (2005,1.67) (2010,2.15) (2015,2.77) (2020,3.42) (2024,4.00)
};
\legend{Douala, Yaoundé}
% Annotation croissance
\node[align=left, font=\small, fill=popcream, draw=popline, rounded corners, inner sep=4pt] at (rel axis cs:0.02,0.95) {
    \textbf{TCAM moyen 1976--2024 :} \\
    Douala : \textbf{4,9 \%/an} \\
    Yaoundé : \textbf{5,1 \%/an} \\
    \textit{(Croissance naturelle ~2,5\% + Migration nette ~2,5\%)}
};
\end{axis}
\end{tikzpicture}
\legende{Évolution comparée des populations de Douala et Yaoundé depuis le premier recensement post-indépendance. La courbe exponentielle témoigne de l'incapacité des infrastructures à suivre le rythme. Sources : BUCREP (recensements 1976, 1987, 2005), INS (projections), ONU-Habitat (World Urbanization Prospects 2024).}
\end{popfigure}

\begin{methode}[Construire un croquis de paysage urbain précaire (type bidonville) — Exigence Bac]
Pour réussir un croquis géographique noté au Probatoire/Bac, suivez cette procédure rigoureuse :
\begin{enumerate}
    \item \textbf{Titre informatif :} Localisation + Thème + Date (ex. : « Croquis du quartier Nylon, Douala — Habitat précaire et risques environnementaux, 2024 »).
    \item \textbf{Cadre et orientation :} Rectangle proportionné, flèche Nord (haut de la feuille par convention), échelle indicative (ex. 1 cm = 200 m).
    \item \textbf{Fond de carte (structure) :} Hydrographie (fleuve, marécage), axes de communication principaux (route goudronnée, voie ferrée), limite administrative si pertinente.
    \item \textbf{Occupation du sol (zonage par couleurs/hachures) :}
    \begin{itemize}
        \item Habitat dense (hachures fines serrées, couleur \textcolor{poporange!70!popdark}{orange/brique}).
        \item Zones non bâties / risques (marécage \textcolor{popblue!30}{bleu clair}, décharge \textcolor{popgreen!40!black}{vert foncé}, carrière \textcolor{popmuted}{gris}).
        \item Équipements rares (école \textcolor{popblue}{bleu}, forage \textcolor{popblue}{point bleu}, centre de santé \textcolor{popred}{croix rouge}, marché \textcolor{popgold}{jaune}).
    \end{itemize}
    \item \textbf{Réseaux et flux :} Ruelle principale (trait discontinu \textcolor{popmuted}{gris}), flux piétons (flèches fines), évacuation eaux usées (flèches \textcolor{popblue}{bleues} vers fleuve).
    \item \textbf{Légende organisée :} Regroupée par thèmes (Habitat, Équipements, Risques, Circulation), symboles identiques au croquis, hiérarchisés.
    \item \textbf{Commentaire géographique (obligatoire) :} 5--7 lignes analysant : localisation (marge, risque), morphologie (densité, anarchie), fonctionnement (économie de survie, insalubrité), enjeux (intégration, réhabilitation, résilience).
\end{enumerate}
\textbf{Astuce :} Utilisez des crayons de couleur (pas de feutre qui bave), règle pour le cadre, main levée pour l'habitat (aspect « organique »). La lisibilité prime sur le détail exhaustif.
\end{methode}

\begin{exercice}[Application : Diagnostic territorial d'un quartier d'accueil]
Choisissez un quartier de votre ville (ou un quartier type connu : Nylon, Messa, Bépanda, Mokolo, Akwa, Bonabéri).
\begin{enumerate}
    \item Réalisez un croquis simplifié (format A4) respectant la méthode ci-dessus.
    \item Rédigez un diagnostic en 15 lignes structurées : \textbf{Atouts} (proximité emploi, marchés, réseaux solidaires, transport) vs \textbf{Contraintes} (insalubrité, insécurité, précarité foncière, déficit services, risques climatiques).
    \item Proposez 3 mesures concrètes d'amélioration (court, moyen, long terme) en distinguant ce qui relève de la \textbf{municipalité}, de l'\textbf{État central} et de la \textbf{population elle-même} (participative).
\end{enumerate}
\end{exercice}

\begin{corrige}[Exercice n°1 : Diagnostic territorial — Éléments attendus]
\textbf{Croquis :} Vérifier présence : titre, nord, échelle, légende structurée, distinction habitat/équipements/risques, commentaire.
\textbf{Diagnostic type (ex. Nylon) :}
\begin{itemize}
    \item \textit{Atouts :} Proximité zone industrielle (Bonabéri) et port (emplois journaliers), marché de gros (revenus commerçants), réseaux ethniques/associatifs forts (entraide, tontines), desserte relative (taxis, motos).
    \item \textit{Contraintes :} Inondations annuelles (Wouri + ruissellement), densité extrême (>40 000 hab/km²), absence voirie/assainissement (maladies hydriques), insécurité (bandes, racket), précarité foncière (menace expulsion), pollution (air, eau, sol).
\end{itemize}
\textbf{Mesures :}
\begin{itemize}
    \item \textit{Court terme (Mairie/État) :} Curage caniveaux avant saison pluies, points d'eau d'urgence, éclairage public solaire, poste de police de proximité, ramassage ordures régulier.
    \item \textit{Moyen terme (Mairie/MINHDU) :} Réhabilitation voirie (pavés/dalles béton), lotissement partiel (permis d'occuper), construction école/CSPS, régularisation foncière progressive (certificats d'occupation).
    \item \textit{Long terme (État/Partenaires) :} Plan d'urbanisme directeur (PDU) intégrant le quartier, relogement zones à risque non constructibles (digues, parcs), raccordement réseaux eau/électricité formels, programmes emploi jeunes (chantiers école, formation qualifiante).
    \item \textit{Participatif :} Comités de quartier (hygiène, sécurité, médiation), cotisations pour micro-réalisations (pontons, dallage ruelles), veille environnementale (signalement décharges, bouchons caniveaux).
\end{itemize}
\end{corrige}

\begin{aretenir}[Synthèse : Les incidences des migrations dans les zones d'accueil — Bilan contrasté]
\begin{center}
\begin{tabularx}{\linewidth}{|>{\bfseries}X|X|X|}\hline
\textbf{Dimension} & \textbf{Coûts / Dysfonctionnements (Négatifs)} & \textbf{Bénéfices / Potentiels (Positifs)} \\ \hline
\textbf{Démographique / Urbain} & Urbanisation anarchique, bidonvilles, surdensité, pression foncière & Croissance du marché urbain, densification économique, vitalité démographique \\ \hline
\textbf{Équipements / Services} & Saturation (école, santé, eau, voirie, transport), dégradation qualité & Demande justifiant investissements publics, innovation services (mobile money, e-santé) \\ \hline
\textbf{Économie / Emploi} & Chômage de masse, sous-emploi, informel de survie, bas salaires & Main-d'œuvre disponible/bon marché, esprit d'entreprise migrant, niches commerciales, transferts \\ \hline
\textbf{Social / Sécurité} & Insécurité, délinquance, prostitution, tensions identitaires, exclusion & Brassage culturel, solidarités communautaires, recomposition identitaire nationale \\ \hline
\textbf{Environnement} & Pollution, déforestation peri-urbaine, inondations, gestion déchets défaillante & Agriculture urbaine (marais), récupération/recyclage informel (économie circulaire) \\ \hline
\end{tabularx}
\end{center}
\textbf{Conclusion pédagogique :} Le migrant n'est ni un « fardeau » ni un « sauveur » : c'est un \textbf{acteur stratégique} dont l'intégration réussie dépend de la \textbf{capacité de gouvernance} de la ville (planification, investissement, inclusion, droit à la ville). Le défi camerounais n'est pas d'arrêter la migration (impossible, indésirable), mais de \textbf{la gérer} : transformer l'informel en formel, le précaire en durable, l'exclusion en citoyenneté.
\end{aretenir}

\coursec{Les migrations internationales : gains et pertes pour le Cameroun}

Dans la section précédente, nous avons étudié les incidences des migrations dans les zones d'accueil à l'intérieur du pays, notamment la pression exercée sur les villes comme Douala et Yaoundé. Mais les mouvements migratoires ne s'arrêtent pas aux frontières du Cameroun. Chaque année, des milliers de Camerounais quittent leur pays pour l'étranger, tandis que des ressortissants des pays voisins viennent s'installer chez nous. Ces \cle{migrations internationales} ont des conséquences ambivalentes : elles appauvrissent le pays en compétences, mais elles l'enrichissent financièrement. C'est ce double visage, fait de \cle{gains} et de \cle{pertes}, que nous allons analyser, avant de dégager un bilan chiffré.

\courssub{La fuite des cerveaux : une perte majeure de compétences}

\textbf{Intuition.} Imaginons un hôpital de district dans l'Extrême-Nord où un seul médecin dessert plusieurs dizaines de milliers d'habitants. Si ce médecin part exercer en France ou au Canada, c'est toute une population qui se retrouve sans soins. Multiplions ce cas par des milliers : c'est le phénomène de la \cle{fuite des cerveaux}.

\begin{definition}[Fuite des cerveaux]
La \cle{fuite des cerveaux} (en anglais \emph{brain drain}) est l'\cle{émigration} des personnes \cle{qualifiées} (médecins, infirmiers, ingénieurs, enseignants, chercheurs, informaticiens) d'un pays en développement vers des pays plus développés, attirées par de meilleurs salaires, de meilleures conditions de travail et de meilleures perspectives de carrière.
\end{definition}

\textbf{Mécanisme.} Le départ des cadres camerounais s'explique par un jeu de \cle{facteurs de répulsion} et de \cle{facteurs d'attraction} :

\begin{itemize}
\item \textbf{Facteurs de répulsion (push factors)} : salaires faibles et parfois impayés, manque d'équipements dans les hôpitaux et les laboratoires, chômage des diplômés, faibles perspectives de promotion.
\item \textbf{Facteurs d'attraction (pull factors)} : salaires élevés en Europe et en Amérique du Nord, conditions de travail modernes, sécurité sociale, possibilités de formation continue.
\end{itemize}

Les secteurs les plus touchés sont la \cle{santé} (médecins et infirmiers partis vers le Canada, la France, l'Allemagne), l'\cle{éducation} (enseignants et universitaires) et l'\cle{ingénierie} (informaticiens, ingénieurs des télécommunications). Le Cameroun compte environ \textbf{1 médecin pour \num{10000} habitants}, très en dessous de la norme recommandée par l'Organisation mondiale de la santé (au moins 1 médecin pour \num{1000} habitants), alors que des milliers de médecins formés au Cameroun exercent à l'étranger.

\textbf{Le coût de la formation perdue pour l'État camerounais.} Former un médecin coûte très cher à l'État : bourses, fonctionnement des facultés de médecine, hôpitaux universitaires. On estime le coût de formation d'un médecin généraliste à plusieurs dizaines de millions de FCFA, entièrement ou largement supporté par les deniers publics. Lorsque ce médecin part s'installer en Europe, c'est un \cle{investissement perdu} : l'État camerounais a payé la formation, mais c'est le pays d'accueil qui récolte les bénéfices du travail. On parle d'une véritable \emph{subvention indirecte} des pays pauvres aux pays riches. La même logique vaut pour les ingénieurs sortis des grandes écoles publiques et les enseignants formés dans les écoles normales supérieures.

\begin{attention}[Ne pas confondre]
Il faut bien distinguer la \cle{migration régulière} (le migrant possède un \cle{visa}, un passeport valide et respecte les lois du pays d'accueil) de la \cle{migration irrégulière} ou \cle{clandestine} (le migrant entre ou séjourne sans autorisation, souvent au péril de sa vie). La fuite des cerveaux concerne le plus souvent des migrations régulières, car les cadres qualifiés sont recrutés officiellement. Ne mélangez pas les deux notions dans vos rédactions.
\end{attention}

\courssub{Les routes de la migration clandestine et leurs drames}

Tous les départs ne se font pas avec un visa. Face à la fermeture des voies légales, de nombreux jeunes Camerounais tentent l'aventure clandestine vers l'Europe, au prix de risques considérables.

\textbf{Le parcours type.} Le candidat au départ quitte le Cameroun, traverse le \cle{Nigeria}, remonte vers le \cle{Niger} (ville d'Agadez, plaque tournante des migrations sahariennes), franchit le \cle{désert du Sahara} dans des conditions extrêmes (chaleur, soif, abandons par les passeurs), arrive en \cle{Libye} où il est souvent victime de détention, de racket, voire de traitements inhumains, avant de tenter la traversée de la \cle{Méditerranée} sur des embarcations de fortune vers l'Italie.

\begin{exemplebox}[Un drame chiffré]
Selon l'Organisation internationale pour les migrations (OIM), \textbf{plusieurs milliers de migrants meurent chaque année en Méditerranée} : plus de \num{3000} décès recensés en 2023, et près de \num{30000} depuis 2014, sans compter les disparus dans le Sahara, dont les corps ne sont jamais retrouvés. Parmi ces victimes figurent des \textbf{Camerounais}, partis chercher en Europe un avenir qu'ils ne trouvaient pas chez eux. Beaucoup sont aussi victimes du \cle{trafic d'êtres humains} : réseaux de passeurs, exploitation, travail forcé, prostitution.
\end{exemplebox}

\begin{popfigure}
\begin{center}
\begin{tikzpicture}[scale=1,font=\small]
% Fond mer
\fill[popblueL!60] (-6,3.2) rectangle (6,6);
\node[popblue!70] at (0,5.6) {\textbf{Mer Méditerranée}};
% Europe
\fill[popgreenL] (-6,6) rectangle (6,7.4);
\node[popdark] at (0,6.9) {\textbf{EUROPE (Italie, France, Allemagne...)}};
% Afrique du Nord
\fill[popgoldL] (-6,1.2) rectangle (6,3.2);
\node[popdark] at (3.5,2.6) {\textbf{Libye}};
\node[popdark] at (-3.5,2.6) {Algérie};
% Sahel
\fill[poporangeL!60] (-6,-1.2) rectangle (6,1.2);
\node[popdark] at (-1.5,0.4) {\textbf{Niger}};
\fill[popink] (-1.5,0.9) circle (0.07);
\node[below,popdark] at (-1.5,0.9) {\footnotesize Agadez};
% Nigeria et Cameroun
\fill[popgreenL!70] (-6,-3.2) rectangle (1,-1.2);
\node[popdark] at (-3.5,-2.2) {\textbf{Nigeria}};
\fill[poppinkL] (1,-3.2) rectangle (6,-1.2);
\node[popdark] at (3.5,-2.2) {\textbf{CAMEROUN}};
\fill[popink] (3.2,-2.6) circle (0.07);
\node[below,popdark] at (3.2,-2.6) {\footnotesize Yaoundé};
% Routes
\draw[->,line width=2.2pt,poporange] (3.2,-2.5) -- (-2.5,-1.9);
\draw[->,line width=2.2pt,poporange] (-2.5,-1.8) -- (-1.5,0.8);
\draw[->,line width=2.2pt,poporange] (-1.5,1.0) -- (2.5,2.3);
\draw[->,line width=2.6pt,poppink] (2.5,2.5) -- (0.5,5.4);
\node[rotate=62,poppink] at (2.0,4.1) {\footnotesize traversée mortelle};
% Nord
\draw[->,thick,popdark] (-5.4,4.6) -- (-5.4,5.6);
\node[above,popdark] at (-5.4,5.6) {\footnotesize N};
% Échelle indicative
\draw[thick,popdark] (-5.6,-3.5) -- (-3.6,-3.5);
\node[below,popdark] at (-4.6,-3.5) {\footnotesize $\approx$ \SI{1000}{km}};
\end{tikzpicture}
\end{center}
\legende{Schéma simplifié des routes migratoires clandestines du Cameroun vers l'Europe : 1. Cameroun $\rightarrow$ Nigeria ; 2. Nigeria $\rightarrow$ Agadez (Niger) ; 3. traversée du Sahara vers la Libye ; 4. traversée de la Méditerranée vers l'Italie. Croquis indicatif, sans précision cartographique.}
\end{popfigure}

\courssub{Les transferts de fonds de la diaspora : une ressource majeure}

Face à ces pertes, il existe un gain considérable : l'argent envoyé par les Camerounais de l'étranger.

\begin{definition}[Transferts de fonds]
Les \cle{transferts de fonds} (ou \emph{remittances}) sont les sommes d'argent envoyées par les migrants installés à l'étranger à leurs familles restées au pays, par les banques, les sociétés de transfert (Western Union, MoneyGram) ou les applications mobiles.
\end{definition}

\textbf{Utilisation des fonds.} Ces envois constituent le premier revenu de nombreuses familles camerounaises. Ils financent principalement :

\begin{enumerate}
\item la \cle{construction} de maisons (les fameuses « maisons de la diaspora » dans les villages et les quartiers urbains) ;
\item l'\cle{éducation} des enfants (frais de scolarité, fournitures, études universitaires) ;
\item la \cle{santé} (paiement des soins, médicaments) ;
\item le \cle{commerce} et les petites entreprises (ouverture de boutiques, achats-reventes, transport) ;
\item les cérémonies et la consommation courante.
\end{enumerate}

\begin{popfigure}
\begin{center}
\begin{tikzpicture}
\begin{axis}[
ybar, bar width=14pt,
width=0.9\linewidth, height=6.5cm,
xlabel={Année}, ylabel={Milliards de FCFA},
symbolic x coords={2010,2013,2016,2019,2021,2023},
xtick=data, x tick label style={rotate=30,anchor=east},
ymin=0, ymax=260,
nodes near coords, nodes near coords style={font=\footnotesize},
axis lines=left,
ymajorgrids, grid style={popline!40},
]
\addplot[fill=popblue!70,draw=popblue] coordinates {(2010,60) (2013,95) (2016,130) (2019,170) (2021,190) (2023,220)};
\end{axis}
\end{tikzpicture}
\end{center}
\legende{Évolution indicative des transferts de fonds de la diaspora camerounaise, 2010--2023 (ordres de grandeur en milliards de FCFA, d'après les données de la Banque mondiale et de la BEAC). La tendance est nettement croissante.}
\end{popfigure}

\begin{savaistu}[La diaspora, premier « investisseur » du Cameroun ?]
Les transferts de la diaspora représentent une \textbf{part importante du PIB camerounais} (plusieurs centaines de milliards de FCFA par an). Ils dépassent, certaines années, le montant des \cle{investissements directs étrangers} (IDE) et rivalisent avec l'aide publique au développement. Autrement dit, les Camerounais de l'étranger financent l'économie nationale plus sûrement que bien des investisseurs étrangers, et cet argent va directement aux familles, sans passer par les circuits administratifs.
\end{savaistu}

\courssub{Les autres apports de la diaspora}

La diaspora ne se résume pas à l'argent envoyé. Elle apporte aussi :

\begin{itemize}
\item des \cle{investissements productifs} : création d'entreprises, d'hôtels, de cliniques, de fermes agricoles par des Camerounais de retour ou restés à l'étranger ;
\item des \cle{transferts de compétences} : médecins de la diaspora qui viennent opérer gratuitement pendant leurs vacances, ingénieurs qui montent des formations, professeurs qui donnent des cours dans les universités camerounaises ;
\item des \cle{retours d'expérience} : certains migrants reviennent définitivement avec un savoir-faire, un réseau de contacts et un capital, et deviennent entrepreneurs ;
\item un \cle{rayonnement culturel} : artistes, sportifs et intellectuels camerounais de l'étranger font connaître le pays dans le monde.
\end{itemize}

On parle alors de \emph{gain de cerveaux} (\emph{brain gain}) lorsque le retour des compétences compense, au moins partiellement, la fuite initiale.

\courssub{L'immigration vers le Cameroun : le pays est aussi une terre d'accueil}

Le Cameroun n'est pas seulement un pays de départ : sa relative stabilité politique et sa position de carrefour en Afrique centrale en font aussi un pays d'\cle{immigration}. On y trouve notamment :

\begin{itemize}
\item des \textbf{travailleurs nigérians}, très présents dans le \cle{commerce} (marchés de Douala, Buea, Bamenda), l'artisanat et les services ;
\item des \textbf{éleveurs et commerçants tchadiens}, actifs dans l'\cle{élevage} bovin, le commerce de bétail et les transports dans le Nord et l'Adamaoua ;
\item des \textbf{Centrafricains}, présents dans le commerce, l'élevage et l'artisanat à l'Est et dans les grandes villes, auxquels s'ajoutent des centaines de milliers de \cle{réfugiés} fuyant les crises de leur pays ;
\item des ressortissants d'autres pays (Maliens, Sénégalais, Congolais) dans le commerce et les petits métiers.
\end{itemize}

Ces immigrés occupent souvent des emplois que les nationaux délaissent, contribuent à l'approvisionnement des marchés et paient des impôts et taxes. Mais leur présence peut aussi susciter des tensions : concurrence commerciale, conflits entre éleveurs et agriculteurs, et parfois des discours de rejet. L'immigration vers le Cameroun est donc, elle aussi, un phénomène à double face.

\courssub{Bilan et application}

\begin{aretenir}[Gains et pertes des migrations internationales pour le Cameroun]
\textbf{Pertes} : fuite des cerveaux (médecins, ingénieurs, enseignants), coût de formation perdu pour l'État, morts et drames de la migration clandestine (Sahara, Méditerranée), trafic d'êtres humains, déstructuration de familles.
\textbf{Gains} : transferts de fonds massifs (construction, éducation, commerce), investissements de la diaspora, transferts de compétences et retours d'expérience, apports des travailleurs immigrés (Nigérians, Tchadiens, Centrafricains) dans le commerce et l'élevage.
\end{aretenir}

\begin{exoresolu}[Calculer et commenter un taux de variation]
\textbf{Énoncé.} Les transferts de fonds de la diaspora camerounaise sont estimés à 60 milliards de FCFA en 2010 et à 220 milliards de FCFA en 2023. 1) Calculez le taux de variation de ces transferts entre 2010 et 2023. 2) Dégagez deux conséquences de cette évolution pour les familles camerounaises.
\tcblower
\textbf{Solution détaillée.}
1) Le taux de variation se calcule par la formule :
\[
t = \frac{V_{2023} - V_{2010}}{V_{2010}} \times 100 = \frac{220 - 60}{60} \times 100 = \frac{160}{60} \times 100 \approx 266,7\,\%.
\]
Les transferts de fonds ont donc augmenté d'environ \textbf{267 \%} entre 2010 et 2023 : ils ont été multipliés par plus de 3,5 (car $220 / 60 \approx 3,7$).

2) Conséquences pour les familles :
\begin{itemize}
\item \textbf{Conséquence positive} : davantage de familles peuvent financer la construction de maisons, la scolarisation des enfants et les soins de santé, ce qui améliore leur niveau de vie et stimule le commerce local.
\item \textbf{Conséquence ambivalente} : cette manne crée une \cle{dépendance} des familles à l'égard des proches de l'étranger et peut décourager l'initiative locale ; elle traduit aussi l'ampleur du départ des forces vives du pays.
\end{itemize}
\end{exoresolu}

\begin{exercice}[Pour s'entraîner]
1) Définissez la fuite des cerveaux et citez deux secteurs camerounais qu'elle affecte.
2) Reconstituez, sur une carte muette de l'Afrique du Nord et de l'Afrique de l'Ouest, l'itinéraire d'un migrant clandestin camerounais vers l'Italie.
3) Rédigez un paragraphe argumenté (10 à 15 lignes) répondant à la question : « La diaspora camerounaise est-elle une chance ou un problème pour le développement du pays ? » Vous justifierez votre conclusion par au moins deux arguments chiffrés.
\end{exercice}

\begin{corrige}[Exercice]
1) La fuite des cerveaux est l'émigration des personnes qualifiées vers des pays plus développés, attirées par de meilleurs salaires et conditions de travail. Secteurs affectés : la santé (médecins, infirmiers) et l'éducation (enseignants, universitaires), ainsi que l'ingénierie.
2) Itinéraire attendu : Cameroun $\rightarrow$ Nigeria $\rightarrow$ Niger (Agadez) $\rightarrow$ traversée du Sahara $\rightarrow$ Libye $\rightarrow$ traversée de la Méditerranée $\rightarrow$ Italie. L'élève doit flécher chaque étape et orienter sa carte (flèche du nord).
3) Le paragraphe doit présenter une thèse nuancée : chance (transferts de fonds de plusieurs centaines de milliards de FCFA par an, supérieurs certaines années aux investissements directs étrangers ; investissements et transferts de compétences) et problème (fuite des cerveaux, coût de formation perdu, drames de la clandestinité avec plus de \num{3000} morts en Méditerranée en 2023). La conclusion personnelle doit être justifiée par les arguments développés.
\end{corrige}

Ainsi, les migrations internationales placent le Cameroun dans une situation paradoxale : le pays perd ses cerveaux et parfois ses enfants dans le désert et la mer, mais il reçoit en retour des milliards de francs et des compétences de sa diaspora, tout en accueillant lui-même des travailleurs étrangers. Reste à savoir comment les États et les organisations internationales tentent d'organiser et de réguler ces mouvements : ce sera l'objet de la section suivante, consacrée aux tentatives de règlement des problèmes migratoires.

\coursec{Les tentatives de règlement des problèmes migratoires}

Nous avons vu dans la section précédente que les migrations internationales procurent au Cameroun des gains réels (transferts d'argent, formation de la diaspora) mais aussi des pertes sensibles (fuite des cerveaux, départ des forces vives). Face à ces effets contradictoires, les États, les organisations régionales et les institutions internationales ne restent pas passifs : ils tentent d'encadrer, de canaliser, parfois de freiner les mouvements migratoires. C'est l'objet de cette dernière section : comprendre \emph{qui} agit, \emph{comment}, et avec \emph{quels résultats}.

\courssub{Les actions menées au niveau national}

\courssub{Retenir les populations rurales : l'aménagement du territoire}

L'intuition de départ est simple : si les jeunes quittent les campagnes, c'est parce que la vie y est difficile (pas d'électricité, pas de routes, pas de marchés pour écouler les récoltes). Pour freiner l'\cle{exode rural}, l'État camerounais cherche donc à rendre les campagnes plus attractives. C'est la logique de l'\cle{aménagement du territoire}.

\begin{definition}[Aménagement du territoire]
L'\cle{aménagement du territoire} désigne l'ensemble des politiques publiques visant à organiser l'espace national (équipements, infrastructures, services) afin de réduire les déséquilibres entre les régions et de fixer les populations sur place.
\end{definition}

Les principales mesures camerounaises sont :

\begin{itemize}
\item l'\cle{électrification rurale} : le Programme national d'électrification rurale, porté notamment par l'Agence d'électrification rurale (AER), vise à raccorder les villages au réseau électrique ; l'électricité permet la conservation des produits, l'artisanat, les petites industries locales ;
\item les \cle{pistes rurales} : leur réhabilitation désenclave les bassins de production agricole et permet aux paysans d'écouler leurs récoltes vers les villes à moindre coût ;
\item les \cle{agropôles} : il s'agit de pôles agro-industriels qui regroupent production, transformation et commercialisation dans une même zone rurale, afin de créer des emplois hors des grandes villes ;
\item l'accès à l'eau potable, aux centres de santé et aux écoles, qui réduit l'écart de niveau de vie entre ville et campagne.
\end{itemize}

\begin{exemplebox}[Les agropôles au Cameroun]
Le programme des agropôles lancé dans les années 2010 prévoyait des pôles de production et de transformation dans plusieurs régions (par exemple dans le Nord et à l'Ouest). L'objectif affiché était de créer des milliers d'emplois ruraux et de retenir les jeunes agriculteurs tentés par le départ vers Douala ou Yaoundé.
\end{exemplebox}

\courssub{Décentralisation et pôles de développement régionaux}

La deuxième stratégie consiste à déconcentrer le développement. Au Cameroun, la \cle{décentralisation} (loi de 2004, régions instituées par la Constitution) transfère des compétences et des ressources aux collectivités territoriales : régions et communes. L'idée est de créer des \cle{pôles de développement régionaux} (chefs-lieux de régions comme Garoua, Maroua, Bafoussam, Bertoua) capables d'offrir emplois et services, afin que les migrations ne se concentrent pas uniquement sur le couple Douala--Yaoundé.

\begin{attention}[Piège fréquent]
Ne confonds pas \emph{déconcentration} (transfert de services de l'État vers ses représentants locaux, comme les gouverneurs) et \emph{décentralisation} (transfert de pouvoirs à des collectivités élues, comme les conseils régionaux et municipaux). Au Probatoire, cette distinction est souvent attendue.
\end{attention}

\courssub{Contrôler les flux : visas, séjour, lutte contre l'émigration clandestine}

L'État agit aussi par le contrôle :

\begin{itemize}
\item les \cle{visas} et les \cle{cartes de séjour} encadrent l'entrée et le séjour des étrangers sur le territoire ;
\item la \cle{lutte contre l'émigration clandestine} : sensibilisation aux dangers de la traversée du Sahara et de la Méditerranée, démantèlement des réseaux de passeurs ;
\item les \cle{rapatriements} : le Cameroun a organisé le retour volontaire de ressortissants bloqués en Libye, avec l'appui de l'OIM (plusieurs milliers de migrants camerounais rapatriés depuis 2017).
\end{itemize}

\courssub{Les actions au niveau régional : CEMAC et Union africaine}

\begin{savaistu}[La libre circulation dans la CEMAC]
La Communauté économique et monétaire de l'Afrique centrale (CEMAC), qui regroupe le Cameroun, le Tchad, la Centrafrique, le Gabon, le Congo et la Guinée équatoriale, autorise \emph{en principe} la libre circulation des personnes entre ses États membres. En pratique, des tracasseries aux frontières et des exigences administratives limitent souvent cette liberté : c'est un bel exemple d'écart entre le droit et la réalité.
\end{savaistu}

Au niveau continental, l'\cle{Union africaine} (UA) a adopté un Cadre de politique migratoire pour l'Afrique (révisé en 2018) et promeut la libre circulation des personnes, notamment à travers le protocole de 2018 sur la libre circulation et le passeport africain. L'objectif est de faire des migrations intra-africaines un facteur de développement plutôt qu'un problème.

\courssub{Les actions au niveau international}

\begin{definition}[Réfugié]
Un \cle{réfugié} est une personne qui a fui son pays en raison de persécutions (pour sa race, sa religion, sa nationalité, ses opinions politiques) ou de conflits, et qui ne peut ou ne veut pas y retourner par crainte pour sa vie. Son statut est protégé par le droit international, notamment la Convention de Genève de 1951 : on ne peut pas renvoyer un réfugié vers un pays où sa vie est menacée (principe de \emph{non-refoulement}).
\end{definition}

Trois acteurs internationaux jouent un rôle majeur :

\begin{itemize}
\item l'\cle{OIM} (Organisation internationale pour les migrations) : elle aide au retour volontaire des migrants, à leur réinsertion, et fournit des données sur les flux ;
\item le \cle{HCR} (Haut-Commissariat des Nations unies pour les réfugiés) : il protège les réfugiés, gère les camps avec les États d'accueil et cherche des solutions durables (retour, intégration locale, réinstallation) ;
\item le \cle{Pacte mondial sur les migrations} (2018, Pacte de Marrakech) : premier accord international non contraignant qui propose un cadre de coopération pour des migrations « sûres, ordonnées et régulières » ;
\item les \cle{accords bilatéraux} : accords de réadmission ou de gestion des flux signés entre deux États (par exemple entre pays africains et européens).
\end{itemize}

\begin{popfigure}
\begin{center}
\begin{tikzpicture}[
  font=\small,
  niveau/.style={draw=popdark, fill=popblueL, rounded corners, minimum width=3.2cm, minimum height=1cm, align=center, font=\bfseries},
  mesure/.style={draw=popline, fill=popcream, rounded corners, minimum width=6.2cm, minimum height=1cm, align=left, text width=6cm},
  lien/.style={-{Stealth[length=2.5mm]}, thick, popblue}
]
\node[niveau, align=center] (nat) at (0,4){Niveau national\\ (État camerounais)};
\node[niveau, align=center] (reg) at (0,2){Niveau régional\\ (CEMAC, UA)};
\node[niveau, align=center] (int) at (0,0){Niveau international\\ (OIM, HCR, ONU)};

\node[mesure, right=1.2cm of nat] (mnat) {Aménagement du territoire (électrification rurale, pistes, agropôles) ; décentralisation ; visas, cartes de séjour, rapatriements};
\node[mesure, right=1.2cm of reg] (mreg) {Libre circulation des personnes (CEMAC) ; cadre migratoire de l'Union africaine ; protocole UA de 2018};
\node[mesure, right=1.2cm of int] (mint) {Retours volontaires (OIM) ; protection des réfugiés (HCR) ; Pacte mondial de 2018 ; accords bilatéraux};

\draw[lien] (nat) -- (mnat);
\draw[lien] (reg) -- (mreg);
\draw[lien] (int) -- (mint);
\end{tikzpicture}
\end{center}
\legende{Les trois niveaux de gestion des problèmes migratoires : à chaque échelle correspondent des acteurs et des instruments différents, qui se complètent.}
\end{popfigure}

\courssub{Étude de cas : la gestion des réfugiés au Cameroun}

Le Cameroun n'est pas seulement un pays de départ : c'est aussi un pays d'\emph{accueil}. Deux grandes crises expliquent la présence de réfugiés sur son sol :

\begin{itemize}
\item la crise liée à la secte \cle{Boko Haram} dans le bassin du lac Tchad, qui a poussé des Nigérians vers l'Extrême-Nord camerounais ;
\item la crise politico-militaire en \cle{Centrafrique}, qui a conduit des dizaines de milliers de Centrafricains vers les régions de l'Est et de l'Adamaoua.
\end{itemize}

\begin{exemplebox}[Le camp de Minawao]
Le camp de \cle{Minawao} (département du Mayo-Tsanaga, Extrême-Nord), ouvert en 2013 et géré par le HCR avec le gouvernement camerounais, accueille des \textbf{dizaines de milliers de réfugiés nigérians} (plus de 60 000 à certains moments) fuyant les exactions de Boko Haram. Les réfugiés y reçoivent abri, nourriture, soins et scolarisation des enfants. Dans l'Est, les réfugiés centrafricains sont répartis entre des sites organisés (comme Gado, Lolo ou Mbile) et les villages d'accueil.
\end{exemplebox}

\begin{popfigure}
\begin{center}
\begin{tikzpicture}[scale=1.05, font=\small]
% Contour très simplifié du Cameroun
\draw[thick, popdark, fill=popcream]
  (2.2,6.4) -- (3.0,6.9) -- (3.6,6.3) -- (4.1,5.6) -- (4.4,4.9)
  -- (4.0,4.2) -- (4.3,3.4) -- (3.9,2.4) -- (3.4,1.2) -- (2.9,0.3)
  -- (2.3,0.1) -- (1.8,0.6) -- (1.2,1.0) -- (1.0,1.9) -- (1.3,2.8)
  -- (1.1,3.7) -- (1.5,4.5) -- (1.9,5.2) -- cycle;
% Nord (flèche)
\draw[-{Stealth[length=2.5mm]}, thick, popdark] (5.4,1.0) -- (5.4,1.8);
\node[below] at (5.4,1.0) {\textbf{N}};
% Camp de Minawao
\node[circle, fill=poporange, inner sep=2.5pt] (min) at (3.15,6.15) {};
\node[right, align=left] at (3.35,6.35) {\textbf{Minawao}\\ (réfugiés nigérians)};
% Sites de l'Est
\node[circle, fill=poppurple, inner sep=2.5pt] at (3.9,2.9) {};
\node[circle, fill=poppurple, inner sep=2.5pt] at (3.6,2.2) {};
\node[right, align=left] at (4.05,2.55) {\textbf{Sites de l'Est}\\ (Gado, Lolo, Mbile :\\ réfugiés centrafricains)};
% Villes repères
\node[circle, fill=popblue, inner sep=1.8pt] at (2.1,1.3) {};
\node[left] at (2.05,1.3) {Douala};
\node[circle, fill=popblue, inner sep=1.8pt] at (2.6,1.9) {};
\node[right] at (2.7,1.9) {Yaoundé};
\node[circle, fill=popblue, inner sep=1.8pt] at (3.2,5.4) {};
\node[left] at (3.15,5.4) {Maroua};
% Flèches d'origine
\draw[-{Stealth[length=2.5mm]}, very thick, poporange, dashed] (2.2,7.4) -- (3.05,6.35);
\node[above left, poporange] at (2.2,7.4) {Nigeria};
\draw[-{Stealth[length=2.5mm]}, very thick, poppurple, dashed] (5.3,2.0) -- (4.05,2.5);
\node[right, poppurple] at (5.3,2.0) {Centrafrique};
\end{tikzpicture}
\end{center}
\legende{Localisation schématique des principaux sites d'accueil de réfugiés au Cameroun : le camp de Minawao (Extrême-Nord) et les sites de l'Est. Croquis simplifié, sans échelle exacte.}
\end{popfigure}

\courssub{Les limites et les difficultés de ces politiques}

Malgré cette multiplication des acteurs et des mesures, les résultats restent mitigés. Plusieurs limites s'imposent à l'analyse :

\begin{itemize}
\item des \cle{moyens financiers et humains insuffisants} : les agropôles tardent, l'électrification rurale progresse lentement, les camps dépendent de l'aide internationale souvent sous-financée ;
\item la \cle{corruption} et les tracasseries (contrôles abusifs aux frontières, détournements), qui vident de leur contenu la libre circulation pourtant garantie par la CEMAC ;
\item la \cle{persistance des causes profondes} : tant que la pauvreté rurale, le chômage des jeunes, les conflits (Boko Haram, crise centrafricaine) et les inégalités entre régions ne seront pas traités, les migrations continueront, y compris clandestines ;
\item la tentation sécuritaire : fermer les frontières déplace les routes migratoires vers des trajets plus dangereux sans tarir les flux.
\end{itemize}

\begin{methode}[Rédiger une conclusion argumentée sur les migrations]
\begin{enumerate}
\item \textbf{Rappeler le problème} : les migrations ont des effets contrastés (positifs et négatifs) sur les régions de départ et d'accueil.
\item \textbf{Présenter les mesures} : citer les réponses aux trois niveaux (national, régional, international) avec un exemple précis chacun.
\item \textbf{Évaluer leur efficacité} : nuancer (moyens limités, corruption, causes profondes persistantes).
\item \textbf{Proposer une piste} : par exemple, investir davantage dans le développement rural et la formation des jeunes pour transformer la migration en choix plutôt qu'en contrainte.
\end{enumerate}
\end{methode}

\begin{exoresolu}[Les politiques migratoires sont-elles efficaces ?]
\textbf{Énoncé.} À partir de vos connaissances, rédigez un paragraphe argumenté (8 à 12 lignes) répondant à la question : « Les tentatives de règlement des problèmes migratoires au Cameroun sont-elles efficaces ? »
\tcblower
\textbf{Solution détaillée.} Le Cameroun et ses partenaires ont mis en place de nombreuses réponses aux problèmes migratoires. Au niveau national, l'aménagement du territoire (électrification rurale, pistes rurales, agropôles) et la décentralisation visent à retenir les populations rurales et à créer des pôles de développement régionaux ; l'État contrôle aussi les flux par les visas, les cartes de séjour et les rapatriements de migrants clandestins. Au niveau régional, la CEMAC garantit en principe la libre circulation des personnes, et l'Union africaine promeut une politique migratoire commune. Au niveau international, l'OIM organise les retours volontaires, le HCR protège les réfugiés — comme les dizaines de milliers de Nigérians du camp de Minawao — et le Pacte mondial de 2018 encadre la coopération. Cependant, ces politiques restent limitées par le manque de moyens, la corruption et la persistance des causes profondes (pauvreté, chômage, conflits). Elles sont donc partiellement efficaces : une solution durable suppose de traiter les racines des migrations, notamment par le développement rural et l'emploi des jeunes.
\end{exoresolu}

\begin{exercice}[À toi de jouer]
\begin{enumerate}
\item Définis les termes suivants : réfugié, aménagement du territoire, agropôle.
\item Cite deux mesures prises au niveau national pour freiner l'exode rural au Cameroun.
\item Quel est le rôle du HCR dans la gestion du camp de Minawao ?
\item Explique en quelques lignes pourquoi la libre circulation dans la CEMAC reste imparfaite en pratique.
\end{enumerate}
\end{exercice}

\begin{corrige}[Exercice]
\begin{enumerate}
\item Réfugié : personne ayant fui son pays en raison de persécutions ou de conflits, protégée par le droit international (Convention de Genève, 1951). Aménagement du territoire : ensemble des politiques publiques d'organisation de l'espace national pour réduire les déséquilibres régionaux et fixer les populations. Agropôle : pôle agro-industriel regroupant production, transformation et commercialisation en zone rurale pour créer des emplois.
\item L'électrification rurale et la réhabilitation des pistes rurales (on peut aussi citer les agropôles, l'eau potable, les centres de santé).
\item Le HCR protège les réfugiés nigérians, co-gère le camp avec l'État camerounais (abri, nourriture, soins, scolarisation) et recherche des solutions durables (retour volontaire, intégration).
\item Parce que malgré les textes, les frontières restent marquées par des tracasseries, des contrôles abusifs et des exigences administratives liées notamment à la corruption et aux moyens limités des États.
\end{enumerate}
\end{corrige}

\begin{aretenir}[L'essentiel de la section]
Les problèmes migratoires sont traités à trois niveaux complémentaires : \textbf{national} (aménagement du territoire, décentralisation, contrôle des flux), \textbf{régional} (libre circulation CEMAC, cadre de l'Union africaine) et \textbf{international} (OIM, HCR, Pacte mondial de 2018, accords bilatéraux). Le Cameroun illustre aussi l'accueil des réfugiés (Minawao, sites de l'Est). Mais ces politiques butent sur le manque de moyens, la corruption et la persistance des causes profondes : réguler les migrations suppose d'abord de développer les régions de départ.
\end{aretenir}

\coursec{TP guidé : cartographier et analyser un flux migratoire}

Après avoir étudié les tentatives de règlement des problèmes migratoires (politiques nationales, accords bilatéraux, rôle de l'OIM), il est temps de passer de la théorie à la pratique. Ce travail pratique guidé vous place dans la peau d'un géographe : à partir de données réelles simplifiées, vous allez \cle{construire une carte de flux}, \cle{réaliser un diagramme statistique} et \cle{rédiger un commentaire argumenté}. Ce sont exactement les compétences exigées à l'épreuve de Géographie du Probatoire et du Baccalauréat technique, où l'exploitation de documents (tableaux, cartes, graphiques) vaut une part importante de la note.

\courssub{Objectifs et déroulement du TP}

\begin{definition}[Objectifs du TP]
À l'issue de ce TP, l'élève doit être capable de :
\begin{itemize}
\item lire et trier un tableau de données de flux migratoires interrégionaux ;
\item construire une \cle{carte de flux} à flèches proportionnelles (fond de carte, échelle, titre, légende, source) ;
\item construire un \cle{diagramme en barres} à partir de données chiffrées ;
\item rédiger un commentaire structuré de 10 à 15 lignes dégageant les problèmes et proposant des solutions.
\end{itemize}
\end{definition}

\textbf{Déroulement :} le TP se déroule en cinq étapes, sur deux séances. La séance 1 est consacrée aux étapes 1 à 3 (lecture des données, carte, diagramme). La séance 2 est consacrée aux étapes 4 et 5 (commentaire écrit, restitution orale et correction collective).

\courssub{Étape 1 : lecture du tableau de données}

Le professeur distribue le tableau suivant, qui présente les flux migratoires interrégionaux vers les deux métropoles du Cameroun (données simplifiées et arrondies à visée pédagogique, inspirées des résultats du RGPH 2005).

\begin{center}
\begin{tabularx}{\linewidth}{|l|X|X|X|}\hline
\rowcolor{popblueL} \textbf{Région de départ} & \textbf{Vers Douala (Littoral)} & \textbf{Vers Yaoundé (Centre)} & \textbf{Total des départs} \\ \hline
Extrême-Nord & 25 000 & 95 000 & 120 000 \\ \hline
Nord & 20 000 & 55 000 & 75 000 \\ \hline
Adamaoua & 10 000 & 40 000 & 50 000 \\ \hline
Nord-Ouest & 85 000 & 60 000 & 145 000 \\ \hline
Ouest & 110 000 & 70 000 & 180 000 \\ \hline
Sud-Ouest & 60 000 & 15 000 & 75 000 \\ \hline
Est & 8 000 & 30 000 & 38 000 \\ \hline
Sud & 12 000 & 45 000 & 57 000 \\ \hline
\end{tabularx}
\end{center}

\begin{experience}[Lecture guidée du tableau]
Avant de cartographier, il faut \emph{comprendre} les données. Répondez par écrit aux questions suivantes :
\begin{enumerate}
\item Quelle région fournit le plus grand nombre de migrants ? (\emph{Réponse attendue : l'Ouest, avec 180 000 départs.})
\item Quelle ville polarise le plus les flux ? (\emph{Réponse : Douala et Yaoundé se partagent les flux ; Douala attire surtout l'Ouest, le Nord-Ouest et le Sud-Ouest, Yaoundé attire surtout le Grand Nord.})
\item Calculez la part des trois régions du Grand Nord (Extrême-Nord, Nord, Adamaoua) dans le total des départs. (\emph{Réponse : $120 000 + 75 000 + 50 000 = 245 000$ sur $740 000$, soit environ 33 \%.})
\item Classez les flux de la région de départ la plus émettrice à la moins émettrice : ce classement servira à l'étape 2.
\end{enumerate}
\textbf{Interprétation :} on observe que les régions à forte pression démographique et à faibles opportunités économiques (Ouest, Grand Nord) sont les principaux pourvoyeurs de migrants, ce qui confirme l'étude des effets migratoires menée dans les sections précédentes.
\end{experience}

\courssub{Étape 2 : construction de la carte de flux}

Une \cle{carte de flux} représente des déplacements par des flèches dont \textbf{l'épaisseur est proportionnelle à l'importance du flux}. C'est l'outil privilégié du géographe pour montrer d'un coup d'œil qui part, où, et en quelle quantité.

\begin{methode}[Construire une carte de flux en cinq étapes]
\begin{enumerate}
\item \textbf{Trier les données} : classer les flux du plus fort au plus faible et repérer les flux majeurs (ceux que la carte doit faire ressortir).
\item \textbf{Choisir une échelle de flèches} : fixer une correspondance simple, par exemple 1 mm d'épaisseur pour 20 000 migrants. Ainsi un flux de 110 000 personnes = 5.5mm, un flux de 10 000 = 0.5mm. Prévoir 3 ou 4 classes d'épaisseur maximum pour garder une carte lisible.
\item \textbf{Tracer le fond de carte} : contour simplifié du Cameroun, limites des régions concernées, points pour les villes d'accueil (Douala, Yaoundé), flèche du nord et échelle graphique indicative.
\item \textbf{Tracer les flèches} : de chaque région de départ vers la ville d'accueil, avec l'épaisseur calculée. Ne représenter que les flux significatifs (par exemple supérieurs à 40 000) pour éviter la surcharge.
\item \textbf{Titrer, légender, sourcer} : un titre précis (quoi, où, quand), une légende donnant l'échelle des flèches, et la source des données.
\end{enumerate}
\end{methode}

\begin{attention}[Erreurs fréquentes à éviter]
L'erreur la plus fréquente au Probatoire et au Bac est d'oublier le \textbf{titre}, la \textbf{légende} ou la \textbf{source} sur la carte : cela coûte systématiquement des points, même si le tracé est correct. Autres pièges : des flèches dont l'épaisseur ne respecte pas l'échelle annoncée, l'absence d'orientation (flèche du nord), ou une carte surchargée de trop nombreux flux devenus illisibles. Vérifiez toujours votre carte avec la check-list : titre ? légende ? source ? nord ? échelle ?
\end{attention}

\begin{popfigure}
\begin{center}
\begin{tikzpicture}[scale=1.0]
% Fond de carte simplifié du Cameroun
\draw[thick, popdark, fill=popcream] (2.2,5.6) -- (3.4,6.6) -- (4.6,6.1) -- (5.0,5.0) -- (5.6,4.2) -- (5.4,3.0) -- (4.8,2.2) -- (4.6,1.0) -- (3.6,0.2) -- (2.4,0.4) -- (1.4,1.2) -- (0.6,1.0) -- (0.4,2.0) -- (1.0,3.0) -- (1.4,4.0) -- (1.8,4.8) -- cycle;
% Villes d'accueil
\fill[popink] (1.0,1.6) circle (0.09) node[below left, font=\small] {Douala};
\fill[popink] (3.0,1.8) circle (0.09) node[below right, font=\small] {Yaoundé};
% Régions de départ (points)
\fill[poporange] (3.6,5.6) circle (0.07) node[above, font=\scriptsize] {Extrême-Nord};
\fill[poporange] (2.6,4.6) circle (0.07) node[above left, font=\scriptsize] {Nord};
\fill[poporange] (4.2,3.6) circle (0.07) node[right, font=\scriptsize] {Adamaoua};
\fill[poporange] (1.2,3.6) circle (0.07) node[left, font=\scriptsize] {Nord-Ouest};
\fill[poporange] (1.9,2.9) circle (0.07) node[above right, font=\scriptsize] {Ouest};
\fill[poporange] (0.7,2.5) circle (0.07) node[left, font=\scriptsize] {Sud-Ouest};
\fill[poporange] (5.0,2.5) circle (0.07) node[right, font=\scriptsize] {Est};
\fill[poporange] (2.5,0.5) circle (0.07) node[below, font=\scriptsize] {Sud};
% Flèches de flux (épaisseurs proportionnelles, échelle : 1 mm = 20 000 migrants)
% Seuls les flux >= 40 000 sont tracés (épaisseur >= 2 mm)
\draw[-{Stealth[length=3mm]}, line width=5.5mm, popblue, opacity=0.75] (1.9,2.9) -- (1.05,1.7);     % Ouest -> Douala 110 000
\draw[-{Stealth[length=3mm]}, line width=4.75mm, popblue, opacity=0.75] (3.6,5.5) -- (3.0,1.95);   % Extrême-Nord -> Yaoundé 95 000
\draw[-{Stealth[length=3mm]}, line width=4.25mm, popblue, opacity=0.75] (1.2,3.55) -- (1.0,1.75); % Nord-Ouest -> Douala 85 000
\draw[-{Stealth[length=3mm]}, line width=3.5mm, popblue, opacity=0.75] (1.9,2.85) -- (2.9,1.9);   % Ouest -> Yaoundé 70 000
\draw[-{Stealth[length=3mm]}, line width=3.0mm, popblue, opacity=0.75] (1.25,3.6) -- (2.9,1.85);  % Nord-Ouest -> Yaoundé 60 000
\draw[-{Stealth[length=3mm]}, line width=3.0mm, popblue, opacity=0.75] (0.75,2.45) -- (1.0,1.7);   % Sud-Ouest -> Douala 60 000
\draw[-{Stealth[length=3mm]}, line width=2.75mm, popblue, opacity=0.75] (2.6,4.55) -- (2.95,1.85); % Nord -> Yaoundé 55 000
\draw[-{Stealth[length=3mm]}, line width=2.25mm, popblue, opacity=0.75] (2.55,0.55) -- (2.95,1.8);  % Sud -> Yaoundé 45 000
\draw[-{Stealth[length=3mm]}, line width=2.0mm, popblue, opacity=0.75] (4.2,3.55) -- (3.05,1.85);  % Adamaoua -> Yaoundé 40 000
% Nord
\draw[->, thick, popdark] (5.9,5.2) -- (5.9,6.0) node[above, font=\small] {N};
% Échelle graphique indicative
\draw[thick] (0.4,-0.5) -- (1.4,-0.5) node[midway, below, font=\scriptsize] {100 km (indicatif)};
% Légende des flèches (échelle 1 mm = 20 000)
\draw[line width=5.5mm, popblue, opacity=0.75] (2.6,-0.7) -- (3.4,-0.7) node[right, font=\scriptsize, black] {110 000 migrants (5,5 mm)};
\draw[line width=3.0mm, popblue, opacity=0.75] (2.6,-1.3) -- (3.4,-1.3) node[right, font=\scriptsize, black] {60 000 migrants (3,0 mm)};
\draw[line width=1.0mm, popblue, opacity=0.75] (2.6,-1.9) -- (3.4,-1.9) node[right, font=\scriptsize, black] {20 000 migrants (1,0 mm)};
\end{tikzpicture}
\end{center}
\legende{Carte de flux corrigée : les principaux flux migratoires interrégionaux vers Douala et Yaoundé (flux $\ge$ 40 000 migrants). Échelle des flèches : 1 mm d'épaisseur pour 20 000 migrants. Source : données simplifiées d'après RGPH 2005.}
\end{popfigure}

\courssub{Étape 3 : le diagramme en barres des transferts de fonds}

Les migrations internationales se traduisent par des \cle{transferts de fonds} (\emph{remittances}) de la diaspora vers le Cameroun. Pour comparer ces montants selon la région de résidence de la diaspora, le graphique adapté est le \textbf{diagramme en barres} : chaque barre a une hauteur proportionnelle au montant transféré.

\begin{methode}[Construire un diagramme en barres]
\begin{enumerate}
\item Tracer deux axes perpendiculaires : en abscisse les catégories (régions de la diaspora), en ordonnée la grandeur mesurée (montants en milliards de FCFA), avec une échelle régulière.
\item Graduer l'axe vertical en partant de 0 (sinon les comparaisons sont faussées).
\item Tracer des barres de même largeur, régulièrement espacées, de hauteur proportionnelle aux valeurs.
\item Indiquer le titre, le nom et l'unité de chaque axe, et la source des données.
\end{enumerate}
\end{methode}

\begin{popfigure}
\begin{center}
\begin{tikzpicture}
\begin{axis}[
    ybar,
    width=0.95\linewidth,
    height=7cm,
    bar width=0.9cm,
    ymin=0, ymax=140,
    ylabel={Montant (milliards de FCFA)},
    symbolic x coords={Europe, Amérique du Nord, Afrique, Asie},
    xtick=data,
    nodes near coords,
    every node near coord/.append style={font=\small},
    ymajorgrids=true,
    grid style={popline, dashed},
    axis line style={popdark},
    tick label style={font=\small},
    label style={font=\small},
]
\addplot[fill=popblue, draw=popdark] coordinates {(Europe,120) (Amérique du Nord,60) (Afrique,35) (Asie,10)};
\end{axis}
\end{tikzpicture}
\end{center}
\legende{Exemple de production finale attendue : transferts de fonds de la diaspora vers le Cameroun par région de résidence (estimations simplifiées, environ 225 milliards de FCFA au total). Source : données pédagogiques inspirées de la Banque mondiale.}
\end{popfigure}

La lecture du diagramme est immédiate : la diaspora installée en \textbf{Europe} (France, Allemagne, Belgique...) fournit à elle seule plus de la moitié des transferts (120 milliards sur 225, soit environ 53 \%), devant l'Amérique du Nord (environ 27 \%). Ces fonds financent la consommation, l'éducation, la santé et la construction dans les familles restées au pays : c'est l'un des principaux \emph{effets positifs} des migrations internationales pour le Cameroun.

\courssub{Étape 4 : rédaction du commentaire (10 à 15 lignes)}

Le commentaire est l'exercice roi de l'épreuve de Géographie. Il ne s'agit pas de décrire mécaniquement le document, mais de \textbf{dégager une idée principale} et de la discuter.

\begin{methode}[Commenter un document statistique ou cartographique]
\begin{enumerate}
\item \textbf{Dégager l'idée principale} : en une phrase, que montre le document ? (ex. : « Les flux migratoires camerounais convergent massivement vers Douala et Yaoundé. »)
\item \textbf{Citer des chiffres précis} : appuyer chaque affirmation sur au moins deux ou trois données du document (valeurs, pourcentages, classements).
\item \textbf{Expliquer} : relier les faits observés aux mécanismes étudiés en cours (exode rural, chômage, attraction des métropoles, transferts de fonds).
\item \textbf{Problématiser} : identifier les problèmes posés (dépeuplement des campagnes, surpopulation urbaine, fuite des cerveaux).
\item \textbf{Conclure par des solutions} : proposer des pistes réalistes de règlement (décentralisation, création d'emplois ruraux, politiques migratoires).
\end{enumerate}
\end{methode}

\begin{exoresolu}[Commentaire rédigé : corrigé type]
\textbf{Énoncé.} À partir de la carte de flux et du diagramme en barres construits aux étapes 2 et 3, rédigez un commentaire de 10 à 15 lignes dégageant les problèmes posés par les migrations au Cameroun et les solutions envisageables.
\tcblower
\textbf{Solution détaillée (corrigé type).}
Les documents étudiés montrent que les migrations au Cameroun sont fortement polarisées par les deux métropoles : Douala et Yaoundé concentrent l'essentiel des flux interrégionaux. Les régions de l'Ouest (180 000 départs), du Nord-Ouest (145 000) et du Grand Nord (245 000 départs cumulés, soit 33 \% du total) sont les principaux foyers d'émigration. Ce mouvement massif pose de graves problèmes : dans les régions de départ, l'exode prive les campagnes de bras valides et accentue le vieillissement ; dans les villes d'accueil, l'afflux de migrants provoque surpopulation, chômage, bidonvilles et insécurité. À l'échelle internationale, le départ des cadres (fuite des cerveaux) appauvrit le pays, même si les transferts de fonds de la diaspora — environ 225 milliards de FCFA, dont 53 \% venant d'Europe — soutiennent fortement les familles et l'économie nationale. Pour régler ces problèmes, l'État peut poursuivre la décentralisation et l'équipement des chefs-lieux régionaux, créer des emplois ruraux (agropastoral, agro-industrie) pour fixer les jeunes, et organiser les migrations internationales par des accords bilatéraux avec l'appui de l'OIM. En conclusion, la migration n'est pas un mal en soi : bien encadrée, elle peut devenir un levier de développement équilibré du territoire camerounais.
\emph{(Ce texte fait 14 lignes manuscrites environ : il dégage l'idée principale, cite des chiffres, explique, problématise et conclut par des solutions.)}
\end{exoresolu}

\courssub{Étape 5 : restitution orale et correction collective}

Chaque groupe présente sa carte et son diagramme en 3 minutes : un élève commente la carte, un autre le diagramme, un troisième lit le commentaire. La classe compare les productions au corrigé projeté, puis chacun s'auto-évalue à l'aide du barème suivant.

\begin{exemplebox}[Barème type (inspiré des épreuves du Bac technique), sur 20 points]
\begin{center}
\begin{tabularx}{\linewidth}{|l|X|c|}\hline
\rowcolor{popblueL} \textbf{Production} & \textbf{Critères évalués} & \textbf{Points} \\ \hline
Carte de flux & Exactitude des flèches et de l'échelle (4) ; titre, légende, source, nord (2) ; lisibilité et soin (2) & 8 \\ \hline
Diagramme en barres & Axes gradués et nommés (2) ; hauteurs exactes et titre (2) & 4 \\ \hline
Commentaire & Idée principale dégagée (2) ; chiffres cités (2) ; explication des problèmes (2) ; solutions proposées et conclusion (2) & 8 \\ \hline
\end{tabularx}
\end{center}
\end{exemplebox}

\begin{exercice}[Pour s'entraîner à la maison]
À partir du tableau de l'étape 1 :
\begin{enumerate}
\item Calculez la part (en \%) des flux à destination de Yaoundé dans le total des départs.
\item Construisez le diagramme en barres des totaux de départs par région (8 barres).
\item Rédigez en 5 lignes la légende complète de votre propre carte de flux (titre, échelle, source).
\end{enumerate}
\end{exercice}

\begin{corrige}[Exercice « Pour s'entraîner à la maison »]
\begin{enumerate}
\item Total vers Yaoundé : $95 + 55 + 40 + 60 + 70 + 15 + 30 + 45 = 410$ milliers de migrants. Total général : 740 000. Part : $\dfrac{410}{740} \times 100 \approx 55,4 \%$. Yaoundé attire donc un peu plus de la moitié des flux.
\item Barres (hauteurs en milliers) : Ouest 180 ; Nord-Ouest 145 ; Extrême-Nord 120 ; Nord 75 ; Sud-Ouest 75 ; Sud 57 ; Adamaoua 50 ; Est 38. Axe vertical gradué de 0 à 200, titre et source obligatoires.
\item Exemple : « Les flux migratoires interrégionaux vers Douala et Yaoundé (Cameroun). Échelle des flèches : 1 mm d'épaisseur = 20 000 migrants. Source : données simplifiées d'après RGPH 2005. »
\end{enumerate}
\end{corrige}

\begin{aretenir}[L'essentiel du TP]
Une carte de flux se construit en cinq étapes : trier les données, choisir une échelle de flèches, tracer le fond de carte, tracer les flèches proportionnelles, puis titrer, légender et sourcer. Un diagramme en barres exige des axes gradués depuis zéro, nommés avec leurs unités. Un bon commentaire dégage une idée principale, cite des chiffres, explique les problèmes et conclut par des solutions. Ces trois savoir-faire — cartographier, graphiquer, commenter — sont au cœur de l'épreuve de Géographie au Probatoire et au Baccalauréat technique.
\end{aretenir}

\newpage

\coursec{Activité d'intégration}

\courssub{Objectifs de l'activité}

À l'issue de cette activité d'intégration, l'élève doit être capable de :
\begin{itemize}
\item identifier et classer les \cle{effets des migrations} sur une région de départ (village) et sur une zone d'accueil (ville) à partir de documents réels ;
\item exploiter des données chiffrées (tableau statistique, montants de transferts) pour calculer des pourcentages et construire un diagramme ;
\item construire un \cle{croquis} ou un diagramme simple résumant le bilan migratoire ;
\item rédiger un \cle{texte argumenté} proposant des solutions réalistes inspirées des politiques de règlement des problèmes migratoires étudiées en cours ;
\item justifier une démarche, une réponse et une conclusion devant la classe (restitution orale et débat).
\end{itemize}

\courssub{Situation}

\textbf{Le village de Bangang-Fokam et l'exode de sa jeunesse.}

Bangang-Fokam est un village de la région de l'Ouest du Cameroun, situé à environ \SI{25}{km} de Bafoussam. La population y vit essentiellement de l'agriculture vivrière (maïs, haricot, pomme de terre) et d'un petit élevage. Depuis une dizaine d'années, les jeunes du village partent massivement vers Douala, la capitale économique, à la recherche d'un emploi. Le chef du village s'inquiète : \og Nos champs vieillissent, nos enfants partent. Est-ce une chance ou un problème pour notre village ? \fg

Le comité de développement du village a réalisé une enquête. Voici le dossier documentaire rassemblé.

\textbf{Document 1 : Les départs de jeunes de Bangang-Fokam vers Douala (2015--2024).}

\begin{center}
\begin{tabularx}{\linewidth}{|l|X|X|X|}\hline
\rowcolor{popblueL} \textbf{Année} & \textbf{Nombre de départs} & \textbf{Dont jeunes de 15 à 29 ans} & \textbf{Population totale du village} \\ \hline
2015 & 18 & 15 & \num{3200} \\ \hline
2017 & 25 & 21 & \num{3100} \\ \hline
2019 & 34 & 29 & \num{2950} \\ \hline
2021 & 47 & 40 & \num{2800} \\ \hline
2024 & 60 & 52 & \num{2600} \\ \hline
\end{tabularx}
\end{center}

\textbf{Document 2 : Photographies d'un quartier précaire de Douala.} Les jeunes arrivés du village s'installent dans un quartier non aménagé (habitat spontané en planches et tôles, ruelles non goudronnées, absence d'eau potable courante et d'évacuation des ordures). Beaucoup survivent grâce au secteur informel : vente à la sauvette, moto-taxi, chargement au port. Sur 10 jeunes partis, environ 4 seulement déclarent avoir un emploi stable.

\textbf{Document 3 : Les transferts d'argent.} En 2024, les migrants originaires de Bangang-Fokam ont envoyé au village environ 18 millions de FCFA (argent envoyé par mobile money ou par les voyageurs). Ces fonds ont permis de rénover 12 maisons, de payer la scolarité de 85 enfants et d'acheter deux motos. À l'échelle nationale, la diaspora camerounaise transfère chaque année plusieurs centaines de milliards de FCFA.

\textbf{Document 4 : Extrait d'un discours officiel.} \og L'émigration clandestine de notre jeunesse, par voie terrestre vers le Maghreb ou par mer, expose nos enfants à la mort et prive le pays de ses forces vives. Le gouvernement œuvre à créer des emplois, à soutenir l'entreprenariat des jeunes et à renforcer la coopération avec les pays de destination. \fg

\textbf{Problème posé :} le départ des jeunes de Bangang-Fokam est-il une chance ou un problème, pour le village comme pour la ville de Douala ? Quelles solutions réalistes pourraient retenir les jeunes au village ?

\courssub{Consignes}

Travail en groupes de 4 à 6 élèves. Chaque groupe remet une production écrite et présente ses résultats en 5 minutes.

\begin{enumerate}
\item À partir du document 1, calcule : a) l'évolution en pourcentage du nombre de départs entre 2015 et 2024 ; b) la part des jeunes de 15 à 29 ans dans les départs de 2024 ; c) la baisse en pourcentage de la population du village entre 2015 et 2024.
\item À l'aide des documents 1, 2 et 3, dresse en deux colonnes la liste des \cle{effets négatifs} et des \cle{effets positifs} des migrations : d'abord pour le village de départ, ensuite pour la ville d'accueil (Douala).
\item Construis un croquis ou un diagramme simple (schéma à flèches ou diagramme en barres) résumant le bilan migratoire entre Bangang-Fokam et Douala. N'oublie pas le titre, la légende et l'orientation.
\item Rédige un texte argumenté d'environ 15 lignes répondant au problème posé. Ton texte doit : a) présenter une conclusion claire (chance et/ou problème) ; b) proposer \textbf{deux solutions réalistes} pour retenir les jeunes au village, en t'appuyant sur les politiques de règlement des problèmes migratoires étudiées en cours ; c) justifier chaque solution par au moins un argument tiré du dossier ou de la leçon.
\item Restitution : chaque groupe présente son croquis et lit sa conclusion. La classe débat : les solutions proposées sont-elles réalistes ? Lesquelles sont prioritaires ?
\end{enumerate}

\courssub{Ressources à mobiliser}

\begin{aretenir}[Ce qu'il faut se rappeler]
\begin{itemize}
\item \cle{Effets dans les régions de départ} : perte de bras valides et de forces vives, vieillissement de la population, baisse de la production agricole, désorganisation des familles ; mais aussi réduction du chômage local et \cle{transferts d'argent} (remises des migrants) qui soutiennent l'économie du village.
\item \cle{Incidences dans les zones d'accueil} : surpopulation, croissance des quartiers précaires, chômage et secteur informel, pression sur les services (eau, santé, écoles), insécurité ; mais aussi main-d'œuvre abondante et bon marché, dynamisme économique, brassage culturel.
\item \cle{Tentatives de règlement} : politiques de création d'emplois et d'appui à l'entreprenariat des jeunes, développement rural (équipement des villages, pistes rurales, électrification), accords et coopération entre pays de départ et d'accueil, lutte contre l'émigration clandestine et sensibilisation.
\item \cle{Calculs utiles} : taux de variation $= \dfrac{\text{valeur finale} - \text{valeur initiale}}{\text{valeur initiale}} \times 100$ ; part en pourcentage $= \dfrac{\text{partie}}{\text{total}} \times 100$.
\end{itemize}
\end{aretenir}

\courssub{Démarche et corrigé commenté}

\begin{exoresolu}[Corrigé commenté de l'activité]
\textbf{Étape 1 : exploitation statistique du document 1.}

a) Évolution du nombre de départs entre 2015 et 2024 :
\[ \frac{60 - 18}{18} \times 100 = \frac{42}{18} \times 100 \approx 233{,}3\,\% \]
Le nombre de départs a été multiplié par plus de 3 (hausse d'environ 233 \%).

b) Part des jeunes de 15 à 29 ans dans les départs de 2024 :
\[ \frac{52}{60} \times 100 \approx 86{,}7\,\% \]
Près de 9 départs sur 10 concernent des jeunes : c'est bien un \cle{exode des forces vives}.

c) Baisse de la population du village entre 2015 et 2024 :
\[ \frac{2600 - 3200}{3200} \times 100 = \frac{-600}{3200} \times 100 = -18{,}75\,\% \]
Le village a perdu près de 19 \% de sa population en dix ans.

\emph{Justification :} ces calculs appliquent la formule du taux de variation et confirment chiffrablement l'inquiétude du chef du village.

\tcblower

\textbf{Étape 2 : tableau des effets (consigne 2).}

\begin{center}
\begin{tabularx}{\linewidth}{|X|X|}\hline
\rowcolor{popblueL} \textbf{Effets négatifs} & \textbf{Effets positifs} \\ \hline
\multicolumn{2}{|l|}{\textbf{Pour le village de départ (Bangang-Fokam)}} \\ \hline
Perte de bras valides (86,7 \% de jeunes) ; vieillissement ; baisse de la production agricole ; population en recul de 18,75 \% ; familles séparées. & Transferts de 18 millions de FCFA en 2024 : 12 maisons rénovées, 85 enfants scolarisés, achat de motos ; moins de chômage au village. \\ \hline
\multicolumn{2}{|l|}{\textbf{Pour la ville d'accueil (Douala)}} \\ \hline
Surpopulation et quartiers précaires ; chômage (6 jeunes sur 10 sans emploi stable) ; explosion du secteur informel ; pression sur l'eau, la santé, les écoles. & Main-d'œuvre jeune et bon marché (moto-taxis, port, petits commerces) ; dynamisme économique ; brassage culturel. \\ \hline
\end{tabularx}
\end{center}

\textbf{Étape 3 : croquis bilan (consigne 3).}

\begin{popfigure}
\begin{tikzpicture}[font=\small]
\draw[fill=popgreenL,draw=popgreen!60!black,rounded corners=6pt] (0,0) rectangle (3.4,1.6);
\node[align=center] at (1.7,0.8) {\textbf{Bangang-Fokam}\\ village de départ\\ (Ouest)};
\draw[fill=popblueL,draw=popblue!60!black,rounded corners=6pt] (9,0) rectangle (12.6,1.6);
\node[align=center] at (10.8,0.8) {\textbf{Douala}\\ ville d'accueil\\ (Littoral)};
\draw[-{Stealth[length=4mm]},line width=3pt,poporange] (3.5,1.1) -- (8.9,1.1) node[midway,above,align=center]{flux de jeunes migrants\\ 60 départs en 2024};
\draw[-{Stealth[length=4mm]},line width=2pt,poppurple] (8.9,0.4) -- (3.5,0.4) node[midway,below,align=center]{transferts d'argent\\ 18 millions FCFA (2024)};
\draw[-{Stealth},popdark] (11.8,2.2) -- (11.8,3.0) node[above]{N};
\end{tikzpicture}
\legende{Croquis bilan du flux migratoire Bangang-Fokam--Douala : 1. flux de personnes (flèche orange, épaisse) ; 2. contre-flux d'argent (flèche violette). Échelle indicative : environ \SI{280}{km} entre les deux localités.}
\end{popfigure}

\textbf{Étape 4 : texte argumenté (consigne 4) — modèle de réponse.}

Le départ des jeunes de Bangang-Fokam est à la fois une chance et un problème, mais le bilan est davantage négatif. C'est une chance car les migrants envoient de l'argent (18 millions de FCFA en 2024) qui rénove les maisons et scolarise 85 enfants, et car le village compte moins de chômeurs. C'est surtout un problème : le village a perdu 18,75 \% de sa population en dix ans, 86,7 \% des partants sont des jeunes de 15 à 29 ans, les champs manquent de bras et la production agricole recule. À Douala, ces jeunes grossissent les quartiers précaires et 6 sur 10 restent sans emploi stable.

Première solution : développer le village pour retenir les jeunes, conformément aux politiques de développement rural. L'État et la commune peuvent électrifier le village, entretenir les pistes rurales et créer des unités de transformation (séchage du maïs, conservation de la pomme de terre) : des jeunes restent quand le travail rapporte, comme le montre le succès des deux motos achetées grâce aux transferts.

Deuxième solution : soutenir l'entreprenariat agricole des jeunes, comme le préconise le discours officiel (création d'emplois, appui à l'entreprenariat). Des microcrédits et une formation permettraient aux jeunes de moderniser les cultures et l'élevage ; les 18 millions de FCFA de transferts pourraient être orientés vers ces projets au lieu de financer seulement la consommation.

En conclusion, la migration n'est une chance que si le village transforme les transferts en investissements ; sinon, elle vide le village de ses forces vives sans garantir aux jeunes une vie meilleure en ville.
\end{exoresolu}

\courssub{Bilan de l'activité}

Cette activité a permis de mettre en évidence que :
\begin{itemize}
\item les effets des migrations sont \cle{ambivalents} : elles appauvrissent les régions de départ en forces vives tout en les soutenant financièrement par les transferts ; elles dynamisent les villes d'accueil tout en y aggravant chômage, habitat précaire et pression sur les services ;
\item les données chiffrées (taux de variation, pourcentages) permettent de passer d'une impression (\og les jeunes partent \fg) à un diagnostic argumenté (perte de 18,75 \% de la population, 86,7 \% de jeunes parmi les partants) ;
\item un croquis simple à flèches proportionnelles suffit à visualiser le double flux (personnes dans un sens, argent dans l'autre) ;
\item les solutions aux problèmes migratoires existent et relèvent des politiques étudiées : développement rural, création d'emplois et appui à l'entreprenariat des jeunes, coopération entre États et lutte contre l'émigration clandestine ;
\item rédiger une conclusion argumentée exige de trancher (chance ou problème), de justifier par des faits chiffrés et de proposer des solutions réalistes et hiérarchisées : c'est exactement la démarche attendue au Probatoire et au Baccalauréat technique.
\end{itemize}

\newpage

\coursec{Méthodes et exercices résolus}

\courssub{Fiches méthodes et exercices résolus}

\begin{methode}[Appliquer les notions de l'unité à une situation concrète de la vie courante]
Pour analyser une situation migratoire concrète (un village qui se vide, un quartier qui s'étend, un départ à l'étranger) :
\begin{enumerate}
\item \textbf{Identifier le type de migration} : interne ou internationale ? rurale-urbaine (exode rural) ou urbaine-urbaine ? définitive ou temporaire ? volontaire ou forcée ?
\item \textbf{Repérer les facteurs} : facteurs de répulsion (pauvreté, chômage, conflits, manque de services) et facteurs d'attraction (emplois, écoles, hôpitaux, sécurité).
\item \textbf{Analyser les effets en deux lieux distincts} : effets dans la \cle{région de départ} (perte de bras valides, vieillissement, mais envois d'argent) et effets dans la \cle{zone d'accueil} (main-d'œuvre, dynamisme, mais surpopulation, chômage, insalubrité).
\item \textbf{Distinguer effets positifs et effets négatifs} dans un tableau à deux colonnes.
\item \textbf{Conclure} par un bilan nuancé et une piste de solution adaptée à la situation.
\end{enumerate}
Réflexe : toujours répondre en mobilisant le vocabulaire du cours (\cle{exode rural}, \cle{solde migratoire}, \cle{transferts de fonds}, \cle{fuite des cerveaux}).
\end{methode}

\begin{exoresolu}[Application : le départ des jeunes d'un village du Nord vers Douala]
\textbf{Énoncé.} Dans un village de la région de l'Extrême-Nord du Cameroun, 60 jeunes de 18 à 30 ans sur 200 actifs partent chaque année chercher du travail à Douala. Les champs sont de moins en moins cultivés, mais les familles reçoivent régulièrement de l'argent de leurs enfants installés en ville.
\begin{enumerate}
\item Identifie le type de migration décrit et ses facteurs.
\item Présente deux effets négatifs et deux effets positifs de cette migration pour le village.
\end{enumerate}
\tcblower
\textbf{Solution détaillée.}
\begin{enumerate}
\item Il s'agit d'une \cle{migration interne}, plus précisément d'un \cle{exode rural} (déplacement campagne $\rightarrow$ ville), à caractère économique. Les facteurs de \textbf{répulsion} sont la pauvreté rurale, le manque d'emplois et l'absence de perspectives au village ; le facteur d'\textbf{attraction} est l'espoir d'un emploi et d'un revenu à Douala, premier pôle économique du Cameroun.
\item \textbf{Effets négatifs pour le village :}
\begin{itemize}
\item la \textbf{perte de la main-d'œuvre jeune et valide} : 60 départs sur 200 actifs, soit $\dfrac{60}{200}\times 100 = 30\,\%$ des actifs qui partent chaque année, ce qui explique l'abandon des champs et la baisse de la production agricole ;
\item le \textbf{vieillissement de la population} restante, composée surtout de personnes âgées, de femmes et d'enfants, ce qui affaiblit le développement local.
\end{itemize}
\textbf{Effets positifs pour le village :}
\begin{itemize}
\item les \textbf{transferts d'argent} (envois de fonds) améliorent le revenu des familles, permettent de payer la scolarité, les soins et parfois de construire des maisons ;
\item la \textbf{réduction de la pression} sur les terres et sur les rares emplois locaux.
\end{itemize}
\end{enumerate}
\textbf{Conclusion.} Cette migration soulage financièrement les familles à court terme, mais elle fragilise l'économie agricole du village à long terme : seul le développement local (pistes rurales, électrification, crédits agricoles) pourrait retenir les jeunes.
\end{exoresolu}

\begin{methode}[Rédiger et justifier une démarche, une réponse ou une conclusion]
Au Probatoire et au Baccalauréat technique, une réponse de géographie se construit en quatre temps :
\begin{enumerate}
\item \textbf{Lire et souligner} les mots-clés du sujet (le lieu, le phénomène, la consigne : « expliquez », « montrez », « justifiez »).
\item \textbf{Faire un plan brouillon} : deux ou trois parties équilibrées, chacune avec des idées et des exemples précis (chiffres, lieux).
\item \textbf{Rédiger} : introduction (définition des termes + annonce du plan), développement (un paragraphe = une idée + un exemple + une explication), conclusion (bilan + ouverture).
\item \textbf{Justifier} chaque affirmation par un fait vérifiable : donnée chiffrée, exemple localisé, mécanisme de cause à effet (« parce que », « ce qui entraîne »).
\end{enumerate}
Réflexe : une affirmation sans exemple ni chiffre ne vaut que la moitié des points.
\end{methode}

\begin{exoresolu}[Rédaction : les incidences des migrations dans les villes camerounaises]
\textbf{Énoncé.} Sujet type Bac : « Les migrations ont des effets contrastés sur les villes camerounaises. Montrez-le à partir d'exemples précis. » Rédigez le développement de ce sujet (deux paragraphes argumentés).
\tcblower
\textbf{Solution détaillée.}

\emph{Démarche :} le sujet porte sur les \textbf{zones d'accueil} (les villes) et exige de montrer un \textbf{contraste} (effets positifs ET négatifs). Le plan s'impose : I. effets positifs ; II. effets négatifs.

\textbf{Paragraphe 1 — Les effets positifs.} Les migrations fournissent aux villes une \cle{main-d'œuvre} abondante et peu coûteuse qui alimente l'industrie, le bâtiment et le secteur informel : à Douala, port et zone industrielle, les migrants ruraux occupent une large part des emplois de manutention et de petits métiers. Elles accroissent aussi le \textbf{marché de consommation} : l'augmentation de la population de Yaoundé stimule le commerce, les transports (moto-taxis) et le BTP. Enfin, les migrants apportent des savoir-faire et un dynamisme culturel qui enrichissent la vie urbaine.

\textbf{Paragraphe 2 — Les effets négatifs.} La croissance urbaine rapide dépasse les capacités d'accueil : il en résulte la \textbf{surpopulation} et l'extension des \textbf{quartiers précaires} non desservis en eau potable et en assainissement (quartiers spontanés de Douala et de Yaoundé). L'offre d'emplois formels étant insuffisante, le \textbf{chômage} et le \textbf{secteur informel} explosent, avec leurs conséquences : insécurité, délinquance juvénile, mendicité. Enfin, la pression sur les services sociaux (écoles surchargées, hôpitaux saturés) et sur l'environnement (déchets, pollution) dégrade les conditions de vie urbaines.

\textbf{Conclusion.} Les migrations alimentent le dynamisme économique des villes camerounaises, mais, faute de planification urbaine, elles y aggravent chômage, insalubrité et inégalités : d'où la nécessité de politiques d'aménagement du territoire.
\end{exoresolu}

\begin{methode}[Construire un tableau bilan des effets migratoires]
Pour comparer les effets des migrations selon les espaces :
\begin{enumerate}
\item Tracer un tableau à \textbf{deux lignes} (région de départ / zone d'accueil) et \textbf{deux colonnes} (effets positifs / effets négatifs).
\item Remplir chaque case avec \textbf{deux ou trois effets précis}, sans répéter la même idée sous deux formulations.
\item Vérifier la \textbf{symétrie logique} : ce qui est perdu au départ (bras valides) est gagné à l'arrivée (main-d'œuvre).
\item Tirer une conclusion équilibrée : les migrations ne sont ni totalement positives ni totalement négatives, tout dépend des politiques menées.
\end{enumerate}
\end{methode}

\begin{exoresolu}[Tableau bilan : effets des migrations internes au Cameroun]
\textbf{Énoncé.} Complète un tableau présentant deux effets positifs et deux effets négatifs des migrations internes, pour la région de départ et pour la zone d'accueil, puis rédige une phrase de conclusion.
\tcblower
\textbf{Solution détaillée.}
\begin{center}
\begin{tabularx}{\linewidth}{|l|X|X|}\hline
\rowcolor{popblueL} & \textbf{Effets positifs} & \textbf{Effets négatifs} \\ \hline
\textbf{Région de départ (campagne)} & Transferts d'argent ; baisse de la pression sur les terres et l'emploi local & Perte de bras valides ; vieillissement et dépeuplement ; baisse de la production agricole \\ \hline
\textbf{Zone d'accueil (ville)} & Main-d'œuvre abondante ; élargissement du marché ; dynamisme économique & Surpopulation et quartiers précaires ; chômage et informalité ; pression sur écoles et hôpitaux \\ \hline
\end{tabularx}
\end{center}
\textbf{Conclusion.} Les migrations internes transfèrent les hommes et une partie des revenus de la campagne vers la ville : elles profitent aux individus et aux villes, mais appauvrissent les campagnes si l'État n'aménage pas les zones rurales.
\end{exoresolu}

\begin{methode}[Exploiter des données chiffrées sur les migrations]
Pour traiter un document statistique (effectifs de migrants, transferts de fonds) :
\begin{enumerate}
\item \textbf{Lire} le titre, l'unité, la source et les dates du tableau.
\item \textbf{Calculer} les grandeurs utiles :
\begin{itemize}
\item part en pourcentage : $\text{part} = \dfrac{\text{effectif de la catégorie}}{\text{effectif total}} \times 100$ ;
\item taux de variation : $\text{taux} = \dfrac{V_f - V_i}{V_i} \times 100$ ;
\item solde migratoire : $\text{solde} = \text{immigrés} - \text{émigrés}$.
\end{itemize}
\item \textbf{Interpréter} : un solde positif signifie que le territoire gagne de la population ; un solde négatif qu'il en perd.
\item \textbf{Rédiger} une phrase de commentaire par résultat, en citant le chiffre.
\end{enumerate}
Réflexe : arrondir à une décimale et toujours écrire l'unité ou le symbole \%.
\end{methode}

\begin{exoresolu}[Calculs : solde migratoire et évolution des transferts de fonds]
\textbf{Énoncé.} En une année, une région du Cameroun enregistre \num{45000} arrivées de migrants et \num{78000} départs. Par ailleurs, les transferts de fonds des Camerounais de la diaspora vers le pays sont passés de 200 milliards à 250 milliards de FCFA en cinq ans.
\begin{enumerate}
\item Calcule le solde migratoire de la région et interprète-le.
\item Calcule le taux de variation des transferts de fonds et commente.
\end{enumerate}
\tcblower
\textbf{Solution détaillée.}
\begin{enumerate}
\item \textbf{Solde migratoire :}
\[ \text{Solde} = \text{arrivées} - \text{départs} = \num{45000} - \num{78000} = -\num{33000} \text{ personnes.} \]
Le solde est \textbf{négatif} : la région perd \num{33000} habitants de plus qu'elle n'en accueille. C'est une région d'\textbf{émigration nette}, ce qui traduit un exode (manque d'emplois, de services) et risque d'entraîner dépeuplement et vieillissement.
\item \textbf{Taux de variation des transferts :}
\[ \text{Taux} = \dfrac{V_f - V_i}{V_i} \times 100 = \dfrac{250 - 200}{200} \times 100 = \dfrac{50}{200} \times 100 = 25\,\%. \]
Les transferts de fonds de la diaspora ont augmenté de \textbf{25\,\%} en cinq ans, soit 50 milliards de FCFA supplémentaires. Cet argent soutient la consommation des familles, l'éducation, la santé et la construction : c'est un \textbf{effet positif majeur} des migrations internationales pour le Cameroun.
\end{enumerate}
\textbf{Conclusion.} Les chiffres montrent le double visage des migrations : perte de population pour certaines régions, mais manne financière croissante venue de la diaspora.
\end{exoresolu}

\begin{methode}[Analyser les migrations internationales : gains et pertes pour le pays d'origine]
Pour traiter un sujet sur l'émigration internationale des Camerounais :
\begin{enumerate}
\item \textbf{Définir} les notions : \cle{émigration} (départ), \cle{immigration} (arrivée), \cle{diaspora}, \cle{fuite des cerveaux} (émigration des diplômés et cadres qualifiés).
\item \textbf{Lister les gains} : transferts de fonds, acquisition de compétences à l'étranger, réduction du chômage interne, rayonnement du pays.
\item \textbf{Lister les pertes} : fuite des cerveaux (médecins, enseignants, ingénieurs formés à grands frais), déstructuration des familles, trafic d'êtres humains et morts en route (désert, Méditerranée).
\item \textbf{Conclure} en proposant des pistes : création d'emplois pour les jeunes, valorisation des diplômes, coopération avec la diaspora.
\end{enumerate}
\end{methode}

\begin{exoresolu}[Analyse : la fuite des cerveaux camerounais]
\textbf{Énoncé.} De nombreux médecins et infirmiers formés au Cameroun partent travailler en Europe et en Amérique du Nord. Expliquez ce phénomène, présentez ses conséquences pour le Cameroun et proposez deux solutions.
\tcblower
\textbf{Solution détaillée.}

\emph{Explication du phénomène.} Il s'agit de la \cle{fuite des cerveaux} : l'émigration internationale des personnes qualifiées. Elle s'explique par des facteurs de répulsion (salaires faibles, mauvaises conditions de travail, manque d'équipements hospitaliers) et des facteurs d'attraction (salaires élevés, carrières valorisées, meilleures conditions de vie dans les pays développés).

\emph{Conséquences pour le Cameroun.}
\begin{itemize}
\item \textbf{Négatives} : pénurie de personnel de santé (déserts médicaux, surcharge des hôpitaux publics) ; perte sèche pour l'État qui a financé des formations coûteuses sans bénéficier des services ; ralentissement du développement.
\item \textbf{Positives} (à nuancer) : les transferts d'argent des diasporas médicales soutiennent les familles ; certains professionnels reviennent avec de nouvelles compétences.
\end{itemize}

\emph{Solutions possibles.}
\begin{enumerate}
\item Améliorer les \textbf{salaires et les conditions de travail} des personnels de santé (équipements, primes, carrières) pour retenir les compétences.
\item Mettre en place des \textbf{accords de coopération} : retour programmé des cadres, stages de la diaspora dans les hôpitaux camerounais, investissements de la diaspora dans le secteur sanitaire.
\end{enumerate}
\textbf{Conclusion.} La fuite des cerveaux prive le Cameroun de ses forces vives qualifiées ; seule une politique volontariste de rétention et de mobilisation de la diaspora peut la transformer en atout.
\end{exoresolu}

\begin{methode}[Traiter un sujet sur les tentatives de règlement des problèmes migratoires]
Pour un sujet sur les solutions aux problèmes migratoires :
\begin{enumerate}
\item \textbf{Rappeler le problème} en une phrase (exode rural, surpopulation urbaine, émigration clandestine).
\item \textbf{Classer les solutions par échelle} :
\begin{itemize}
\item \emph{locale et nationale} : aménagement des campagnes (électrification rurale, pistes, écoles, centres de santé), création d'emplois, décentralisation, planification urbaine (lotissements, logements sociaux) ;
\item \emph{internationale} : accords bilatéraux sur les migrations de travail, lutte contre l'émigration clandestine et le trafic d'êtres humains, coopération avec les organisations internationales (OIM, HCR), politiques de rapatriement volontaire.
\end{itemize}
\item \textbf{Évaluer} brièvement les limites de ces mesures (moyens financiers insuffisants, application partielle).
\item \textbf{Conclure} : aucune solution unique ; il faut combiner développement des zones de départ et encadrement des flux.
\end{enumerate}
\end{methode}

\begin{exoresolu}[Sujet de synthèse : régler les problèmes de l'exode rural au Cameroun]
\textbf{Énoncé.} « L'exode rural pose de graves problèmes au Cameroun. Après avoir rappelé ces problèmes, présentez les tentatives de règlement possibles. »
\tcblower
\textbf{Solution détaillée.}

\emph{Introduction.} L'\cle{exode rural} est le départ massif des populations rurales, surtout jeunes, vers les villes. Au Cameroun, il alimente la croissance rapide de Douala et Yaoundé. Nous rappellerons ses problèmes, puis nous présenterons les solutions envisageables.

\emph{I. Les problèmes posés.} Dans les campagnes : dépeuplement, vieillissement, champs abandonnés et baisse de la production agricole. Dans les villes : surpopulation, quartiers précaires, chômage, informalité, insalubrité et insécurité.

\emph{II. Les tentatives de règlement.}
\begin{itemize}
\item \textbf{Retenir les populations en milieu rural} : électrification et adduction d'eau, construction de pistes rurales, d'écoles et de centres de santé, appui à l'agriculture (crédits, intrants, encadrement par des organismes comme la SODECOTON dans le Nord ou la SEMRY à Yagoua) ;
\item \textbf{Créer des emplois} en dehors des grandes villes : décentralisation, installation d'administrations et d'industries dans les chefs-lieux de régions, promotion des petites et moyennes entreprises ;
\item \textbf{Aménager les villes} : lotissements, logements sociaux, extension des réseaux d'eau et d'assainissement, structuration du secteur informel ;
\item \textbf{Encadrer les flux} : politiques d'aménagement du territoire et, pour les migrations internationales, accords bilatéraux et coopération avec l'OIM.
\end{itemize}

\emph{Conclusion.} Ces mesures se heurtent au manque de moyens financiers. Régler durablement les problèmes migratoires suppose de rendre la campagne attractive et les villes mieux planifiées : c'est tout l'enjeu de l'aménagement du territoire camerounais.
\end{exoresolu}

\begin{attention}[Les 5 erreurs qui coûtent des points au Bac]
\begin{enumerate}
\item \textbf{Confondre émigration et immigration.} L'\cle{émigration} est le départ (vu du pays d'origine), l'\cle{immigration} l'arrivée (vu du pays d'accueil). Inverser les deux fausse tout le devoir.
\item \textbf{Oublier un des deux espaces.} Un sujet sur les effets des migrations exige TOUJOURS de traiter la région de départ ET la zone d'accueil ; n'en traiter qu'une, c'est perdre la moitié des points.
\item \textbf{Ne voir que le négatif (ou que le positif).} Le correcteur attend un bilan \textbf{contrasté} : effets positifs ET négatifs, même si la consigne semble orientée.
\item \textbf{Affirmer sans chiffrer ni localiser.} « Beaucoup de migrants » ne vaut rien sans exemple précis (Douala, Yaoundé, la diaspora) ou sans calcul (solde migratoire, taux de variation, part en \%).
\item \textbf{Se tromper dans les formules.} Le taux de variation se calcule par rapport à la valeur \emph{initiale} : $\dfrac{V_f - V_i}{V_i}\times 100$, et non $\dfrac{V_f - V_i}{V_f}\times 100$. De même, un solde migratoire négatif signifie plus de départs que d'arrivées : ne pas inverser l'interprétation.
\end{enumerate}
\end{attention}

\newpage

\coursec{Exercices}

\courssub{Je vérifie mes connaissances}

\begin{exercice}[QCM : définitions et acteurs]
Pour chaque question, choisis la bonne réponse et recopie-la sur ta copie.
\begin{enumerate}
\item L'\cle{exode rural} désigne :
\begin{enumerate}
\item le départ des paysans vers les villes ;
\item le départ des citadins vers les villages ;
\item l'expulsion des étrangers d'un pays.
\end{enumerate}
\item La \cle{fuite des cerveaux} est :
\begin{enumerate}
\item l'émigration des cadres et diplômés qualifiés vers l'étranger ;
\item la baisse du niveau scolaire dans les campagnes ;
\item le retour des migrants dans leur pays d'origine.
\end{enumerate}
\item Un \cle{réfugié} est une personne qui :
\begin{enumerate}
\item cherche du travail dans un pays riche ;
\item a fui son pays à cause des persécutions, de la guerre ou des violences ;
\item se déplace à l'intérieur de son propre pays.
\end{enumerate}
\item Un \cle{bidonville} est :
\begin{enumerate}
\item un quartier précaire construit sans autorisation, sans eau ni assainissement ;
\item une ville nouvelle construite par l'État ;
\item un camp de réfugiés géré par le HCR.
\end{enumerate}
\item L'organisation qui s'occupe spécifiquement des réfugiés est :
\begin{enumerate}
\item l'OIM ; \item le HCR ; \item la CEMAC.
\end{enumerate}
\end{enumerate}
\end{exercice}

\begin{exercice}[Vrai ou faux justifié]
Réponds par vrai ou faux, puis justifie ta réponse en une ou deux phrases.
\begin{enumerate}
\item Les migrations ne posent des problèmes que dans les zones d'accueil.
\item L'exode rural rajeunit la population des campagnes camerounaises.
\item L'argent envoyé par les migrants camerounais de l'étranger (transferts de fonds) constitue une ressource pour leurs familles restées au pays.
\item La CEMAC garantit en principe la libre circulation des personnes entre ses États membres.
\end{enumerate}
\end{exercice}

\begin{exercice}[Définitions]
Définis en deux phrases maximum chacun des termes suivants, en donnant pour chacun un exemple camerounais ou africain :
\begin{enumerate}
\item migration interne ;
\item migration internationale ;
\item bidonville ;
\item fuite des cerveaux.
\end{enumerate}
\end{exercice}

\begin{exercice}[Calculs directs sur les transferts de fonds]
En 2020, la diaspora camerounaise a envoyé environ \num{334} millions de dollars à ses familles restées au pays. On admet que 1 dollar vaut 600 FCFA.
\begin{enumerate}
\item Calcule le montant total de ces transferts en milliards de FCFA.
\item Une famille de Yaoundé reçoit \num{120000} FCFA par mois d'un fils installé en France. Quelle somme annuelle cela représente-t-il ?
\item Si cette famille dépense 70 \% de cette aide pour la scolarité et la santé, calcule le montant annuel consacré à ces deux postes.
\end{enumerate}
\end{exercice}

\courssub{Je m'entraîne}

\begin{exercice}[Les effets de l'exode rural sur un village]
Le village de Nkometou (région du Centre) a vu partir en dix ans 60 \% de ses jeunes de 15 à 35 ans vers Yaoundé et Douala.
\begin{enumerate}
\item Cite trois problèmes que ce départ pose au village (main-d'œuvre, production, population).
\item Cite deux avantages que ces départs peuvent malgré tout apporter au village.
\item Propose une mesure concrète qui pourrait retenir les jeunes au village.
\end{enumerate}
\end{exercice}

\begin{exercice}[Les incidences des migrations sur Douala]
La ville de Douala accueille chaque année des dizaines de milliers de migrants venus des régions de l'Ouest, du Nord et de l'Extrême-Nord. Sa population est estimée à plus de 3,5 millions d'habitants.
\begin{enumerate}
\item Cite trois problèmes que cette forte immigration pose à la ville (logement, emploi, hygiène, circulation).
\item Explique en quoi ces migrants contribuent aussi à la vie économique de la ville.
\item Rédige un court paragraphe (5 à 8 lignes) intitulé : « Douala, ville attractive mais saturée ».
\end{enumerate}
\end{exercice}

\begin{exercice}[Gains et pertes des migrations internationales pour le Cameroun]
Recopie et complète le tableau suivant avec deux éléments dans chaque case, en t'aidant de ton cours :
\begin{center}
\begin{tabularx}{\linewidth}{|X|X|}
\hline
\rowcolor{popblueL} \textbf{Gains pour le Cameroun} & \textbf{Pertes pour le Cameroun} \\
\hline
 & \\
\hline
 & \\
\hline
\end{tabularx}
\end{center}
Rédige ensuite une phrase de conclusion qui donne ton avis argumenté : les migrations internationales sont-elles globalement positives ou négatives pour le Cameroun ?
\end{exercice}

\begin{exercice}[Les acteurs de la gestion migratoire]
Associe chaque acteur à son rôle, puis réponds à la question de synthèse.
\begin{center}
\begin{tabularx}{\linewidth}{|l|X|}
\hline
\rowcolor{popblueL} \textbf{Acteur} & \textbf{Rôle} \\
\hline
1. OIM & a. Protection et assistance des réfugiés \\
\hline
2. HCR & b. Aide à la gestion des migrations et au retour volontaire des migrants \\
\hline
3. CEMAC & c. Organisation de la libre circulation des personnes en Afrique centrale \\
\hline
4. État camerounais & d. Contrôle des frontières, délivrance des titres de séjour, politique migratoire nationale \\
\hline
\end{tabularx}
\end{center}
Question de synthèse : pourquoi la gestion des problèmes migratoires nécessite-t-elle la coopération de plusieurs acteurs ? (5 lignes maximum)
\end{exercice}

\courssub{Je me prépare au Bac}

\begin{exercice}[Dissertation de type Baccalauréat]
Sujet : \emph{« Les migrations internes au Cameroun : un problème ou une opportunité pour le développement ? »}

Tu rédigeras une dissertation complète comprenant :
\begin{enumerate}
\item une introduction (accroche, présentation du sujet, problématique, annonce du plan) ;
\item un développement en deux parties :
\begin{itemize}
\item première partie : les migrations internes comme source de problèmes (dans les régions de départ et dans les villes d'accueil) ;
\item deuxième partie : les migrations internes comme opportunité de développement (pour les régions de départ et pour les villes) ;
\end{itemize}
\item une conclusion avec une ouverture.
\end{enumerate}
Barème indicatif : introduction 4 pts ; développement 12 pts ; conclusion 2 pts ; présentation 2 pts.
\end{exercice}

\begin{exercice}[Exercice documentaire : les flux migratoires interrégionaux]
Le tableau ci-dessous présente les principaux flux migratoires interrégionaux au Cameroun (données simplifiées, en milliers de migrants).

\begin{center}
\begin{tabularx}{\linewidth}{|X|X|c|}
\hline
\rowcolor{popblueL} \textbf{Région de départ} & \textbf{Région d'accueil} & \textbf{Migrants (milliers)} \\
\hline
Nord et Extrême-Nord & Centre (Yaoundé) & 180 \\
\hline
Ouest & Littoral (Douala) & 250 \\
\hline
Nord-Ouest & Sud-Ouest & 120 \\
\hline
Est et Adamaoua & Centre (Yaoundé) & 90 \\
\hline
\end{tabularx}
\end{center}

\begin{enumerate}
\item Sur la carte schématique du Cameroun fournie par ton professeur (ou reproduite sur ta copie), construis une carte de flux : représente chaque flux par une flèche d'épaisseur proportionnelle au nombre de migrants (échelle : 1 mm d'épaisseur pour 30 000 migrants). N'oublie pas le titre, la légende et l'orientation.
\item Calcule la part (en \%) du flux Ouest $\rightarrow$ Littoral dans l'ensemble des flux du tableau.
\item Quelles sont les deux régions les plus attractives ? Justifie ta réponse à l'aide des chiffres.
\item Rédige un commentaire de la carte en une dizaine de lignes : tu montreras que les migrations internes au Cameroun s'orientent surtout vers les grandes métropoles du Sud, tu en donneras les causes et tu évoqueras une conséquence pour les régions de départ.
\end{enumerate}
\end{exercice}

\begin{exercice}[Étude de cas : la traversée de la Méditerranée]
Document : \emph{« Chaque année, des milliers de jeunes Camerounais quittent le pays pour tenter de rejoindre l'Europe. Beaucoup traversent le Nigeria, le Niger puis le désert du Sahara jusqu'en Libye, avant d'embarquer sur des embarcations de fortune pour traverser la Méditerranée. Ce voyage, qui peut durer plusieurs mois, coûte entre 1,5 et 3 millions de FCFA, souvent réunis par toute la famille. Sur la route, les migrants risquent l'arrestation, l'exploitation, la violence et la mort : des milliers de personnes se noient chaque année en Méditerranée. Ceux qui arrivent en Europe vivent souvent dans l'irrégularité. Pourtant, certains réussissent et envoient régulièrement de l'argent à leurs familles restées au Cameroun. »}

\begin{enumerate}
\item Dégage du document trois causes qui poussent ces jeunes à quitter le Cameroun.
\item Relève dans le document trois risques encourus pendant le voyage.
\item Explique pourquoi les familles acceptent de financer un voyage aussi dangereux.
\item Cite deux solutions que l'État camerounais et les organisations internationales (OIM, HCR) peuvent mettre en œuvre pour limiter ces migrations dangereuses.
\item Rédige une conclusion personnelle argumentée (8 à 10 lignes) : « Faut-il décourager les départs ou organiser des migrations sûres et légales ? »
\end{enumerate}
\end{exercice}

\newpage

\coursec{Corrigés des exercices}

\coursec{Corrigés des exercices — Leçon 02 : Les effets des mouvements migratoires}

\courssub{Exercice 1 — QCM : définitions et acteurs}
\begin{corrige}[Exercice 1 — QCM : définitions et acteurs]
\begin{enumerate}
\item Réponse \textbf{a} : l'exode rural est le départ des populations rurales (surtout les jeunes) vers les villes, à la recherche d'emplois et de services.
\item Réponse \textbf{a} : la fuite des cerveaux est l'émigration définitive ou durable des cadres, diplômés et personnes qualifiées vers l'étranger (médecins, ingénieurs, enseignants).
\item Réponse \textbf{b} : un réfugié est une personne qui a fui son pays en raison de persécutions, de guerre ou de violences (définition du droit international, HCR).
\item Réponse \textbf{a} : un bidonville est un quartier précaire, construit sans autorisation et sans équipements (eau potable, assainissement, électricité), comme certains quartiers de Douala ou Yaoundé.
\item Réponse \textbf{b} : le HCR (Haut-Commissariat des Nations unies pour les réfugiés) s'occupe spécifiquement de la protection des réfugiés. L'OIM gère les migrations en général ; la CEMAC est une organisation économique régionale.
\end{enumerate}
\textbf{Point de vigilance} : ne pas confondre réfugié (cause politique ou guerre) et migrant économique (recherche d'un emploi) ; cette distinction est souvent demandée au Bac.
\end{corrige}

\courssub{Exercice 2 — Vrai ou faux justifié}
\begin{corrige}[Exercice 2 — Vrai ou faux justifié]
\begin{enumerate}
\item \textbf{Faux.} Les migrations posent aussi des problèmes dans les zones de départ : perte de main-d'œuvre jeune, vieillissement des campagnes, baisse de la production agricole.
\item \textbf{Faux.} L'exode rural fait partir les jeunes actifs : il \emph{vieillit} la population des campagnes, où ne restent souvent que les personnes âgées, les femmes et les enfants.
\item \textbf{Vrai.} Les transferts de fonds de la diaspora (environ 334 millions de dollars en 2020) financent la scolarité, la santé, la construction et le petit commerce des familles restées au pays.
\item \textbf{Vrai.} La CEMAC prévoit en principe la libre circulation des personnes et des biens entre ses États membres, même si son application reste incomplète (contrôles, tracasseries aux frontières).
\end{enumerate}
\end{corrige}

\courssub{Exercice 3 — Définitions}
\begin{corrige}[Exercice 3 — Définitions]
\begin{enumerate}
\item \textbf{Migration interne} : déplacement durable de population à l'intérieur des frontières d'un même pays. Exemple : l'exode rural des jeunes de l'Adamaoua vers Yaoundé.
\item \textbf{Migration internationale} : déplacement durable de population d'un pays vers un autre pays. Exemple : les Camerounais qui s'installent en France, au Gabon ou au Canada.
\item \textbf{Bidonville} : quartier précaire construit sans autorisation, avec des matériaux de récupération, sans eau potable ni assainissement. Exemple : les quartiers marécageux de New-Bell ou Nylon à Douala.
\item \textbf{Fuite des cerveaux} : émigration des personnes qualifiées (médecins, ingénieurs, enseignants) vers des pays offrant de meilleurs salaires. Exemple : les infirmiers camerounais partis travailler en Allemagne ou au Canada.
\end{enumerate}
\textbf{Point de vigilance} : une définition complète comporte l'idée générale + une précision ; l'exemple demandé doit être camerounais ou africain, pas européen.
\end{corrige}

\courssub{Exercice 4 — Calculs directs sur les transferts de fonds}
\begin{corrige}[Exercice 4 — Calculs directs sur les transferts de fonds]
\begin{enumerate}
\item Montant total en FCFA :
\[ 334\,000\,000 \times 600 = 200\,400\,000\,000 \text{ FCFA} \]
soit \textbf{200,4 milliards de FCFA}.
\item Somme annuelle reçue par la famille :
\[ 120\,000 \times 12 = 1\,440\,000 \text{ FCFA par an.} \]
\item Montant consacré à la scolarité et à la santé :
\[ 1\,440\,000 \times \dfrac{70}{100} = 1\,008\,000 \text{ FCFA.} \]
La famille consacre donc \textbf{1 008 000 FCFA} par an à ces deux postes.
\end{enumerate}
\textbf{Vérification} : $334 \times 600 = 200\,400$ (millions), c'est-à-dire 200,4 milliards ; $1\,440\,000 \times 0,7 = 1\,008\,000$. Les calculs sont exacts.
\end{corrige}

\courssub{Exercice 5 — Les effets de l'exode rural sur un village}
\begin{corrige}[Exercice 5 — Les effets de l'exode rural sur un village]
\begin{enumerate}
\item Trois problèmes pour le village de Nkometou :
\begin{itemize}
\item \textbf{Main-d'œuvre} : pénurie de bras pour les travaux champêtres ; les champs sont mal entretenus ou abandonnés.
\item \textbf{Production} : baisse de la production agricole et vivrière, donc baisse des revenus du village.
\item \textbf{Population} : vieillissement du village (il ne reste que les vieillards, les femmes et les enfants) et fermeture possible de services (école, centre de santé).
\end{itemize}
\item Deux avantages possibles :
\begin{itemize}
\item les migrants envoient de l'argent et des biens à leurs familles (transferts de fonds) ;
\item les départs réduisent la pression sur les terres et l'émigration peut financer des projets du village (forage, école, cases).
\end{itemize}
\item Mesure de rétention (une parmi d'autres) : créer des activités rémunératrices au village, par exemple électrifier et désenclaver le village (route, eau potable) et soutenir l'agropastoralisme des jeunes par des crédits et des intrants. Toute mesure concrète et justifiée est acceptée.
\end{enumerate}
\end{corrige}

\courssub{Exercice 6 — Les incidences des migrations sur Douala}
\begin{corrige}[Exercice 6 — Les incidences des migrations sur Douala]
\begin{enumerate}
\item Trois problèmes :
\begin{itemize}
\item \textbf{Logement} : surpopulation, cherté des loyers, extension des bidonvilles et quartiers précaires ;
\item \textbf{Emploi} : chômage et explosion du secteur informel (petits métiers, motos-taxis) ;
\item \textbf{Hygiène et circulation} : insalubrité, ordures, embouteillages permanents, pression sur l'eau et les hôpitaux.
\end{itemize}
\item Contribution économique : les migrants fournissent une main-d'œuvre abondante et bon marché (chantiers, port, marchés), consomment et font marcher le commerce, créent de petites entreprises et paient des taxes : ils dynamisent l'économie de la ville.
\item Paragraphe attendu (exemple de rédaction) : \emph{Douala, capitale économique du Cameroun, attire chaque année des dizaines de milliers de migrants venus de l'Ouest, du Nord et de l'Extrême-Nord, séduits par le port, les industries et les espoirs d'emploi. Cette attractivité fait de Douala une ville de plus de 3,5 millions d'habitants, dynamique et cosmopolite. Mais la ville est saturée : les logements manquent, les bidonvilles s'étendent, les embouteillages et l'insalubrité s'aggravent. Douala illustre ainsi le double visage des migrations : richesse économique et crise urbaine.}
\end{enumerate}
\textbf{Point de vigilance} : le paragraphe doit présenter les deux aspects (attractive \emph{et} saturée) et rester dans les 5 à 8 lignes.
\end{corrige}

\courssub{Exercice 7 — Gains et pertes des migrations internationales pour le Cameroun}
\begin{corrige}[Exercice 7 — Gains et pertes des migrations internationales pour le Cameroun]
Tableau complété (exemples attendus) :
\begin{center}
\begin{tabularx}{\linewidth}{|X|X|}
\hline
\rowcolor{popblueL} \textbf{Gains pour le Cameroun} & \textbf{Pertes pour le Cameroun} \\
\hline
Transferts de fonds de la diaspora (environ 334 millions de dollars en 2020) qui soutiennent les familles & Fuite des cerveaux : départ des médecins, ingénieurs et enseignants formés aux frais de l'État \\
\hline
Acquisition de compétences et investissements des migrants de retour ; réduction du chômage & Morts et drames des migrations clandestines (Sahara, Méditerranée) ; éclatement des familles \\
\hline
\end{tabularx}
\end{center}
Conclusion argumentée (exemple) : \emph{Les migrations internationales sont globalement positives pour le Cameroun à condition d'être organisées, car les transferts de fonds et les compétences de la diaspora dépassent les pertes ; mais la fuite des cerveaux et les drames de l'émigration clandestine montrent qu'elles restent un défi majeur.} Toute position nuancée et justifiée est acceptée.
\end{corrige}

\courssub{Exercice 8 — Les acteurs de la gestion migratoire}
\begin{corrige}[Exercice 8 — Les acteurs de la gestion migratoire]
Associations correctes :
\begin{itemize}
\item 1. OIM $\rightarrow$ \textbf{b} (aide à la gestion des migrations et au retour volontaire) ;
\item 2. HCR $\rightarrow$ \textbf{a} (protection et assistance des réfugiés) ;
\item 3. CEMAC $\rightarrow$ \textbf{c} (libre circulation des personnes en Afrique centrale) ;
\item 4. État camerounais $\rightarrow$ \textbf{d} (contrôle des frontières, titres de séjour, politique nationale).
\end{itemize}
Synthèse (exemple) : \emph{Les migrations dépassent les frontières d'un seul pays : elles concernent les pays de départ, de transit et d'accueil. Aucun État ne peut donc les gérer seul. La coopération entre l'État camerounais, les organisations internationales (OIM, HCR) et les organisations régionales (CEMAC) permet de protéger les migrants, de lutter contre l'émigration clandestine et d'organiser une migration sûre et légale.}
\end{corrige}

\courssub{Exercice 9 — Dissertation de type Baccalauréat}
\begin{corrige}[Exercice 9 — Dissertation de type Baccalauréat]
\textbf{Corrigé-type (plan détaillé et éléments de rédaction).}

\textbf{Introduction (4 pts)} : accroche (le Cameroun connaît de forts mouvements de population, notamment l'exode rural vers Yaoundé et Douala) ; présentation du sujet (les migrations internes, c'est-à-dire les déplacements de population à l'intérieur du pays) ; problématique : \emph{les migrations internes constituent-elles un problème ou une opportunité pour le développement du Cameroun ?} ; annonce du plan en deux parties.

\textbf{Première partie (6 pts) : les migrations internes comme source de problèmes.}
\begin{itemize}
\item Dans les régions de départ : perte de main-d'œuvre jeune, vieillissement des campagnes, baisse de la production agricole, abandon des villages (exemple : départs massifs du Nord et de l'Extrême-Nord).
\item Dans les villes d'accueil : surpopulation (Douala plus de 3,5 millions d'habitants), bidonvilles, chômage et secteur informel, insalubrité, embouteillages, pression sur l'eau, les écoles et les hôpitaux.
\end{itemize}

\textbf{Deuxième partie (6 pts) : les migrations internes comme opportunité.}
\begin{itemize}
\item Pour les régions de départ : transferts d'argent des migrants, financement de projets villageois, réduction de la pression sur les terres.
\item Pour les villes : main-d'œuvre abondante pour l'industrie, le port et le bâtiment ; dynamisme commercial ; brassage culturel et développement des marchés.
\end{itemize}

\textbf{Conclusion (2 pts)} : bilan nuancé (les migrations internes sont à la fois un problème et une opportunité ; tout dépend de leur encadrement) ; ouverture : \emph{vers un meilleur équilibre de l'aménagement du territoire camerounais, ou vers la question des migrations internationales.}

\textbf{Présentation (2 pts)} : orthographe, paragraphes aérés, exemples chiffrés.

\textbf{Points de vigilance} : une problématique sous forme de question ; deux parties équilibrées ; au moins un exemple camerounais chiffré par partie ; éviter le hors-sujet sur les migrations internationales (le sujet porte sur les migrations \emph{internes}).
\end{corrige}

\courssub{Exercice 10 — Exercice documentaire : les flux migratoires interrégionaux}
\begin{corrige}[Exercice 10 — Exercice documentaire : les flux migratoires interrégionaux]
\begin{enumerate}
\item \textbf{Carte de flux} : épaisseurs des flèches à l'échelle 1 mm pour 30 000 migrants :
\begin{itemize}
\item Nord/Extrême-Nord $\rightarrow$ Centre : $180/30 = 6$ mm ;
\item Ouest $\rightarrow$ Littoral : $250/30 \approx 8,3$ mm ;
\item Nord-Ouest $\rightarrow$ Sud-Ouest : $120/30 = 4$ mm ;
\item Est/Adamaoua $\rightarrow$ Centre : $90/30 = 3$ mm.
\end{itemize}
La carte doit comporter : un titre (« Les principaux flux migratoires interrégionaux au Cameroun »), une légende (signification des flèches et échelle d'épaisseur), une flèche d'orientation (nord) et les noms des régions et des villes (Yaoundé, Douala).

\begin{center}
\begin{tikzpicture}[scale=0.85]
% Regions simplified polygons
\draw[popmuted, thick] (-3,3) -- (-1,4) -- (1,3.5) -- (0,1) -- (-2,1) -- cycle; % Nord/EN
\draw[popmuted, thick] (-1,1) -- (2,1) -- (1.5,-0.5) -- (-0.5,-0.5) -- cycle; % Centre/Adamaoua/Est
\draw[popmuted, thick] (-3,-0.5) -- (-1,-0.5) -- (-1.5,-2) -- (-3,-1.5) -- cycle; % Ouest/NO
\draw[popmuted, thick] (-1,-0.5) -- (1.5,-0.5) -- (1,-2) -- (-1,-2) -- cycle; % Littoral/SO

% Cities
\node[circle, fill=popdark, inner sep=1.5pt, label=above:Yaoundé] (Y) at (0,0) {};
\node[circle, fill=popdark, inner sep=1.5pt, label=below right:Douala] (D) at (0.5,-1.5) {};
\node[circle, fill=popdark, inner sep=1.5pt, label=left:Bamenda] (Bam) at (-2.5,0) {};
\node[circle, fill=popdark, inner sep=1.5pt, label=above:Maroua] (Mar) at (-1,3.5) {};

% Arrows (line width in mm converted to pt: 1mm = 2.845pt)
% 1. Nord/EN -> Centre (6mm = 17pt)
\draw[popblue, line width=17pt, -{Stealth[length=4mm]}, opacity=0.8] (-0.5,2.5) -- (0,0.5);
% 2. Ouest -> Littoral (8.3mm = 23.6pt)
\draw[poporange, line width=23.6pt, -{Stealth[length=4mm]}, opacity=0.8] (-2,-1) -- (0.2,-1.3);
% 3. NO -> SO (4mm = 11.4pt)
\draw[popgreen, line width=11.4pt, -{Stealth[length=4mm]}, opacity=0.8] (-2.2,0.2) -- (-1.2,-1.8);
% 4. Est/Adamaoua -> Centre (3mm = 8.5pt)
\draw[poppurple, line width=8.5pt, -{Stealth[length=4mm]}, opacity=0.8] (-1.5,0.5) -- (-0.2,0.2);

% North Arrow
\node at (2.5,3) {\begin{tikzpicture} \draw[popdark, thick, -{Stealth[length=3mm]}] (0,0) -- (0,0.8); \node[below] at (0,0) {N}; \end{tikzpicture}};

% Legend
\begin{scope}[shift={(3.5,-2)}]
\node[align=left, font=\small] at (0,1.5) {\textbf{Légende}};
\draw[popblue, line width=17pt, -{Stealth[length=3mm]}] (0,1) -- (1.5,1); \node[right, font=\small] at (1.5,1) {Nord/EN $\rightarrow$ Centre (180k)};
\draw[poporange, line width=23.6pt, -{Stealth[length=3mm]}] (0,0.5) -- (1.5,0.5); \node[right, font=\small] at (1.5,0.5) {Ouest $\rightarrow$ Littoral (250k)};
\draw[popgreen, line width=11.4pt, -{Stealth[length=3mm]}] (0,0) -- (1.5,0); \node[right, font=\small] at (1.5,0) {NO $\rightarrow$ SO (120k)};
\draw[poppurple, line width=8.5pt, -{Stealth[length=3mm]}] (0,-0.5) -- (1.5,-0.5); \node[right, font=\small] at (1.5,-0.5) {Est/Adamaoua $\rightarrow$ Centre (90k)};
\node[font=\small] at (0,-1.2) {Échelle : 1 mm $\equiv$ 30 000 migrants};
\end{scope}
\end{tikzpicture}
\end{center}
\legende{Les principaux flux migratoires interrégionaux au Cameroun (effectifs en milliers, épaisseur proportionnelle).}

\item Part du flux Ouest $\rightarrow$ Littoral :
\[ \text{Total des flux} = 180 + 250 + 120 + 90 = 640 \text{ milliers de migrants.} \]
\[ \text{Part} = \dfrac{250}{640} \times 100 = 39{,}0625 \approx \textbf{39,1 \%}. \]

\item Les deux régions les plus attractives sont le \textbf{Centre} (Yaoundé) avec $180 + 90 = 270$ milliers de migrants accueillis, et le \textbf{Littoral} (Douala) avec 250 milliers. Ce sont les deux métropoles du Sud, sièges de l'administration (Yaoundé) et de l'économie portuaire et industrielle (Douala).

\item Commentaire (exemple) : \emph{La carte montre que les migrations internes au Cameroun s'orientent essentiellement vers les grandes métropoles du Sud : le Centre (Yaoundé) accueille 270 000 migrants et le Littoral (Douala) 250 000, soit plus de 80 \% des flux recensés. Le flux le plus important relie l'Ouest au Littoral (250 000 migrants, environ 39 \% du total). Ces mouvements s'expliquent par la recherche d'emplois, par la présence du port et des industries à Douala, de l'administration et des universités à Yaoundé, et par la pauvreté ou l'insécurité dans certaines régions de départ (Nord, Extrême-Nord). Conséquence : les régions de départ perdent leur jeunesse et leur main-d'œuvre, ce qui aggrave leur retard de développement.}
\end{enumerate}
\textbf{Points de vigilance} : respecter l'échelle d'épaisseur (calculs vérifiés : $250/30 \approx 8,3$ mm) ; ne pas oublier titre, légende et orientation, qui valent des points ; le commentaire doit suivre la consigne (orientation, causes, une conséquence).
\end{corrige}

\courssub{Exercice 11 — Étude de cas : la traversée de la Méditerranée}
\begin{corrige}[Exercice 11 — Étude de cas : la traversée de la Méditerranée]
\begin{enumerate}
\item Trois causes des départs (dégagées ou déduites du document) : la recherche d'une vie meilleure en Europe (espoir de travail et de revenus) ; le manque de perspectives d'emploi pour les jeunes au Cameroun ; la pression et l'investissement de la famille qui réunit l'argent du voyage et attend un retour (les transferts de ceux qui réussissent).
\item Trois risques relevés dans le document : l'arrestation et l'exploitation sur la route ; les violences et la mort (noyades en Méditerranée, dangers du désert du Sahara) ; la vie dans l'irrégularité pour ceux qui arrivent en Europe.
\item Les familles financent le voyage (1,5 à 3 millions de FCFA) parce qu'elles le considèrent comme un \emph{investissement} : si le migrant réussit, il enverra régulièrement de l'argent qui fera vivre toute la famille, paiera la scolarité, la santé et la construction. L'espoir de réussite l'emporte sur la conscience du danger.
\item Deux solutions possibles :
\begin{itemize}
\item État camerounais : créer des emplois et des formations pour les jeunes, informer sur les dangers des migrations clandestines, ouvrir des voies de migration légale (accords de travail) ;
\item OIM et HCR : organiser les retours volontaires, sensibiliser aux risques, protéger les migrants en transit et les réfugiés, coopérer avec les pays de transit.
\end{itemize}
\item Conclusion personnelle (exemple) : \emph{Décourager les départs par la seule répression ne suffit pas : tant que le chômage des jeunes persistera, les candidats au départ prendront tous les risques, comme le montrent les drames de la Méditerranée. La meilleure solution est double : d'une part créer au Cameroun des emplois et des perspectives pour retenir les jeunes, d'autre part organiser avec les pays d'accueil et les organisations internationales des migrations sûres et légales (visas de travail, formation, mobilité encadrée). Une migration organisée profite au migrant, à sa famille grâce aux transferts de fonds, et au pays tout entier, tandis que la migration clandestine ne produit que des victimes.} Toute prise de position argumentée et cohérente est acceptée.
\end{enumerate}
\textbf{Point de vigilance} : pour les questions 1 et 2, les réponses doivent s'appuyer sur le document (le relever ou le reformuler) ; pour la question 5, l'avis personnel doit être justifié par au moins deux arguments.
\end{corrige}

\newpage

\coursec{Fiche bilan}

\coursec{Leçon 02 : Les effets des mouvements migratoires}

\begin{popfigure}
\begin{tikzpicture}[
  mindmap,
  concept color=popblue,
  level 1/.style={level distance=4.5cm, sibling angle=60},
  level 2/.style={level distance=3.2cm, sibling angle=45},
  every node/.style={font=\small, text width=2.8cm, align=center},
  grow cyclic
]
\node[concept, text width=3.5cm, font=\bfseries\large] {Effets des mouvements migratoires}
  child[concept color=popblue] {
    node[concept, text width=2.5cm] {1. Typologie des effets}
    child { node[concept, text width=2.2cm] {Positifs / Négatifs} }
    child { node[concept, text width=2.2cm] {Régions départ / arrivée} }
    child { node[concept, text width=2.2cm] {Court / Long terme} }
  }
  child[concept color=poporange] {
    node[concept, text width=2.5cm] {2. Problèmes régions de départ}
    child { node[concept, text width=2.2cm] {Perte main-d'œuvre active} }
    child { node[concept, text width=2.2cm] {Vieillissement population} }
    child { node[concept, text width=2.2cm] {Déséquilibre sexes} }
    child { node[concept, text width=2.2cm] {Transferts financiers (gain)} }
  }
  child[concept color=popgreen] {
    node[concept, text width=2.5cm] {3. Incidences zones d'accueil}
    child { node[concept, text width=2.2cm] {Pression services/logement} }
    child { node[concept, text width=2.2cm] {Main-d'œuvre disponible} }
    child { node[concept, text width=2.2cm] {Tensions sociales/xénophobie} }
    child { node[concept, text width=2.2cm] {Enrichissement culturel} }
  }
  child[concept color=poppurple] {
    node[concept, text width=2.5cm] {4. Migrations int. : Cameroun}
    child { node[concept, text width=2.2cm] {Gains : transferts, compétences retour} }
    child { node[concept, text width=2.2cm] {Pertes : fuite cerveaux, main-d'œuvre qualifiée} }
    child { node[concept, text width=2.2cm] {Solde migratoire souvent négatif} }
  }
  child[concept color=poppink] {
    node[concept, text width=2.5cm] {5. Tentatives de règlement}
    child { node[concept, text width=2.2cm] {Politiques nationales migration} }
    child { node[concept, text width=2.2cm] {Accords bilatéraux / multilatéraux} }
    child { node[concept, text width=2.2cm] {Développement local (création emplois)} }
    child { node[concept, text width=2.2cm] {Engagement diaspora / co-développement} }
  }
  child[concept color=popgold] {
    node[concept, text width=2.5cm] {6. Méthode d'analyse (TP)}
    child { node[concept, text width=2.2cm] {Cartographier flux (flèches proportionnelles)} }
    child { node[concept, text width=2.2cm] {Quantifier (volume, taux)} }
    child { node[concept, text width=2.2cm] {Qualifier (profil, causes, effets)} }
    child { node[concept, text width=2.2cm] {Conclure (bilan coûts/avantages)} }
  };
\end{tikzpicture}
\legende{Carte mentale récapitulative : effets des migrations (départ, arrivée, cas camerounais, régulation, méthode).}
\end{popfigure}

\courssub{Définitions clés}
\begin{itemize}
  \item \cle{Migration} : déplacement durable d'individus impliquant un changement de résidence habituelle.
  \item \cle{Émigration} / \cle{Immigration} : sortie / entrée sur un territoire donné.
  \item \cle{Solde migratoire} = Immigrés $-$ Émigrés (positif = gain, négatif = perte).
  \item \cle{Fuite des cerveaux (brain drain)} : départ de personnel hautement qualifié vers l'étranger.
  \item \cle{Transferts de fonds (remittances)} : envois d'argent des migrants vers le pays d'origine.
  \item \cle{Co-développement} : partenariat migrants / pays d'origine / pays d'accueil pour le développement local.
  \item \cle{Facteurs d'expulsion (push)} / \cle{Facteurs d'attraction (pull)} : causes du départ / attraits de la destination.
\end{itemize}

\courssub{Tableau de synthèse : Effets départ vs Arrivée}
\begin{center}
\begin{tabularx}{\linewidth}{|l|X|X|}\hline
\rowcolor{popblueL}
\textbf{Dimension} & \textbf{Régions de départ} & \textbf{Zones d'accueil} \\ \hline
\textbf{Démographique} & Perte population active, vieillissement, déséquilibre sexes & Croissance pop., rajeunissement, recomposition familles \\ \hline
\textbf{Économique} & $-$ Main-d'œuvre qualifiée ; $+$ Transferts fonds, investissements retour & $+$ Main-d'œuvre (souvent peu qualifiée), consommation ; $-$ Coûts intégration \\ \hline
\textbf{Social} & Déstructuration familles, perte repères culturels & Tensions logement/emploi, xénophobie, enrichissement culturel \\ \hline
\textbf{Politique} & Pression pour réformes, rôle diaspora & Gestion flux, intégration, droits migrants \\ \hline
\end{tabularx}
\end{center}

\courssub{Le cas camerounais : Bilan gains / pertes}
\begin{itemize}
  \item \textbf{Gains} : Transferts \num{300} à \num{400} milliards FCFA/an (source BEAC / Banque Mondiale), acquisition compétences (retour, formation), réseaux diaspora (affaires, santé, éducation).
  \item \textbf{Pertes} : Fuite cerveaux (médecins, ingénieurs, universitaires), coût formation non récupéré, pénurie compétences secteurs stratégiques (santé, éducation, industrie).
  \item \textbf{Solde} : Structurellement négatif pour les qualifiés ; enjeu : transformer \emph{brain drain} en \emph{brain circulation} / \emph{brain gain}.
\end{itemize}

\begin{methode}[Analyser un flux migratoire (TP guidé)]
\begin{enumerate}
  \item \textbf{Localiser \& Cartographier} : Carte de base, flèches proportionnelles aux volumes, sens (origine $\to$ destination), légende complète (titre, échelle, nord, source).
  \item \textbf{Quantifier} : Volume annuel, taux d'émigration/immigration, part dans population totale, évolution temporelle.
  \item \textbf{Qualifier} : Profil migrants (âge, sexe, diplôme, secteur), causes (push/pull), statut (régulier/irrégulier).
  \item \textbf{Évaluer les effets} : Tableau coûts/avantages pour départ et arrivée (cf. tableau de synthèse).
  \item \textbf{Conclure \& Proposer} : Bilan global, pistes de régulation adaptées au contexte.
\end{enumerate}
\end{methode}

\courssub{Principales pistes de règlement (niveau national \& international)}
\begin{itemize}
  \item Politiques nationales : quota, visa, intégration, lutte trafic, incitation retour (ex: loi sur la diaspora camerounaise).
  \item Accords bilatéraux : gestion flux circulaires, reconnaissance diplômes, protection sociale (ex: Cameroun--France, CEMAC).
  \item Cadre multilatéral : Pacte mondial pour des migrations sûres, ordonnées et régulières (ONU 2018), OIM, UA (Protocole sur la libre circulation des personnes, 2018).
  \item Développement territorial : créer emplois \emph{in situ} (zones économiques, agriculture, numérique) pour réduire les \cle{facteurs d'expulsion}.
  \item Mobilisation diaspora : transferts productifs, investissement, transfert savoir, mentorat.
\end{itemize}

\begin{aretenir}[Checklist avant le Bac]
  $\square$ Je définis migration, émigration, immigration, solde migratoire, fuite des cerveaux, transferts de fonds.\par
  $\square$ Je distingue effets positifs et négatifs dans les régions de départ (démo, éco, social).\par
  $\square$ Je distingue effets positifs et négatifs dans les zones d'accueil (marché travail, services, cohésion).\par
  $\square$ Je cite des chiffres clés pour le Cameroun (montant transferts, secteurs touchés par fuite cerveaux).\par
  $\square$ J'explique le paradoxe : pertes de compétences vs gains financiers pour le pays d'origine.\par
  $\square$ Je connais 3 instruments de régulation (national, bilatéral, multilatéral / ONU-UA).\par
  $\square$ Je sais construire un croquis de flux : flèches proportionnelles, légende, titre, source.\par
  $\square$ J'applique la méthode en 5 étapes pour analyser un flux migratoire concret.\par
  $\square$ J'argumente une réponse structurée : thèse / antithèse / synthèse avec exemples camerounais/africains.\par
  $\square$ Je relie migration et développement (co-développement, ODD 10.7).\par
  $\square$ Je maîtrise le vocabulaire technique : push/pull factors, brain drain/gain/circulation, migrant/réfugié/déplacé.
\end{aretenir}

\findecours

\end{document}
